\documentclass[11pt,oneside,openany,notitlepage]{report}
\usepackage[letterpaper,margin=1in]{geometry}
\usepackage[T1]{fontenc}
\usepackage{lmodern}
\usepackage{amsmath,amssymb,amsthm}
\usepackage{booktabs,longtable,array}
\usepackage{titlesec}
\titleformat{\chapter}[display]{\normalfont\bfseries}{\large\chaptertitlename\ \thechapter}{.4em}{\huge}
\titlespacing*{\chapter}{0pt}{0pt}{1.3em}
\setcounter{tocdepth}{0}
\raggedbottom
\usepackage{tikz}
\usepackage[hidelinks]{hyperref}
\hypersetup{pdftitle={Reductions for exact equilibrium in Bayesian simple security games},
  pdfauthor={Research project coordinated by Vincent Lin}}
\urlstyle{same}
\newtheorem{theorem}{Theorem}[chapter]
\newtheorem{lemma}[theorem]{Lemma}
\newtheorem{corollary}[theorem]{Corollary}
\newcommand{\Market}{\mathsf{M1}}
\newcommand{\Traffic}{\mathsf{Traffic}}
\newcommand{\BNE}{\mathsf{BNE}}
\newcommand{\Arrival}{\mathsf{ARRIVAL}}
\newcommand{\PLCP}{\mathsf{P\mbox{-}LCP}}
\newcommand{\Unique}{\mathsf{Unique}}
\newcommand{\SSG}{\mathsf{SSG}}
\newcommand{\dd}{\mathrel{\leq_{\mathrm m}}}
\newcommand{\ttred}{\mathrel{\leq_{\mathrm T}}}
\newcommand{\reviewdate}{November 26, 2026 (UTC)}
\newcommand{\reviewmark}[1]{\par\nopagebreak[4]\noindent
  {\footnotesize\itshape AI self-review, \reviewdate. #1\par}\nobreak\smallskip}
\DeclareMathOperator{\caparc}{cap}
\DeclareMathOperator{\bit}{bit}
\DeclareMathOperator{\supp}{supp}
\title{Reductions for exact equilibrium\\in Bayesian simple security games}
\author{Research project coordinated by Vincent Lin}
\date{Working draft \(\cdot\) \reviewdate}

\begin{document}
\maketitle

\begin{abstract}
Computing an exact Bayes--Nash equilibrium in a Bayesian simple security
game, clearing a divisible matching market with linear transfers, and
computing a multiclass Wardrop equilibrium on parallel roads are equivalent
under polynomial-time many-one reductions. The number of traveler classes
and roads is unrestricted, and each class may have its own strictly
increasing affine cost on each road.

Two constructions yield this equivalence. The first translates directly
between traffic and matching. The second encodes a polynomial sequence of
adaptive market queries in one market, strengthening the game-to-market
direction of Lin's earlier manuscript. That manuscript also gives the
one-resource market-to-game embedding, the reduction from ARRIVAL, and the
upper reduction to P-matrix linear complementarity.

We also recover the entire vector of binary Multi-Arrival terminal counts
from one security equilibrium, using one switch construction per original
nonterminal. Given one equilibrium of a game, we can decide full-profile
uniqueness and produce a second equilibrium whenever one exists. Further
chapters treat restricted market inputs, attack supports, and stopping
stochastic games. All reductions specify an output recoverable from every
admissible oracle answer.
\end{abstract}
\clearpage
{\small\tableofcontents}

\chapter{Which problems reduce to equilibrium search?}\label{sec:intro}
\reviewmark{Reduction directions, oracle calls and prior-work attributions.}

Travelers pick roads of minimum cost at the prevailing loads. Buyers pick
goods of maximum utility at the prevailing prices. Attackers pick targets
of maximum payoff at the prevailing coverage. Translating the choices
also requires translating the other equilibrium condition. Road loads
sum the class flows, prices clear fixed capacities, and coverage maximizes
the defender's payoff against the attacks.

Write \(\Traffic\) for exact parallel-road equilibrium, \(\Market\) for the
divisible affine matching market in Chapter~\ref{sec:models}, and
\(\BNE\) for exact equilibrium in Conitzer's simultaneous Bayesian simple
security game~\cite{conitzer}. Numerical inputs are rationals encoded in binary.
For search problems, \(A\dd B\) uses one polynomial-time transformation
and a polynomial-time decoder that works for \emph{every} solution of the
transformed instance. We use standard rational encodings and a fixed
polynomial witness-length bound for each search problem, large enough
to include a basic feasible solution of every feasible support LP and,
for a game, its marginal realization. Oracle algorithms have a
polynomial clock in the original input length on every admissible
transcript, including the time to read answers. The notation
\(A\ttred B\) allows polynomially many adaptive calls. A nonadaptive
batch fixes every query in advance.
Decision reductions recover a bit. Function reductions recover the entire
specified output, including every component of a terminal-count vector.

Dated notes identify what the drafting AI checked in this self-review.
The UTC date corresponds to November 25 in the project's Pacific time zone.

\begin{theorem}[Traffic and matching]\label{thm:traffic}
\(\Traffic\dd\Market\) and \(\Market\dd\Traffic\).
The traffic oracle takes instances with two vertices, complete menus,
positive slopes and nonnegative intercepts. Both slope and intercept may depend on both the
class and the road.
\end{theorem}

Chapters~\ref{sec:forward} and~\ref{sec:device} prove both directions
directly. A reader interested in the single-query improvement for security
games can instead follow Chapter~\ref{sec:bridge} through
Theorem~\ref{thm:compile} in Appendix~\ref{sec:compile}. The earlier
adaptive reduction already gives Turing equivalence of all three problems
when combined with Theorem~\ref{thm:traffic}.

The principal chain is
\[
 \Arrival\ \dd\ \Market
 \ \mathrel{\equiv_{\rm m}}\ \Traffic
 \ \mathrel{\equiv_{\rm m}}\ \BNE
 \ \dd\ \PLCP .
\]
An arrow points from the problem being solved to the problem queried.
Lin's earlier manuscript gives the ARRIVAL arrow, the one-resource
embedding \(\Market\dd\BNE\), the adaptive reduction \(\BNE\ttred\Market\),
and the direct P-LCP upper bound~\cite[Chapters~3--5]{book}.

\paragraph{What this answers about Conitzer's question.}
An algorithm for finding one exact solution to any of the three equivalent
search problems gives an algorithm for the other two by one call and
polynomial additional bit work. The same holds if the available algorithm
handles only the one-resource games of Theorem~\ref{thm:forms}.
ARRIVAL supplies a lower bound on what a BNE oracle can compute. The
P-matrix reduction places the search problem in \(\mathsf{UEOPL}\).
A polynomial-time algorithm for general equilibrium search, or a hardness
reduction from a complete problem for \(\mathsf{UEOPL}\), would answer a
further question. Neither follows by composing the arrows above.
The completeness result in Chapter~\ref{sec:decisions} concerns decision
algorithms \emph{already given} an equilibrium oracle.

\begin{longtable}{@{}>{\raggedright\arraybackslash}p{.31\linewidth}@{\hspace{.03\linewidth}}>{\raggedright\arraybackslash}p{.50\linewidth}@{\hspace{.03\linewidth}}r@{}}
\toprule
Problem & What a query supplies & Chapter\\ \midrule
\endfirsthead
\toprule Problem & What a query supplies & Chapter\\ \midrule
\endhead
\bottomrule\endfoot
Affine matching & One security equilibrium gives prices and trades. One market gives a full security equilibrium. & \ref{sec:bridge}, \ref{sec:compile}\\
Parallel-road traffic & One market gives class flows. One traffic instance gives market prices and trades. & \ref{sec:forward}--\ref{sec:device}\\
Hide-and-seek & One matching solution or one mixed Nash equilibrium suffices in the respective direction. & \ref{sec:matching-links}\\
Fixed-total rent division & One security equilibrium gives signed rents and a perfect assignment on favorite rooms. & \ref{sec:rent}\\
Market updates and extreme prices & An equilibrium after a unit entry, or at a specified price extreme, encodes ordinary matching. Each converse uses one market. & \ref{sec:update}\\
Prices on specified vertices & One residual market and a flow test decide whether pinned buyer utilities and good prices extend. & \ref{sec:pins}\\
Splitting graph edges & With one price pinned, one clearing split gives a market answer and conversely. & \ref{sec:pinned}\\
Quotas and cycles & A half-quota clearing gives one market answer. One cycle query per good recovers prices. The reverse uses an adaptive market oracle. & \ref{sec:quota}\\
Stable matching and Boolean search & A market gives stable partners, or an ordered fixed-point or crossing witness for the stated Boolean families. & \ref{sec:boolean}\\
ARRIVAL & One one-resource game yields the destination bit and a switching-flow certificate on the augmented graph. & \ref{sec:arrival}\\
Binary Multi-Arrival & One ordinary ARRIVAL instance gives a count bit. Terminal purchases in one market or game give the entire vector directly. & \ref{sec:multi}\\
Attack forests and sparse games & A positive forest can encode a full equilibrium. A game with two attainable targets per type and three types per target suffices for search. & \ref{sec:forests}--\ref{sec:sparse}\\
Full-profile uniqueness & One equilibrium of the input game permits a uniqueness test and a second profile when one exists. A batch of uniqueness bits recovers an equilibrium. & \ref{sec:unique}--\ref{sec:decisions}\\
P-LCP and contractions & One P-LCP solution or one exact contraction fixed point gives a BNE. The reverse market construction covers a family with a direct convex quadratic algorithm. & \ref{sec:upper}, \ref{sec:cayley}\\
Stopping stochastic games & One explicit start value suffices for two targets. General inputs use polynomial-bit singleton values of a succinct game, or adaptive threshold decisions. & \ref{sec:ssg}\\
\end{longtable}

\paragraph{Relation to earlier work.}
Galichon and Jacquet~\cite{gj} study matching with linearly transferable
utility and singles. The label \(\Market\) is the earlier security-game
manuscript's notation for their market. Chapter~\ref{sec:matching-links}
identifies the corresponding family of hide-and-seek games and states what
the equilibrium encodes. Meunier and Pradeau~\cite[Theorem~1]{mp} give a polynomial algorithm
for affine multiclass traffic when both the number of vertices and the
number of classes are fixed. Here the network always has two vertices and
both the number of classes and the number of parallel roads are input
parameters. The counting network in Chapter~\ref{sec:device} is due to
Disser and Skutella~\cite[Section~3.1]{ds}. We use its marginal cost to
represent a traffic load by a clearing price. For switching graphs we use the
standard ARRIVAL decision problem~\cite{arrival} and order independence of
Multi-Arrival~\cite[Proposition~4.3]{glm}. The project's AI agents proposed
reducing ARRIVAL to exact BNE with one defender resource and developed
the construction in Lin's earlier manuscript~\cite[Chapter~4]{book}.
In a later comment on that manuscript, quoted here from Lin's record,
Conitzer wrote ``Incidentally someone else suggested ARRIVAL here''
and linked to Jeonseong Kim's post~\cite{kim}.
Chapter~\ref{sec:arrival} gives the construction and decoder, with a direct
proof that fractional solutions give the correct terminal counts. A note on
interchangeability~\cite[Corollary~5]{generic} establishes uniqueness of the
defender strategy for generic priors with one resource, fixed positive
defender gains, and positive attacker coverage effects.
Our test decides full-profile uniqueness on a specified input, including
ties, using a connectedness theorem of Jones and Gowda~\cite{jg}.

\chapter{Models and exact recovery from a support}\label{sec:models}


\section{The simultaneous security game}
\reviewmark{Type conventions, cutoff optimality and lottery representation.}
There are \(n\ge1\) targets, \(m\ge1\) attacker types of prior
\(p_i\ge0\), \(\sum_i p_i=1\), and an integral defender budget \(0\le r\le n\).
We write \([n]=\{1,\ldots,n\}\).
A pure defense is an \(r\)-subset. Its mixed strategy has target marginals
in the hypersimplex
\[
 H_r=\{x\in[0,1]^n:\ \textstyle\sum_jx_j=r\}.
\]
Write \(\sigma_{ij}\) for the conditional probability with which type \(i\)
attacks \(j\). At a marginal equilibrium, each \(\sigma_i\) is supported on
maximizers of \(U_{ij}(x)=a_{ij}-b_{ij}x_j\), \(b_{ij}\ge0\), and the
defender maximizes \(\pi(\sigma)\cdot x\) on \(H_r\), where
\(\pi_j(\sigma)=d_j\sum_i p_i\sigma_{ij}\), \(d_j\ge0\). We require a
conditional best response even at a listed zero-prior type. A \emph{full
profile} specifies the defender's distribution on pure defenses and the
attacker's conditional distribution for each type.

Uniqueness uses this conditional-strategy convention, including choices
at listed zero-prior types. Different lotteries on pure defenses count
as different full profiles even when they have the same marginals.

For the positive-prior rows let \(f_{ij}=p_i\sigma_{ij}\) denote joint
attack mass, so the score of target \(j\) is \(d_j\sum_i f_{ij}\).
A coverage \(x\) is optimal against these scores precisely when a cutoff
\(\lambda\) makes \(x_j>0\) imply \(\pi_j\ge\lambda\) and \(x_j<1\)
imply \(\pi_j\le\lambda\). Moving coverage from a lower to a higher
score proves necessity, and allocating to the largest scores proves
sufficiency. Every point of \(H_r\) is realized by a lottery on at most
\(n\) pure defenses with rational, polynomial-size support when the
marginals are rational. For example, partition \([0,r)\) consecutively
into intervals of lengths \(x_j\). A uniform \(s\in[0,1)\) selects the
half-open intervals hit by \(s,s+1,\ldots,s+r-1\). Each interval has
length at most one, so each selected set has size \(r\) and each marginal is \(x_j\).
Only fractional parts of the interval endpoints change the selected set.
The consecutive subinterval lengths in \([0,1)\) give the lottery weights.
For \(r=0\) take the empty defense.
\section{Parallel roads}
\reviewmark{Class-dependent costs, existence, and the distinction between total loads and class flows.}
There are two vertices joined by parallel directed roads. Class \(i\) has
mass \(h_i>0\) of nonatomic travelers. If \(f_{ij}\ge0\) travel on road \(j\)
and \(X_j=\sum_i f_{ij}\), its cost to class \(i\) is
\(\ell_{ij}(X_j)=A_{ij}+B_{ij}X_j\), where \(A_{ij}\ge0\) and \(B_{ij}>0\).
A Wardrop equilibrium is a matrix of class flows with row sums \(h_i\)
such that \(f_{ij}>0\) implies
\(\ell_{ij}(X_j)\le\ell_{ik}(X_k)\) for every \(k\). The total mass is
\(T=\sum_i h_i\). Intermediate instances can have restricted nonempty menus,
constant costs, and signed intercepts. Existence follows from the continuous
variational inequality on the compact product of the class simplices
\cite{mp}.

For the strictly increasing costs of Theorem~\ref{thm:traffic}, the
aggregate road loads \(X\) are unique. Suppose equilibria \(f,g\) give
different loads \(X,Y\) and put \(S=\{j:X_j>Y_j\}\). Since road loads
have fixed total, \(S\) gains flow. Some class \(i\) sends more to \(S\)
under \(f\) and hence less to its complement. Choose \(j\in S\) used
by \(i\) under \(f\), and \(k\notin S\) used by \(i\) under \(g\).
Optimality at both profiles gives
\[
 \ell_{ij}(Y_j)<\ell_{ij}(X_j)\le\ell_{ik}(X_k)
       \le\ell_{ik}(Y_k)\le\ell_{ij}(Y_j),
\]
a contradiction. This parallel-road case fits the network uniqueness
results of Milchtaich~\cite{milchtaich}. The \emph{class flows} can vary.
For example, two unit classes with cost \(X_j\) on each of two roads
have loads \((1,1)\) under either assignment of one class per road.
Given \(X\), a transportation LP recovers class flows with the prescribed
row and column sums and permits only minimum-cost class--road pairs.

\section{Affine matching}
\reviewmark{Outside options, the two repair LPs and their denominator bounds.}
Buyer \(i\) has mass \(m_i>0\), good \(j\) has capacity \(c_j>0\), and
a listed pair \(ij\) has utility \(a_{ij}-b_{ij}P_j\) with \(b_{ij}>0\).
The outside utility and reserve prices are zero. An equilibrium consists of
prices \(P_j\ge0\), values \(u_i\ge0\), and trades \(g_{ij}\ge0\), satisfying
\begin{align}
 u_i&\ge a_{ij}-b_{ij}P_j
       &&\text{with equality if }g_{ij}>0,\nonumber\\
 \sum_jg_{ij}&\le m_i &&\text{with equality if }u_i>0,\label{eq:m1}\\
 \sum_i g_{ij}&\le c_j &&\text{with equality if }P_j>0.\nonumber
\end{align}
Equivalently, the no-blocking inequality is
\(\lambda_{ij}u_i+(1-\lambda_{ij})P_j\ge\Phi_{ij}/2\) with
\(\lambda_{ij}=1/(1+b_{ij})\) and \(\Phi_{ij}=2a_{ij}/(1+b_{ij})\),
as in~\cite{gj,book}. An unlisted option may be given a negative intercept.
Isolated buyers and goods have value or price zero and can be reinstated
after solving. We may assume at least one listed pair.

\begin{lemma}[Recovery by linear programming]\label{lem:repair}
For either model, one can compute a positive integer \(Q\) of polynomial bit
length with the following guarantees.
\begin{enumerate}
\item For traffic, suppose a feasible flow uses roads of excess cost strictly
less than \(1/Q\) for their classes. Its support determines an exact rational
equilibrium in polynomial time.
\item For matching, suppose prices, values and trades satisfy the payoff and
row conditions in~\eqref{eq:m1}. Allow an additive capacity excess on any
column and an additive shortfall on positive-price columns, each less than
\(1/Q\). Assume \(P_j\le P_*:=1+\max(0,\max_{ij}a_{ij}/b_{ij})\).
The trade support and the zero/nonzero labels of prices and values then
determine an exact rational equilibrium in polynomial time.
\end{enumerate}
\end{lemma}
\begin{proof}
For traffic fix the supplied support \(S\). Keep nonnegativity, the row sums,
the load identities, and zeros off \(S\). For each \(ij\in S\) impose
\(\ell_{ij}(X_j)\le\ell_{ik}(X_k)+t\) for every allowed \(k\). Minimize
\(0\le t\le1\). For matching set unrecorded trades and zero-labeled
coordinates to zero, keep all payoff inequalities and equality on the
recorded support, and keep the row bounds with equality on positive-value
rows. Minimize \(t\) subject to \(0\le t\le1\). Impose
\(\sum_i g_{ij}\le c_j+t\) for all goods and
\(\sum_i g_{ij}\ge c_j-t\) on positive-price goods. Bound prices by \(P_*\)
and values by \(1+\max(0,\max_{ij}a_{ij})\). The latter bound holds because
a buyer with positive value uses some listed pair. No label imposes a strict
inequality in either LP.

In each model there are only polynomially many \emph{individual} rows
that can occur over all labels. Clear denominators row by row, including
the right-hand side. If \(D\) is the number of variables including \(t\)
and \(H\ge1\) bounds the resulting absolute integer entries, use
\(Q=(DH)^D\). Each feasible LP is bounded. Cramer's rule bounds the
denominator of \(t\) at an optimal vertex by \(Q\). The supplied point
permits \(t<1/Q\), so the optimum is zero. Exact rational linear programming
returns a solution. If a recorded flow or positively labeled coordinate
becomes zero, the equations imposed on it still imply the defining
equilibrium conditions.
\end{proof}

\chapter{Between security games and matching markets}\label{sec:bridge}

In a market, a positive price signals a full good. In a game with one
resource and unit defender gains, positive coverage signals a target
receiving maximum attack mass. A market good can therefore be represented
by a target if a prescribed amount of attack on that target makes the two
conditions coincide.

\begin{theorem}[Market and game]\label{thm:bridge}
\(\Market\dd\BNE\) with one defender resource and unit defender gains.
Conversely \(\BNE\dd\Market\) for an arbitrary legal budget and weakly
decreasing attacker payoffs.
\end{theorem}

\section{One coverage vector records prices and utilities}
\reviewmark{Private loads, the inverse construction and the slope-preserving price scale.}

We give the construction of \cite[Theorem~3.16]{book} to specify exactly
which parts of the oracle answer later decoders use. Omit isolated market
vertices. Rescale each buyer's entire payoff row by a positive power of
two so all its slopes are at least one. This rescales its utility by that
factor and preserves prices and trades. Use the rescaled \(a_{ij},b_{ij}\)
below. For any rational \(0<\omega\le1\), set
\[
 D=1+\max\bigl(\textstyle\sum_i m_i,\max_j c_j\bigr),\quad
 W=(2+2/\omega)D,\quad A=1+\max(0,\max_{ij}a_{ij}),\quad
 S=1+(|I|+|J|)A .
\]
Use one real target per good, one private target \(g_i\) per buyer, an
outside target \(o\), and a control target \(h\). Normalize the following
type weights by their sum to form the priors.

\begin{itemize}
\item Source \(i\), of weight \(W+m_i\), receives \(a_{ij}/S-b_{ij}x_j\)
at a listed real good, \(A/S-x_{g_i}\) at its private target, and
\(-x_o\) outside.
\item Dedicated type \(j\), of weight \(W-c_j\), gets \(2-x_j\) at its
own real target. The control type, of weight \(W\), gets \(2-x_h\) at \(h\).
\end{itemize}
Each unspecified payoff at \(k\) is \(-2-x_k\). Every weight is between
\(W-D\) and \(W+D\), so the largest prior is at most \(1+\omega\) times
the smallest.

Measure attacks in the original weight units. The control and dedicated
types always use their designated targets. The control load gives a
lower bound \(W\) on maximum load. At least \(W\) of each source type
goes to its private target. If that target is covered, it must have
maximum load, and if uncovered, it strictly attracts its entire source.
Fewer than \(W\) units are left for the outside target across all sources,
so \(x_o=0\) and each source's best payoff is nonnegative. A covered real
target needs a source buyer on it, since its dedicated weight is below
\(W\). Its coverage is therefore at most \(A/S\), using \(b_{ij}\ge1\).
The same bound holds on the private targets. The choice of \(S\) leaves
positive coverage at \(h\), so the maximum load is exactly \(W\).
Every private target is therefore covered and attracts exactly \(W\).

Read the good price as \(P_j=Sx_j\), the rescaled buyer utility as \(S\)
times its source payoff, and \(g_{ij}\) as the source's weight attacking
the real target \(j\). A real target now has at most \(c_j\) units of
source attack, exactly \(c_j\) at positive coverage. A positive-utility
source sends all its remaining \(m_i\) units to real targets. The source
payoff comparisons give the remaining conditions of \eqref{eq:m1}.
Undo the initial row rescaling.

Conversely, from a rescaled market equilibrium take \(x_j=P_j/S\),
\(x_{g_i}=(A-u_i)/S,\ x_o=0\), and assign the remaining coverage to \(h\).
Give \(W\) source units to each private target, the real trades to the
real targets, and each source's unused mass to \(o\). These formulas
recover the complete one-resource game profile and invert the decoder.

\begin{corollary}[Keeping the original slopes]\label{cor:original-slopes}
The market embedding may retain every listed positive attacker slope
\(b_{ij}\), including a slope smaller than one. Use the original offers
and replace \(S\) throughout the game and decoder by
\[
 S'=\frac{1+(|I|+|J|)A}{\gamma},\qquad
 \gamma=\min(1,\min_{ij}b_{ij}).
\]
\end{corollary}
\begin{proof}
The control load, private priority and outside-load bound in the preceding
argument apply at any positive scale. In particular the outside target is
uncovered. A covered real good needs a source with nonnegative payoff
there, hence \(x_j\le A/(\gamma S')\). A covered private target satisfies
\(x_{g_i}\le A/S'\). The sum of both kinds of coverage is at most
\((|I|+|J|)A/(\gamma S')<1\), leaving the control covered. The real
price is \(S'x_j\) and comparison on the real offers recovers the original
row without any rescaling.
Conversely, in any market clearing, a positively priced good has a
positive tight purchase, giving \(P_j\le A/\gamma\). Every buyer value
lies in \([0,A)\). The inverse formulas with \(S'\) therefore give
nonnegative coverages with positive control coverage. Unused source mass
has outside payoff zero, the designated loads are as above, and both
sides best respond. The input minimum and \(S'\) have polynomial binary
length. With no listed edge the all-zero market answer is immediate.
\end{proof}

\section{How the reverse direction uses one oracle answer}
\reviewmark{Earlier adaptive theorem and its conversion to one query in Appendix~\ref{sec:compile}. The cutoff bounds in the projection argument were checked again on November 27, 2026 (UTC).}

For positive defender gain \(d_j\), an optimal defense has a cutoff
\(\lambda\). At interior coverage the attack load on \(j\) equals
\(\lambda/d_j\). At zero coverage it is at most this amount and at full
coverage it is at least this amount. Fixing \(\lambda\) therefore gives
a market with prices in \([0,1]\). Its buyer comparisons are the
attacker comparisons. To represent an equilibrium of the game, its prices
must also sum to the defender budget \(r\).

Here is how the adaptive construction in \cite[Theorem~3.1]{book} handles
the price sum. Assume first that the priors, slopes and defender gains are
positive and \(0<r<n\). A box \([l,h]\) with capacities
\(c_j=\lambda/d_j\) requires full best-response rows, load at least
\(c_j\) at \(x_j>l_j\), and load at most \(c_j\) at \(x_j<h_j\).
Fixing \(l_j=h_j\) leaves that column load free.

For \(\lambda>0\) one ordinary market answers this box query. Give buyer
\(i\) outside value \(v_i^0=\max_k(a_{ik}-b_{ik}h_k)\) and good \(j\)
reserve \(l_j\) by using ordinary intercepts
\(a_{ij}-b_{ij}l_j-v_i^0\) and reading \(x_j=P_j+l_j\).
If \(x_j>h_j\), strict decrease makes every offer on \(j\) worse than
its buyer's outside option; the positive-price column cannot fill.
A buyer with spare mass has value \(v_i^0\) in the original units.
Every maximizer \(k\) in its reserve formula is at its ceiling, since
otherwise its offer would exceed \(v_i^0\). Assign spare mass to one
such \(k\). Only ceiling columns receive the extras, as the box requires.

At \(\lambda=0\) use the all-one corner with any best responses. At
\(\lambda=1+\max_jd_j\) the all-zero corner works because each
capacity exceeds the total attack mass. Bisect, retaining one answer
whose price sum is at least \(r\) and one whose price sum is at most
\(r\). The pattern of each endpoint fixes permitted positive edges,
zero or unit prices, and equality of load to \(\lambda/d_j\) at free
prices. Include the endpoint's sum inequality and the initial cutoff bounds
\(0\le\lambda\le1+\max_jd_j\). These closed constraints define a rational
polytope in \((\lambda,x,f)\). Its projection on \(\lambda\) is an
interval with endpoints of denominator at most a uniform \(Q\) from
the original row data, by the determinant bound of Lemma~\ref{lem:repair}.
When the sampled cutoffs are closer than \(1/Q^2\) the two intervals
meet, since distinct rational endpoints would be at least \(1/Q^2\)
apart. This yields both sum inequalities at one exact cutoff.

To order and then narrow these prices, use the following demand
comparison. On \(D=\{j:x_j>y_j\}\) the decrease from \(x\) to \(y\)
strictly improves every offer, while off \(D\) no offer improves.
Any buyer using \(D\) at \(x\) consequently uses only \(D\) at \(y\).
For two clearings at the same capacities, the box inequalities give the
opposite load comparison on each \(j\in D\). Thus the total loads on
\(D\) agree, and precisely the same buyers send their full mass to \(D\)
at both endpoints. The coordinatewise minimum takes the cheaper
columns of each clearing and their flows. The buyers confined to \(D\)
prefer their \(y\)-choices to their \(x\)-choices; the other buyers can
obtain their \(y\)-value off \(D\) at \(x\). The maximum takes the
dearer columns, where each selected choice keeps its price and no
competitor becomes cheaper. Both patched flows are equilibria. This
gives an ordered bracket \(L\le H\) with \(\sum L_j\le r\le\sum H_j\).

Pin a wide coordinate \(k\) to a dyadic price in the middle third of
\([L_k,H_k]\) and solve the smaller box \([L,H]\) with this pin.
If its load on \(k\) is at least \(c_k\), compare with \(L\) on
\(D=\{j:x_j>L_j\}\). Each column there has load at least \(c_j\) in the
new solution and at most \(c_j\) at \(L\). The demand comparison forces
equality and isolates the buyers on \(D\). Use the new columns on \(D\)
and the old ones elsewhere. If the pinned load is at most \(c_k\),
compare with \(H\) on \(\{j:x_j<H_j\}\) and use the new columns on
that set. At the new prices the other buyers have a best choice off
the set, whose prices agree with \(H\); their old best choices there
remain optimal. Both cases recover a full-box equilibrium at the given
pin. When the cutoff is zero the upper corner of the pinned box already
answers the smaller-box query. Keep the side of the bracket containing
the desired total. Its \(k\)-width falls to at most two thirds of its old value
and other widths cannot grow.

Finally fix the lower endpoint's flow, its zero and unit labels and
the payoff comparisons on its support. Minimize the absolute budget
error over the resulting price LP in \([0,1]^n\). Its rows use only
original data and its positive optimum, if any, is at least \(1/Q'\)
by a second determinant bound. Once every width is at most
\(\eta=1/(4nQ')\), the lower endpoint exhibits error at most \(n\eta\).
The optimum is zero and the repaired prices with the saved flow are a
BNE. There are \(O(n\log(1/\eta))\) narrowing calls, with pins on a
fixed dyadic grid finer than \(\eta/4\). To bound all query lengths,
first replace each returned price by a vertex on its support and box
labels, then patch. In a tree of tight comparisons every price is
an affine function of one root, with coefficients from the original
payoff data along a path of length at most \(n-1\). At a vertex the root
is either one box endpoint or pin, or is fixed by a further tight
comparison with original coefficients. Thus each round adds polynomially
many bits to the maximum previous endpoint length. The corresponding
flow vertex depends only on margins and the fixed cutoff.

Endpoint budgets have fixed coverages. Zero-prior types can be restored
with a conditional best response. For zero slopes or gains add a common
positive dyadic \(\epsilon\) to those payoff coefficients (adding it to
all slopes and gains also works). In the original game the perturbed
support and price labels have attacker comparison error at most
\(2\epsilon\) and defender score error at most \(\epsilon\), since
\(x_j\in[0,1]\) and each load is at most one. Relax the original closed
pattern by one slack in these inequalities, leaving rows and budget
exact. The same basis argument bounds positive optimum slack below by
\(1/Q_0\) for an input-computable integer \(Q_0\) of polynomial bit
length. Taking \(2\epsilon<1/Q_0\)
recovers an exact original-game solution by LP. These are the pattern
and bit bounds of \cite[Section~3.1]{book}.

Appendix~\ref{sec:compile} isolates each query, encodes its comparisons
in goods prices, and uses those bits to select the next query. Every
equilibrium of the single market yields a valid transcript and a full game
profile, proving the reverse direction of Theorem~\ref{thm:bridge}.

\chapter{From matching to traffic}\label{sec:forward}
\reviewmark{Capacity filler, strict road costs and both support-recovery LPs.}

To read a price from a traffic equilibrium, give each good its own road
and add classes of auxiliary travelers. The auxiliary travelers absorb
capacity left by original buyers. On a used road, a small change in load
then represents the good's price. At the end a support LP removes the
small errors introduced by making every road cost strictly increasing.

Choose a positive dyadic \(\tau\) with \(\tau P_*<1/(2Q_{\rm M})\),
where \(Q_{\rm M}\) comes from the matching input in Lemma~\ref{lem:repair}.
Create one road per good and a common outside road \(o\). Each original
buyer becomes a class with costs
\[
 \ell_{ij}(X_j)=-a_{ij}+b_{ij}(X_j-c_j)/\tau,\qquad \ell_{io}=0 .
\]
For each \(j\), a further class of mass \(c_j\) can use road \(j\) at cost
\((X_j-c_j)/\tau\), or \(o\) at cost zero. This class fills unused capacity
while the price is zero.

At equilibrium \(X_j<c_j\) is impossible, since then the added class alone
would send \(c_j\) there. Put \(P_j=(X_j-c_j)/\tau\). If \(P_j>0\), the
added class uses only \(o\), so the original buyers supply
\(X_j=c_j+\tau P_j>0\). Some original buyer weakly prefers \(j\) to the
outside option, giving \(P_j\le P_*\). If \(P_j=0\), original traffic on
the road is at most \(c_j\). The original flows use best offers at prices
\(P\), and a buyer with strictly positive best utility uses its entire mass
on real goods. Lemma~\ref{lem:repair} recovers an exact matching equilibrium
from the support and labels, since every column error is at most \(\tau P_*\).

It remains to obtain the intermediate equilibrium from an oracle with the
strict costs and full menus of Theorem~\ref{thm:traffic}. Compute
\(Q_{\rm W}\) for this intermediate instance and write \(T'\) for its total
mass, including the added classes. Add \(\zeta X_j\) to every allowed cost,
including that on \(o\), where \(0<\zeta<1/[2(1+T')Q_{\rm W}]\) is dyadic.
For each class add the same constant to all allowed intercepts to make them
nonnegative. On each missing road give the class a positive slope and an
intercept greater than its maximum allowed cost on \([0,T']\). No traveler
can use such a road at equilibrium. Removing the common shift and then the
perturbation leaves every used road within \(\zeta T'<1/Q_{\rm W}\) of
minimum cost for its class. The traffic LP first recovers an exact
intermediate equilibrium, and only then do we form the prices \(P_j\) and
apply the matching LP. Both dyadic precisions require polynomially many
bits. This proves \(\Market\dd\Traffic\).

\chapter{From traffic to matching}\label{sec:device}

The reverse reduction needs a good whose price a traveler can read as a
function of the mass assigned to its road. To achieve accuracy \(\delta\),
a direct table of prices could require \(T/\delta\) entries. We instead
use a network with \(O(1+\log(1+T)+\log(1/\delta))\) arcs. Its minimum
cost as a function of flow has one slope for each price level.

\section{A flow network with consecutive marginal costs}
\reviewmark{Inner-flow recurrence, its feasible interval and slope sequence.}

Let \(0<\delta\le1\) be dyadic and choose the least \(k\ge1\) with
\(K=2^k\delta>T+\delta\). The network \(G_0\) has a single zero-cost arc
\(s_0\to t_0\) of capacity \(\delta\). To obtain \(G_h\), add a source
\(s_h\), a sink \(t_h\), and the four arcs in Figure~\ref{fig:count}.
Each has capacity \(c=2^{h-1}\delta\). The two positive fees both equal
\(w_h=(2^h-1)/2\). Thus \(G_k\) is acyclic and has \(4k+1\) arcs.

\begin{figure}[ht]
\centering
\begin{tikzpicture}[>=stealth,thick,every node/.style={font=\small}]
  \node (s) at (0,0) {$s_h$};
  \node (a) at (2.7,.88) {$s_{h-1}$};
  \node (b) at (2.7,-.88) {$t_{h-1}$};
  \node (t) at (5.4,0) {$t_h$};
  \draw[->] (s) -- node[above,sloped] {$0$} (a);
  \draw[->] (s) -- node[below,sloped] {$w_h$} (b);
  \draw[->] (a) -- node[above,sloped] {$w_h$} (t);
  \draw[->] (b) -- node[below,sloped] {$0$} (t);
  \draw[->] (a) -- node[right] {$G_{h-1}$} (b);
\end{tikzpicture}
\caption{One level of the Disser--Skutella network. The four outer arcs
have equal capacity. The vertical arrow stands for the entire preceding
network.}\label{fig:count}
\end{figure}

\begin{lemma}\label{lem:slopes}
Let \(F_h(y)\) be the minimum cost of a flow of value \(y\) in \(G_h\).
Its domain is \([0,2^h\delta]\). On successive open intervals of width
\(\delta\), its slopes are \(0,1,\ldots,2^h-1\).
\end{lemma}
\begin{proof}
The base arc gives the statement for \(h=0\). Write \(z\) for the flow
through the inner network. At outer value \(y\in[0,2c]\) the two costly
arcs together carry \(y-z\), and the two free arcs together carry \(y+z\).
Thus \(0\le z\le\min(y,2c-y)\). Every such \(z\) is feasible by placing
\((y-z)/2\) on each costly arc and \((y+z)/2\) on each free arc.
Consequently
\[
 F_h(y)=\min_{0\le z\le\min(y,2c-y)}
             \{w_h(y-z)+F_{h-1}(z)\}.
\]
The largest inner slope is \(2^{h-1}-1<w_h\), so the largest feasible \(z\)
is optimal. The first half keeps the inner slopes. The second half has the
slopes \(2w_h-(2^{h-1}-1),\ldots,2w_h\) in that order.
\end{proof}

\section{Representing the marginal cost by a market price}
\reviewmark{Residual assignment dual, cut argument and prices at breakpoints.}

For one road construct a set of private buyers and goods and a single
\emph{entry good} \(p\). Only \(p\) is offered to the traffic classes.
Let \(X\in[0,T]\) denote their total purchases of \(p\), and \(\pi\) its
price. Put \(s=s_k\), \(t=t_k\), \(C=T+\delta+1\), and write
\(\kappa_e\) for an arc's fee. Define
\[
 \kappa'_e=\kappa_e+\mathbf1_{\{\operatorname{head}(e)=t\}},
 \qquad M=4+\sum_e\kappa'_e,\qquad M_p=M+1.
\]
Each arc \(e\) becomes a good of capacity \(\caparc(e)\), available only
to its two endpoints. All private offers have the form
\(\text{intercept}-\text{price}\). The available offers are as follows.

\begin{center}
\small
\begin{tabular}{@{}lll@{}}
\toprule
Buyer & Mass & Goods and their intercepts\\
\midrule
Internal node \(v\) & Sum of incoming capacities &
  Incoming \(M\), outgoing \(M-\kappa_e\)\\
Source \(s\) & \(C\) & Outgoing \(M-\kappa_e\), entry \(p\) at \(M_p\)\\
Sink \(t\) & \(K\) & Incoming \(M+1\), discard good at \(M\)\\
Additional buyer & \(\delta\) & Entry \(p\) at \(M_p+1\) only\\
\bottomrule
\end{tabular}
\end{center}
The entry good has capacity \(C\), and the private discard good has
capacity \(K+1\). The latter always has price zero.

\begin{lemma}[Price as a function of load]\label{lem:price}
In every equilibrium of a market containing this construction, provided
the external buyers can purchase only \(p\) among these goods and contribute
at most \(T\), the entry price satisfies
\[
                         X\ \le\ \delta(\pi-2)\ \le\ X+\delta .
\]
\end{lemma}
\begin{proof}
If \(\pi>M_p\), the source buys none of \(p\), and all other buyers together
can supply at most \(T+\delta<C\), contradicting clearance at a positive
price. Hence the additional buyer strictly prefers to trade and buys
\(\delta\). Set \(r=X+\delta\). In particular
\(0<r<\min(C,K)\), including at the endpoints \(X=0,T\).

Remove all buyers at \(p\) except the source, keeping the other private node
buyers, and reduce its capacity to \(u=C-r\). The private allocations and
prices remain an equilibrium of this residual market. Since every remaining
offer has unit slope, these are primal and dual optima for the assignment
LP maximizing the sum of intercepts times trades, subject to buyer mass and
good capacity upper bounds. Indeed, the dual variable of a buyer is its
value, the dual variable of a good is its price, and~\eqref{eq:m1} is
complementary slackness. Write \(V(u)\) for the optimal value. We compute
it for every \(r\) in the open interval above.

\emph{Private arc goods fill.}
Take any optimal primal and dual pair for the residual market. A slack arc
good has price zero. At its head, the purchase of that arc is smaller than
its capacity. More generally, consider any zero-price incoming arc with this
property. If its head is internal, its buyer attains its largest possible
utility \(M>0\), so must spend its full mass. Its incoming purchases alone
cannot do so. It buys some outgoing arc of zero fee and zero price. At the
next head that arc again has head purchase below capacity. At the sink, a
zero-price incoming arc gives utility \(M+1\), strictly above the discard
offer. The sink must buy its full mass \(K\) on incoming arcs, whose total
capacity is exactly \(K\). This forces every head purchase there to equal
capacity. Acyclicity now contradicts the existence of the first slack arc.

\emph{Welfare of the residual market.}
Write \(f_e\) for the tail's purchase of arc \(e\). Its head buys
\(\caparc(e)-f_e\). Let \(y\) be the source's entry purchase and \(z\) the
sink's discard purchase. After subtracting a constant independent of \(u\),
the objective is
\[
                M_p y-\sum_e\kappa'_e f_e+Mz .
\]
The remaining restrictions are \(0\le f_e\le\caparc(e)\),
\(0\le y\le u\), \(y+\operatorname{out}(s)\le C\),
\(\operatorname{out}(v)\le\operatorname{in}(v)\) internally, and
\(0\le z\le\operatorname{in}(t)\). These conditions are also sufficient
for an allocation with full arc goods. In an optimum
\(z=\operatorname{in}(t)\).

First \(y=u\). Otherwise use spare source mass to increase \(y\), or, if
the source is full, follow a positive-flow path from it to a vertex with
surplus inflow or to the sink. Reduce a small amount of flow along this
path and assign the freed mass to \(p\). If the path ends at the sink,
reduce \(z\) by the same amount. The sacrifice per unit is at most \(M\)
and the entry purchase gains \(M_p>M\). This proves the assertion, leaving
\(\operatorname{out}(s)\le r<K\).

There can be no internal surplus inflow. From a surplus vertex the residual
network reaches \(s\) by tracing a positive-flow path backwards. If it could
not reach \(t\), its reachable set would contain \(s\) and exclude \(t\).
Every arc leaving the set would be saturated and every arc entering it
empty. Its cut capacity would equal its net outflow, at most
\(\operatorname{out}(s)\le r<K\), since internal nodes have no net outflow.
This contradicts the feasible flow of value \(K\) from Lemma~\ref{lem:slopes}.
Push a small amount on a simple residual path from the surplus vertex to
\(t\), increasing \(z\). The gain per unit is at least
\(M-\sum_e\kappa'_e>0\). The same cut argument starting at \(s\) shows
that \(\operatorname{out}(s)<r\) would allow a profitable push using spare
source mass. Thus the optimum conserves at every internal node and sends
exactly \(r\) to \(t\). Conversely, every feasible \(s\)-\(t\) flow of
value \(r\) yields a residual allocation with \(y=u\) and \(z=r\).
Its minimum fee sum is \(F_k(r)\), and the extra unit in \(\kappa'_e\)
is paid once per unit entering \(t\). We have proved
\begin{equation}\label{eq:value}
 V(C-r)=\mathrm{constant}+M_p(C-r)+(M-1)r-F_k(r).
\end{equation}

\emph{Price at every dual optimum.}
The dual with entry multiplier \(\pi\) stays feasible if the capacity
changes to \(u'\). Hence \(V(u')\le V(u)+\pi(u'-u)\). Substituting
\eqref{eq:value} and varying \(r\) on both sides gives
\(\pi=M_p-M+1+\mu=2+\mu\) for some \(\mu\in\partial F_k(r)\), the interval
of supporting slopes. In particular when \(r/\delta\in(l,l+1)\) we have
\(\mu=l\), and at an interior breakpoint \(r/\delta=l\) we have
\(\mu\in[l-1,l]\). Both cases give \(r-\delta\le\delta\mu\le r\).
Since \(r=X+\delta\), this is the desired inequality.
\end{proof}

\section{Recovering exact class flows}
\reviewmark{Uniform error, exact recovery and the two-road example.}

For a traffic instance, compute its \(Q_{\rm W}\) and put \(B_*=\max_{ij}B_{ij}\).
Choose a positive dyadic \(\delta\le1\) with \(2B_*\delta<1/Q_{\rm W}\).
Make one independent copy of the preceding market construction per road.
For each class add a buyer of mass \(h_i\) whose offers on the entry goods
are
\[
 R_i-A_{ij}-B_{ij}\delta(\pi_j-2),\qquad
 R_i=1+\max_j\{|A_{ij}|+B_{ij}(T+1)\}.
\]
These offers have strictly positive price slopes \(B_{ij}\delta\) in
the matching convention \(a-bP\). The a priori bound \(T\) suffices for
Lemma~\ref{lem:price} even before establishing that the class buyers spend
all their mass. That lemma now implies that all displayed offers are
strictly positive. Every class spends its entire mass.

Read the class purchases as \(f_{ij}\). The cost used by the market lies
between the true cost \(A_{ij}+B_{ij}X_j\) and that cost plus \(B_*\delta\).
For a used road \(j\) and any competing road \(k\), the true cost on \(j\)
is at most the displayed cost on \(j\), hence at most the displayed cost
on \(k\). The excess relative to the true cost on \(k\) is therefore at
most \(B_*\delta<1/Q_{\rm W}\). The traffic LP recovers exact class flows.
The depth \(k\) is polynomial in the input length, as are the bit lengths of
all fees, capacities and offers. This proves \(\Traffic\dd\Market\) and
completes the proof of Theorem~\ref{thm:traffic}.

\paragraph{An example of support recovery.}
Let classes of masses \(1/4\) and \(3/4\) have costs
\((X_L,2+X_R)\) and \((2X_L,1/4+X_R)\) on two roads. The first class uses
only \(L\). If the second uses both roads, its indifference equation and
\(X_L+X_R=1\) force \(X_L=5/12\), \(X_R=7/12\). The second class sends
\(1/6\) to \(L\) and \(7/12\) to \(R\). This calculation illustrates what the
decoder takes from the market. Knowing which roads carry each class is
enough to reconstruct the exact loads without rounding their numerical values.

\chapter{Matching and two-player hide-and-seek}\label{sec:matching-links}
\reviewmark{Both normalizations and the treatment of nonpositive offers.}
The hider selects a listed buyer--good pair \(e=ij\). The seeker inspects
one vertex. A miss pays both players zero. On a catch the seeker receives
\(r_v>0\), which depends on the inspected vertex \(v\), and the hider
loses \(l_{ev}>0\), which may depend on both the pair and the vertex.
Write \(\mathsf{Hide}\) for finding a Nash equilibrium of this explicit
family. Galichon and Jacquet give a matching interpretation of generalized
hide-and-seek~\cite[Theorem~1]{gj}. The following normalization makes both
directions on binary input explicit.

\begin{theorem}\label{thm:hide}
\(\mathsf{Hide}\mathrel{\equiv_{\rm m}}\Market\).
\end{theorem}
\begin{proof}
Given a market, discard edges with \(a_{ij}\le0\). An equilibrium on the
remaining edges extends by zero flow, since \(u_i,P_j\ge0\). If no edges
remain, the all-zero profile suffices. Otherwise set
\[
  r_i=1/m_i,\quad r_j=1/c_j,\qquad
  l_{ei}=1/(m_i a_{ij}),\quad l_{ej}=b_{ij}/(c_j a_{ij}).
\]
Conversely, from these game data read \(m_i=1/r_i,\ c_j=1/r_j,\)
\(a_{ij}=1/(m_i l_{ei})\) and \(b_{ij}=c_jl_{ej}/(m_i l_{ei})\).
From a market solution form \(F=\sum_e g_e\) and
\(D=\sum_i m_i u_i+\sum_j c_jP_j\). Both are positive: a positive offer
forces a positive utility or price, whose margin must fill. Hide on \(e\)
with probability \(g_e/F\), and inspect \(i,j\) with respective
probabilities \(m_i u_i/D,c_jP_j/D\). The loss at \(ij\) is
\((u_i+b_{ij}P_j)/(a_{ij}D)\ge1/D\), with equality on used edges.
The payoff from inspecting any vertex is at most \(1/F\), with equality
on inspected vertices.

Conversely, write \(p_e,q_v\) for equilibrium probabilities and \(\ell,h\)
for the expected loss and reward. Some vertex meets a positive-probability
edge, so the best inspection reward \(h\) is positive. Thus there is a
positive probability of a catch and \(\ell>0\). Set
\[
   g_e=p_e/h,\qquad u_i=q_i/(m_i\ell),\qquad P_j=q_j/(c_j\ell).
\]
The hider's best-response inequalities are precisely the offer
inequalities, with equality on positive \(g_e\). The seeker's are precisely
the margin bounds, tight on positive utilities and prices. All transforms
use polynomial-bit rational arithmetic.
\end{proof}

\chapter{Fixed-total envy-free rent division}\label{sec:rent}
\reviewmark{Interior coverage, equal loads and the Hall matching decoder.}

There are \(n\) tenants and \(n\) rooms. Tenant \(i\)'s utility for room
\(j\) at rent \(P_j\) is \(v_{ij}-c_{ij}P_j\), where \(c_{ij}>0\).
The output is a bijective assignment and rents summing to an input value
\(R\), such that each tenant receives a favorite room at those rents.
Rents may be signed. This is the affine fixed-total variant of the
envy-free rent-division problem discussed by Arunachaleswaran, Barman and
Rathi~\cite[\S6]{abr}. Unlike \eqref{eq:m1}, the rent sum is prescribed
and every tenant receives one room.

\begin{theorem}\label{thm:rent}
Fixed-total affine envy-free rent division reduces many-one to \(\BNE\)
with one resource, equal type priors and unit defender gains.
\end{theorem}
\begin{proof}
For \(n=1\) assign the room at rent \(R\). Otherwise put
\[
 P_0=R/n,\quad h_{ij}=v_{ij}-c_{ij}P_0,\quad
 K=\frac{n(\max h-\min h+1)}{\min_{ij} c_{ij}} .
\]
Create one target per room and one type per tenant, with prior \(1/n\).
Give type \(i\) payoff \(h_{ij}-Kc_{ij}(x_j-1/n)\) at coverage \(x_j\).
This equals the tenant's utility at \(P_j=P_0+K(x_j-1/n)\). The one
resource gives \(\sum_j x_j=1\), so these rents always sum to \(R\).

Suppose an equilibrium had \(x_\ell=0\). Every type would receive at
least \(\min h+K\min c/n>\max h\) there. Some \(k\) has \(x_k\ge1/n\),
where every type receives at most \(\max h\). Thus \(k\) has positive
coverage and zero attack load. With unit gains a positively covered
target must have maximum load, and the total load is one. This
contradiction proves \(x_j>0\) for all \(j\). Their loads are therefore
all \(1/n\).

Let \(f_{ij}\) be the joint attack mass from the returned equilibrium.
The row and column sums of \(nf\) are one. For any set of rows its
positive neighbors can accommodate their mass only if they are at least
as numerous. Hall's theorem gives a perfect matching on positive entries,
which a bipartite matching algorithm finds. Every matched edge was a best
response and hence gives its tenant a favorite room at the quoted rents.
All coefficients and recovery steps have polynomial binary cost.
\end{proof}

For example, if every tenant ranks two identical rooms solely by rent,
equal loads can be realized by each tenant splitting evenly over both
rooms. The decoder may return either perfect matching. It uses the
positive support because the output of the rent problem is an assignment
even when the equilibrium supplied to the decoder is fractional.

\chapter{Updating a market and selecting its prices}\label{sec:update}
\reviewmark{Unit-entry construction and decoding after a margin update.}

A clearing allocation can be easy to write down before an update. What
information about the next equilibrium does it give? Even an update of one
unit of buyer mass can encode an independently specified market.

\begin{theorem}[One additional unit]\label{thm:update}
A market with integer masses and capacities and rational offers is
supplied with its full equilibrium at its least
prices. One existing buyer receives one additional unit of mass. Finding
any new equilibrium is many-one equivalent to \(\Market\), even when the
supplied prices are all zero, the flow at those prices is unique and every
margin fills. The equivalence also holds when the update first doubles all
margins and then adds the unit.
\end{theorem}
\begin{proof}
Start with a nonempty market and choose, for each buyer,
\[
 t_i=4\max\bigl(\{1\}\cup\{b_{ij}:ij\text{ listed}\}\bigr),\qquad
 A_{ij}=a_{ij}/t_i,\qquad B_{ij}=b_{ij}/t_i .
\]
Thus \(B_{ij}\le1/4\). Take \(H=1+\max(\{0\}\cup\{A_{ij}\})\).
Multiply the original margins by twice their common denominator \(s_0\),
writing \(s=2s_0\), \(M_i=sm_i,\ C_j=sc_j,\ M=\sum_iM_i,\ C=\sum_jC_j\).
All margins are even integers. Add a good \(D\) of capacity \(M\) and
a buyer \(d\) of initial mass \(C\). To specify the offers as functions
of nonnegative physical prices, write
\[
 q_j=P_j-4H,\qquad q_D=P_D-3H.
\]
Each original row has offer \(A_{ij}-B_{ij}q_j\) on each of its real
options and \(-q_D\) on \(D\). Row \(d\) offers \(-q_j\) on every real
good and \(-q_D\) on \(D\). Each buyer retains the outside utility zero.

At physical prices zero the original rows strictly choose \(D\), since
every real offer is less than \(2H\) and \(D\) pays \(3H\).
Buyer \(d\) strictly prefers the real goods and fills all of them.
The allocation is determined and the all-zero price vector is least.
For the doubling formulation, supply these flows and margins divided by
two. The intermediate doubled table is exactly the table just described.

Increase the mass of \(d\) to \(C+1\). In any new equilibrium \(q_D<0\)
would force every buyer to spend in full, exceeding total capacity by
one. If \(q_D>0\), the positively priced good \(D\) would receive no
buyers. Thus \(q_D=0\). If any \(q_j<0\), buyer \(d\) would have to spend
its entire \(C+1\) on real goods of total capacity \(C\). Consequently all
\(q_j\ge0\) and the dummy buyer has utility zero.

Use \(q_j\) as the original prices. Every original buyer with positive
utility spends its full mass on real goods. Every real good with
\(q_j>0\) is physically positively priced and is rejected by \(d\), so
it fills from real buyers. Divide real trades by \(s\) and multiply
each real buyer's utility by \(t_i\). These are precisely the
conditions of \eqref{eq:m1}, for every returned answer. The reverse
reduction submits the updated market. All constants have polynomial
binary length.
\end{proof}

\section{Least and greatest clearing prices}
\reviewmark{Circulation cuts and both mass perturbations.}

An equilibrium returned by a market oracle can have many alternative
prices. A single perturbation makes the desired extreme independent of
which answer the oracle selects.

\begin{theorem}\label{thm:extreme}
Finding a clearing allocation at the componentwise greatest equilibrium
prices, or at the componentwise least equilibrium prices, is many-one
equivalent to \(\Market\).
\end{theorem}
\begin{proof}
Fix posted prices and evaluate each buyer's maximum over its options and
outside utility zero. Let \(T_i\) be the goods attaining this maximum,
\(I^+\) the buyers with positive maximum, and \(J^+\) the positive-price
goods. A clearing flow exists at those prices precisely when
\begin{align*}
 \sum_{\{i\in I^+:T_i\subseteq S\}}m_i&\le c(S)
     &&(S\subseteq J),\\
 c(S)&\le\sum_{\{i:T_i\cap S\ne\varnothing\}}m_i
     &&(S\subseteq J^+).
\end{align*}
Indeed, form a circulation from the source through buyers, their tight
goods, and the sink, with an unbounded return arc. Row and column arcs
have their prescribed upper bounds, which are also lower bounds for
\(I^+\) and \(J^+\). The tight arcs have infinite capacity. For a cut
containing neither terminal the first inequality bounds the incoming
mandatory rows. For a cut containing both, the second bounds the
mandatory goods in the complement. The other two cut types have zero
incoming lower bound or unbounded outgoing capacity. Circulation
feasibility gives both the characterization and an allocation algorithm.

Let \(D\) be a common denominator of the margins. With \(m\) buyers,
increase every buyer mass by \(1/(2mD)\). Take any price vector \(P'\) from
an equilibrium of this table. Restoring the original masses makes the
first inequalities easier. A deficit in the second would be at most
\(1/(2D)\) and would be a multiple of \(1/D\). Hence \(P'\) clears the
original market. If an original equilibrium price \(P\) were higher on a
nonempty set \(S=\{j:P_j>P'_j\}\), all of \(S\) would fill at \(P\).
Every buyer then using \(S\) strictly improves there at \(P'\) and can
only buy in \(S\). Their original mass was already at least \(c(S)>0\).
Their increased mass cannot fit. Thus \(P'\) is the greatest original
price vector.

For the least prices, increase each of the \(n\) capacities by \(1/(2nD)\).
On restoring the old capacities the second inequalities become easier
and the first can have no positive deficit by the same denominator bound.
If \(P'\) exceeded an original price \(P\) on a nonempty set \(S\),
the buyers filling the enlarged capacities of \(S\) at \(P'\) would
strictly prefer \(S\) at \(P\), exceeding the old capacities. This proves
leastness. In either construction recover an original allocation by the
circulation above. Any extremal solution itself answers an ordinary market
query, proving the other direction.
\end{proof}

For unit buyer masses and unit capacities this recovery also gives an
integral assignment at the chosen prices. The circulation has integral
lower and upper bounds and one can take an integral flow. This applies
with the outside options and reserve prices of \eqref{eq:m1}.

\section{Starting from the greatest old prices}
\reviewmark{Predecessor price interval and both successor guarantees.}

\begin{corollary}\label{cor:update-greatest}
Suppose a full equilibrium at the greatest old prices is supplied.
Finding any equilibrium after two units are added to one buyer is
many-one equivalent to \(\Market\). Finding a greatest-price equilibrium
after one added unit is also many-one equivalent to \(\Market\).
\end{corollary}
\begin{proof}
Use the table of Theorem~\ref{thm:update}, starting with dummy mass
\(C-1\). The integer scaling ensures \(C\ge2\). Let
\(V=\max(\{0\}\cup\{A_{ij}+4HB_{ij}\})<2H\).
The possible predecessor prices are exactly
\[
 P_j=0\text{ on real goods},\qquad 0\le P_D\le3H-V .
\]
Indeed, if all real prices were positive, all real goods would fill,
and the total mass shortfall gives \(P_D=0\). But then every real row
strictly prefers \(D\), leaving at most \(C-1\) real units to fill \(C\).
Thus some real price is zero and \(d\) can use only zero-price real
goods. If \(P_D>0\), filling its capacity \(M\) takes every real row's
entire mass. The same allocation of real rows is forced by strict
preference when \(P_D=0\). No positive-price real good could fill.
With all real prices zero, the row comparisons give the displayed
interval, and filling all real capacities except one unit from \(d\)
realizes every price in it. This gives an explicit greatest predecessor.
At mass \(C+1\) the first theorem's decoder works unchanged.

At mass \(C\) no clearing has \(P_D>3H\). Every ordinary source
equilibrium with real prices \(p_j\) supplies one with \(P_D=3H\)
and \(P_j=4H+p_j\): put unmatched real rows on \(D\), fill the real
vacancies from \(d\), and put the remaining dummy mass on \(D\).
Thus the greatest successor prices attain this known value of \(P_D\)
and satisfy \(P_j\ge4H\) on every real good. The projection of
Theorem~\ref{thm:update} now applies. Conversely, an unrestricted
market call handles the two-unit update, and greatest-price selection
in Theorem~\ref{thm:extreme} handles the one-unit update.
\end{proof}

\chapter{Fixing prices and utilities on both sides}\label{sec:pins}
\reviewmark{Raised reservations, exchange of allocations and the extension test.}

An exchange of roles in \(\Market\) is exact. Divide
\(u_i+b_{ij}P_j\ge a_{ij}\) by \(b_{ij}\). In the transposed market,
good \(j\) becomes a buyer of mass \(c_j\), the old buyer \(i\) is a good
of capacity \(m_i\), and the new intercept and slope are
\(a_{ij}/b_{ij},1/b_{ij}\). Transposing the trades exchanges the margin
conditions and gives utility \(P_j\) and price \(u_i\). This also works
with buyer outside utilities \(o_i\) and good reserves \(\ell_j\), which
exchange roles. Such reservations can be translated to zero by shifting
\(u_i,P_j\) and replacing \(a_{ij}\) by \(a_{ij}-o_i-b_{ij}\ell_j\).

\begin{theorem}\label{thm:pins}
Given a market together with prescribed utilities \(s_i\) on buyers
\(i\in A\) and prescribed prices \(t_j\) on goods \(j\in B\), one query
on the remaining buyers and goods and a transportation test decide if the
pins extend to an equilibrium. On a yes instance they return an extension.
Any residual equilibrium can be used for the query answer.
\end{theorem}
\begin{proof}
First check reservations on the pinned vertices and
\(s_i+b_{ij}t_j\ge a_{ij}\) on every pair with both ends pinned.
Delete \(A,B\), leaving each remaining buyer and good its full original
margin. Raise their respective reservations to
\begin{align*}
 R_i&=\max\bigl(\{o_i\}\cup\{a_{ij}-b_{ij}t_j:j\in B,\ ij\text{ listed}\}\bigr),\\
 L_j&=\max\bigl(\{\ell_j\}\cup
           \{(a_{ij}-s_i)/b_{ij}:i\in A,\ ij\text{ listed}\}\bigr).
\end{align*}
These lower bounds enforce every comparison against a pinned option.
A full equilibrium extending the pins restricts to a residual one: a
remaining row strictly above \(R_i\) cannot use a pinned good or the
outside, and a remaining column strictly above \(L_j\) cannot receive
a pinned buyer or have spare capacity.

Suppose such a full witness exists, with residual prices \(p^*\) and
flow \(g\). Compare it to any returned residual prices \(p\) and flow
\(h\). Let \(D=\{j:p_j<p_j^*\}\). Every good of \(D\) fills in \(g\).
All its serving buyers strictly prefer \(D\) at \(p\) and must spend
their entire mass there. Since the total capacity is fixed, equality
forces exactly those buyers to fill exactly those goods in both flows.
Exchanging the two price vectors gives the same isolated structure on
\(\{j:p_j>p_j^*\}\). These two buyer groups are disjoint. Their entire
masses were spent on interior goods even in the full witness.

Replace those buyers' rows in the full witness by their returned rows
from \(h\). All changed-price goods still fill. The remaining buyers
use only unchanged-price interior goods, their residual outside
options, or the pinned boundary in the full witness. Their best
residual utilities coincide: each attains its maximum on the unchanged
menu or residual outside in both profiles. The retained purchases
therefore remain optimal. All comparisons on unused boundary edges
hold by \(R_i,L_j\). This proves extension at the returned potentials,
even preserving every boundary-edge quantity of the witness.

To compute the extension one does not need that witness. Fix the
returned prices and utilities and the pins, keep all original tight
pairs, and test a transportation circulation with the \emph{original}
reservations and fill requirements. The extension argument proves that
a failure certifies infeasibility of the pins. All data have polynomial
binary length. Empty residual sides use the remaining reservations and
zero trade.
\end{proof}

\chapter{Splitting edge weights at a pinned price}\label{sec:pinned}

Consider a graph in which each edge can distribute its weight between its
two endpoints. The vertices announce prices. An edge sends weight to the
endpoint making the better offer and may split its weight if the two offers
tie. The prescribed load at every vertex must be met exactly. One vertex
price is fixed in advance. This chapter relates the joint choice of the
prices and the split to equilibrium search.

Let \(G=(V,E)\) be a loopless multigraph. An edge \(e=vw\) has weight
\(m_e>0\) and offer \(a_{ev}-b_{ev}P_v\) at endpoint \(v\), with
\(b_{ev}>0\). Each vertex has a required load \(c_v>0\). The input specifies
\(q\in V\) and \(t\in\mathbb Q\). A solution gives real prices \(P_q=t\),
\(P_v\in\mathbb R\) elsewhere, and nonnegative endpoint amounts \(z_{ev}\)
with
\[
 z_{ev}+z_{ew}=m_e,\qquad \sum_{e\ni v}z_{ev}=c_v,\qquad
 z_{ev}>0\ \Longrightarrow\
 a_{ev}-b_{ev}P_v\ge a_{ew}-b_{ew}P_w .
\]
The input promise is that the two load equations admit a nonnegative
solution. A transportation feasibility test checks this promise.

\begin{theorem}\label{thm:pinned}
Finding a price vector and a split on this domain is many-one equivalent to
\(\Market\) and hence to \(\BNE\). The equivalence holds even on graphs of
maximum degree three with slopes in \(\{1,2\}\), where the edges whose two
slopes differ have pairwise disjoint endpoints.
\end{theorem}

\section{Turning market rows into edges}
\reviewmark{Dummy normalization, private vertices and degree-three port paths.}

Use Lemma~\ref{lem:rectangle} to obtain a market with slopes one and two.
Write \(M=\sum_i m_i\) and \(C=\sum_j c_j\). Add a good \(D\) of load
\(M+1\) and a buyer \(d\) of prescribed load \(C+1\). For this intermediate
table require \emph{all} margins to fill. Write \(p\) for its prices and
\(u\) for its row maxima. Set \(u_d=0\). Give \(d\) offer \(1-p_j\) on
every good, including \(D\). A real buyer \(i\) has offer \(2-p_D\) on
\(D\) and \(1+a_{ij}+b_{ij}-b_{ij}p_j\) on a listed real good.

Every source equilibrium lifts to these margins at \(p_j=1+P_j\),
\(u_i=1+u_i^{\rm old}\), and \(p_D=1\). Send each real buyer's unused
mass to \(D\), fill each real good's vacancy from \(d\), and give \(D\)
the remaining \(1+\sum_{ij}g_{ij}\) units from \(d\). Conversely, in
every solution \(d\) uses \(D\), since all real buyers together have
only \(M\) units. Its used offer gives \(p_D=1\). Its other comparisons
give \(p_j\ge1\), and the real buyers' comparison with \(D\) gives
\(u_i\ge1\). A real row with \(u_i>1\) uses only real goods and a real
column with \(p_j>1\) receives only real trades. Subtracting one from the
real prices and utilities and retaining real trades yields \eqref{eq:m1}.

For each augmented row \(i\) of prescribed load \(r_i\) and with \(k_i\)
options, create a private vertex \(h_i\). Put one edge of weight
\(K_i=r_i+1\) between \(h_i\) and each option \(j\). The offer on the
goods side is its augmented offer. On the private side use \(A-P_{h_i}\)
with a common rational \(A\), for example zero. Prescribe load
\(k_iK_i-r_i\) at \(h_i\) and pin \(P_{h_d}=A\).
Each edge must send positive weight to \(h_i\), because even \(k_i-1\)
whole edges cannot fill it. Consequently \(A-P_{h_i}\) dominates every
offer to this row, with equality at every option receiving weight.
Exactly \(r_i\) units leave the private vertex's side of these edges.
This recovers all augmented rows with \(u_i=A-P_{h_i}\). To lift, place
each trade on its goods side and its edge's remaining mass on the private
side. Taking \(P_{h_i}=A-u_i\) satisfies every comparison.

\paragraph{Degree three.}
Replace a vertex of degree \(k\) and load \(c\) by a path of \(k\) ports,
each incident to exactly one of its original edges. Between adjacent ports
put a new edge with equal offers \(-P\) at its two ends. Work from the leaf
toward the root. At each nonroot port let \(U\) be the sum of its external
edge weight and its child edge weight, if any. Prescribe load \(U+1\) and
give its parent edge weight \(U+2\). Prescribe load \(c+k-1\) at the root.
The parent edge must send positive mass to the port and is too large to
fit there in full. Its two port prices therefore agree. Summing the port
loads and subtracting the internal edge weights restores \(c\).
Conversely, if \(z\) units have already reached a nonroot port from the
external edge and child, its parent edge supplies \(U+1-z\) there and
\(1+z\) toward the root. The root then receives its prescribed load.
Copy the old price to every port, retaining the pin on one of its copies.
Every port has at most one external edge, so unequal-slope edges have
disjoint endpoints. The recurrence for edge weights adds original weights
and integers along a path and has polynomial binary length.

\section{Why physical feasibility suffices}
\reviewmark{Price bounds and the overload-to-underload path argument.}

Fix any feasible split \(z^0\) ignoring preferences. For each ordered pair
of endpoints \(v,w\) of edge \(e\), buying at \(v\) implies
\[
 P_w\ge (b_{ev}/b_{ew})P_v+(a_{ew}-a_{ev})/b_{ew}.
\]
Choose an integer \(K\ge2\) bounding every ratio \(b_{ev}/b_{ew}\) and
every absolute intercept \(|a_{ew}-a_{ev}|/b_{ew}\) over ordered
endpoint pairs.
Write \(n=|V|\) and \(B_0=1+K^n(|t|+nK^n+1)\).
A composition along a simple path has slope between \(K^{-n}\) and
\(K^n\) and absolute intercept at most \(nK^n\). It sends \(B_0\) strictly
above \(|t|\) and sends \(t\) strictly above \(-B_0\).

Let each edge choose a full split maximizing its payoff, and let each
unpinned vertex choose a price in \([-B_0,B_0]\) maximizing price times
its excess load at the current split. The price at \(q\) stays \(t\).
On the product of these compact convex strategy sets the best-response
correspondence has nonempty convex values and closed graph. Kakutani's
fixed-point theorem gives simultaneous best responses. An overloaded
unpinned vertex then has price \(B_0\), and an underloaded one has
price \(-B_0\).

If some load differs from its prescription, compare the split to \(z^0\).
On each edge direct its net correction from the endpoint with excess
allocation to the endpoint with a shortfall. Flow decomposition gives a
simple path from an overloaded vertex to an underloaded one. Each edge
buys at its departure vertex, so its price comparison holds along the
path. A path starting off the pin ends at a price strictly above \(|t|\),
which neither possible endpoint price attains. A path starting at the pin
ends above \(-B_0\) at an underloaded unpinned vertex. Both are
contradictions. Thus the fixed point realizes every prescribed load.

\section{Recovering a split from an ordinary market}
\reviewmark{Signed translation, price pin and exact filling in every answer.}

Prices and row maxima for the split can have either sign. Compute an
integer \(B>|t|\) bounding their absolute values strictly in \emph{some}
clearing witness. This needs no knowledge of its support. Write every
possible offer inequality, its reverse for a used edge, and the pin
equation. On any feasible support split free variables into differences
and add slacks. If \(h\ge1\) bounds the number of rows and \(H\ge1\)
bounds all integer entries and right-hand sides after clearing denominators,
a basic solution has coordinates at most \((hH)^h\).
For example \(B=2(hH)^h+2+\lceil|t|\rceil\) suffices. The load constraints
on that support also admit a polynomial-bit rational split.

Make one ordinary buyer of mass \(m_e\) per edge, listing its two
endpoints. Add \(B(1+b_{ev})\) to each intercept. This translates a
witness's prices and utilities by \(B\), making all of them positive.
Increase good \(q\)'s capacity by \(S=1+\sum_e m_e\), and add a buyer of
mass \(c_q+S+1\) whose sole offer is \(t+B-P_q\). The translated witness
extends by sending \(S\) units of this buyer to \(q\). Every market
equilibrium has \(P_q=t+B>0\). A lower price would make the new buyer
overfill \(q\), and a higher price would leave its enlarged capacity unfilled.

Compare an arbitrary market answer to the strictly positive witness
\((P^*,u^*,z^*)\). On \(D=\{v:P_v<P_v^*\}\) every original buyer using
\(D\) in the witness strictly improves there and can now buy only there.
Their original masses already cover \(c(D)=\sum_{v\in D}c_v\), so all
goods of \(D\) fill. Every nonpin good outside \(D\) has positive price
and fills as well. For the rows put \(N=\{e:u_e<u_e^*\}\) and
\(T=\{v:P_v>P_v^*\}\). All witness purchases of \(N\) lie in \(T\).
Every good of \(T\) fills in the returned solution, and only a buyer of
\(N\) can purchase there, since its used offer was strictly better at the
witness prices. The pin is outside both \(D\) and \(T\). Thus
\[
 \sum_{e\in N}m_e\ \le\ c(T)\
 \le\sum_{e\in N}\sum_v z_{ev}\ \le\sum_{e\in N}m_e .
\]
All these buyers fill. The remaining original buyers have positive utility
and fill. Total balance now gives exactly \(c_q\) original units at \(q\).
Removing the auxiliary buyer and subtracting \(B\) from prices gives the
required split and prices. The transforms and decoder have polynomial
binary cost, completing Theorem~\ref{thm:pinned}.

\section{A sparse ordinary market with outside options}
\reviewmark{Copy masses, unused supply and port-price equality.}

\begin{corollary}\label{cor:sparsemarket}
\(\Market\) is many-one equivalent to its rectangular-slope subclass in
which a buyer has at most two listed goods and a good has at most three
listed buyers. After a two-option buyer is drawn as an edge between
goods, the unequal-slope edges have disjoint endpoints.
\end{corollary}
\begin{proof}
Start with Lemma~\ref{lem:rectangle}. For each of buyer \(i\)'s \(k_i\)
options create a buyer of mass \(K_i=m_i+1\) with that option and a
private good quoted at \(A-P\), where \(A=1+\max(0,\max a_{ij})\).
Its capacity \(k_iK_i-m_i\) is too small for all \(k_i\) buyers at
price zero and too large for \(k_i-1\) buyers to fill. It therefore
fills and every buyer uses it. Their common utility bounds all real
offers, with equality on positive trades. Real purchases total at most
\(m_i\), exactly \(m_i\) at positive utility. For the lift, give each copy its real trade \(g_{ij}\). Distribute
\(s_i=m_i-\sum_jg_{ij}\) as unspent amounts \(s_{ij}\in[0,K_i-g_{ij}]\). These upper bounds sum to at least \(s_i\). The remainder
\(K_i-g_{ij}-s_{ij}\) buys the private good at \(A-u_i>0\), filling it
with \(k_iK_i-m_i\). Any unspent mass has \(u_i=0\).

Use the port paths above. Each connecting buyer offers \(H-P\) at both
ends, with \(H>\max(0,\max a/b)\) over base offers. Above \(H\) a
positively priced port could not fill. A child dearer than its parent
would need to fill without its parent buyer. A cheaper child would
attract that buyer's entire mass, exceeding capacity. Thus port prices
agree. At \(H\) only connecting buyers could fill the path, whose
total capacity exceeds their masses by \(c>0\). The common price is
below \(H\), so all connecting buyers fill. Subtract their masses to
recover the original column. In the other direction the leaf-to-root
allocation leaves precisely the old slack at the root. Real ports
inherit their good class, while private goods and connecting buyers
are ordinary. Each port meets only one external edge.
\end{proof}

\chapter{A prescribed flow quota and a price-setting cycle}\label{sec:quota}

A rectangle market has slope two on \(B\times E\) and slope one
elsewhere. In this chapter every row and column fills exactly, the
margins on the two sides have equal totals, and both prices and
utilities may be signed. Write \(W(f)=\sum_{i\in B,j\in E}f_{ij}\).
There are two different pieces of information one can request. The
first is a clearing with a prescribed value of \(W\). The second is a
single cycle of tight edges from which to read a specified source price.

\section{How prices order the feasible quotas}
\reviewmark{Strict price ordering, quota monotonicity and exact bisection.}

Assume there is no proper subset balance: no nonempty proper buyer set
and nonempty proper good set have equal total margins. Every full
transport has connected positive support, since each support component
balances internally. Consider two full clearings with price vectors
\(P,P'\), and take \(S=\{j:P'_j>P_j\}\). For a buyer using a good
of \(S\) at \(P'\), returning to \(P\) strictly improves that offer
and weakly worsens every offer outside \(S\). The buyer therefore
sends its entire mass to \(S\) at \(P\). These buyers supplied all
of \(S\) at \(P'\), and its total load is fixed, so they exhaust
its capacity at \(P\). No edge of the first positive support joins
these buyers or goods to their complements. By connectedness a
nonempty \(S\) is the whole good side. Interchanging the clearings
handles the case \(S=\varnothing\). Thus distinct prices are strictly ordered at
every coordinate, and utilities move strictly in the opposite direction.

For \(P<P'\) put \(d_j=P'_j-P_j\) and \(t_i=u_i-u'_i\). The offer
inequalities imply \(t_i\le b_{ij}d_j\) on edges used by \(f\) and
\(t_i\ge b_{ij}d_j\) on edges used by \(f'\). Taking logarithms and
summing with weights \(f\) and \(f'\) gives
\begin{equation}\label{eq:quota-order}
 W(f)\log 2\ \ge\
 \sum_i m_i\log t_i-\sum_jc_j\log d_j\
 \ge W(f')\log 2 .
\end{equation}
The common middle expression uses only the margins. Strict positivity
of the differences is why connected support matters here.

\begin{lemma}\label{lem:quota-back}
On tables with no proper subset balance, finding a full clearing at a
promised feasible rational quota reduces to polynomially many adaptive
\(\Market\) queries.
\end{lemma}
\begin{proof}
Pin any one good price to a trial \(z\). Represent each row of mass
\(m_i\) by the private vertex and two-endpoint edges of
Chapter~\ref{sec:pinned}, using the signed full margins directly.
Pin the chosen good, keep the private prices free, and use
Theorem~\ref{thm:pinned}. Physical feasibility is independent of \(z\).
The returned prices are the unique ones at that pin by the ordering
argument. At those prices two transportation LPs on tight edges give
the minimum and maximum of \(W\). Every intermediate value is attained
by convexity. If the desired value is above this interval, decrease \(z\).
If below, increase \(z\), by \eqref{eq:quota-order}.

For bit termination, each possible support at the desired quota is
described by a rational polyhedron: retain full margins, impose the
quota, set unused flows to zero, and require equality at the supported
pairs. Its projection onto \(z\) is a closed interval, point or ray.
Finite endpoints and a finite witness in any nonempty projection have
uniform polynomial bit bounds from the possible individual rows, using
basic solutions after splitting free variables. Compute a bound \(D\) on
their denominators and start with a symmetric integer bracket containing
a witness strictly inside. Test its endpoints as well. Feasibility of the
quota is promised, and each missed probe discards only a side containing
no solution. If no probe hits, bisect until the bracket has width less
than \(1/(2D^2)\). Both ends have missed, so a projected endpoint lies
inside. Rational reconstruction gives that unique bounded-denominator
number. Query there and solve the tight-edge transportation LP with the
quota imposed. All probes, including the last, have polynomial bit size.
\end{proof}

\section{Half the exceptional buyer mass in one query}
\reviewmark{Subset balances, quota calibration and removal of the dummy margins.}

\begin{theorem}\label{thm:halfquota}
One clearing with prescribed quota
\(Q=\tfrac12\sum_{i\in B_{\rm new}}m_i\) suffices to solve an ordinary
market. The new margins are positive integers without proper subset
balance, \(Q\) is a half-integer, and all four edges of a designated
two-by-two square carry positive flow in every answer.
\end{theorem}
\begin{proof}
First use Lemma~\ref{lem:rectangle}. Complete missing pairs with
strictly negative original intercepts and the designated slopes. For
\(m\) buyers and common margin denominator \(D\), add \(1/(2mD)\) to
each buyer mass, then multiply all margins by \(4mD\). A nonempty buyer
subset now has residue in \(\{2,4,\ldots,2m\}\) modulo \(4m\), and
all good capacities have residue zero. Write \(p,c\) for the totals,
\(M\) for the exceptional buyer mass, and \(K=1+p+c\).

Add an ordinary dummy row of mass \(c+K\) and an ordinary dummy good
\(D_0\) of initial capacity \(p+K\), with zero intercepts on these new
pairs. Full clearing with \(P_{D_0}=0\) represents an ordinary clearing:
the dummy--dummy flow is at least \(K\), and each real unused margin can
be traded with its dummy at zero utility or price. Every transport is
connected. A component with neither dummy would balance original subsets,
one containing both would balance their complements, and either dummy
alone outweighs the opposite real side.

Add ordinary and exceptional buyers of masses \(8K,2K+1\) and an
exceptional good \(S\) of capacity \(5K\). Enlarge \(D_0\) by \(5K+1\).
The two new buyers have intercept zero at \(S,D_0\). Every other new
incidence has intercept \(-1\). Slopes throughout follow the rectangle.
The quota is \(Q=K+(M+1)/2\). To check subset balance, let \(a,b,s,h\)
indicate the presence of the two new rows, \(S,D_0\). Relative to the old
augmented table the change in the subset difference is
\[
 K(8a+2b-5s-5h)+(b-h).
\]
The absolute old difference is at most \(2K-1\). A nonzero coefficient
of \(K\) has magnitude at least two. In the two cases of magnitude two
the unit correction has the same sign, and magnitudes of three or more
also dominate. If the coefficient vanishes all four indicators coincide,
the correction vanishes, and old nonbalance forces the empty or full
choice.

At the ordinary clearing extended at zero pin, give \(S\) and the
two new buyers zero potentials. All \(-1\) edges are strict. If the old
exceptional flow is \(w\in[0,M]\), let the new exceptional buyer send
\(y\) to \(S\) and \(2K+1-y\) to \(D_0\). The ordinary buyer sends
\(5K-y\) and \(3K+y\) there, respectively. The same potentials support
the entire interval \([w,w+2K+1]\), which contains \(Q\) strictly
inside for every possible \(w\). By \eqref{eq:quota-order} strictly
higher prices have quotas at most \(w\), and strictly lower ones have
quotas at least \(w+2K+1\). All answers at \(Q\) therefore have the
reference potentials. At those potentials the old contribution to the
quota is between zero and \(M\), so all four new square flows are positive.
Delete the two rows and \(S\), undo the dummy augmentation and scale, and
restore the original margins by the fixed-price circulation in
Theorem~\ref{thm:extreme}.
\end{proof}

Thus one quota query suffices in this direction. The reverse algorithm
in Lemma~\ref{lem:quota-back} can use several queries. At a vertex of the
tight transportation face intersected with the quota, the support has
at most as many edges as vertices, by the rank of the margin equations
plus one quota row. The four positive square edges already give a cycle,
and the support is connected. It follows that this square is its only
cycle. The next construction encodes an unknown price on the cycle itself.

\section{Reading a designated good price from a cycle}
\reviewmark{Tree perturbations, the reflected copy, cycle gains and decoding.}

\begin{theorem}\label{thm:distcycle}
Consider signed rectangle markets with a feasible rational quota, no
proper subset balance, and the promise that the tight graph at every
quota clearing has a unique simple cycle. A polynomial batch of cycle
queries, one per source good, suffices for \(\Market\). Conversely a
\(\Market\) oracle finds the promised cycle with polynomially many
adaptive calls.
\end{theorem}
\begin{proof}
Use the preceding perturbed ordinary margins and the full dummy
augmentation at \(P_D=0\), before adding the calibration square. Call
the number of vertices \(v\). Perturb the intercepts by distinct
positive powers \(\theta^e\) so that the tight graph at that pin is a
tree \(T\), and its edges still support the original intercepts at pin
zero. Here is a uniform choice argument. A vertex of the tight transport
is a positive spanning tree by nonbalance. Propagate its equalities
from the pin. Each extra edge's slack is an affine expression in the
perturbed intercepts, with a nonzero coefficient on its own perturbation
power. All tree path coefficients and denominators use at most \(v\)
factors of two. Clearing the original intercept denominators therefore
bounds the integer polynomial coefficients by some \(H\) of polynomial
bit length, computable without the tree. Taking \(0<\theta<1/(2+2vH)\)
makes the first nonzero term dominate and preserves the sign of every
originally nonzero slack. Each candidate slack uses at most \(v+1\)
intercepts. Hence no extra edge ties, and the generic tree is feasible
with original intercepts. A fixed-price solution for the generic table
exists by the pinned construction of Lemma~\ref{lem:quota-back}.

Mark the ordinary anchor \(D\) and the requested source good \(j\).
If \(r_i,c_k\in\{0,1\}\) are the exceptional row and good labels, set
\(h_D=0,\ h_i-h_k=r_ic_k\) on the tree. Its prices at a nearby anchor
\(z\) are \(P_k(0)+2^{h_k}z\), and its row utilities change by
\(-2^{h_i}z\). All \(|h|\le v\). An interval \(|z|<\rho\) on which
every other edge stays strict is computable without \(T\). For instance,
with a common denominator \(D'\) of the generic intercepts use
\(\rho=1/(D'2^{3v+6})\), by propagating each slack and its derivative
along the tree. Its positive tree flows are at least a common margin
quantum \(g>0\), since deleting a tree edge expresses its flow as a
margin difference.

Make a transposed copy. Buyer \(k'\) has group \(c_k\), good \(i'\)
has class \(1-r_i\), and their intercept is
\(2^{c_k}a_{ik}/b_{ik}\). Transposing a source clearing gives prices
\(\bar P_{i'}=u_i\), utilities \(\bar u_{k'}=2^{c_k}P_k\), and flow
\(f^{\mathsf T}\), as is seen by multiplying its offer inequality by
\(2^{c_k}/b_{ik}\). The sum of the exceptional flows of the two copies
is the known capacity of the source's exceptional goods. Connect mirror
buyer \(D'\) to source good \(D\) by an ordinary edge of intercept zero,
adding one to each of their margins so the edge can carry one unit.
The equality on this bridge compares opposite anchor values \(z,-z\).

Make a path from \(D\) to \(j\). One stage uses an ordinary row between
the current ordinary good and a fresh exceptional good, then an exceptional
row from that good to a fresh ordinary one. All intercepts on the path
are zero. If both options of a row are used, the stage doubles the price.
After \(K\) stages use an ordinary row to copy the final price to \(j\).
Choose \(K>v+1\) so that \(U/(2^K-2^v)<\rho\) for an upper bound \(U\)
on \(|P_j(0)|+1\), obtained by tree propagation. This requires only
polynomially many stages. Each path row and internal good has margin
\(\kappa\) with \(0<\kappa<g/4\). Enlarge \(D\)'s capacity by another
\(\kappa\). When every path row sends \(t\) forward and \(\kappa-t\)
backward, the source tree must put \(t\) more into \(D\) and \(t\) less
into \(j\). Alternating changes of \(t\) on its \(D,j\) path do this.

On the tree, \(W(f)=\sum_i m_ih_i-\sum_k C_kh_k\), where \(C_k\)
are its capacities. Thus \(W\) changes by \(h_jt\), and the new path
contributes \(K(\kappa-t)\). Take the known quota \(q_0-\kappa/2\),
where \(q_0=\sum_{k:c_k=1} C_k+K\kappa\).
It is attained by \(t=\kappa/[2(K-h_j)]\), strictly between zero and
\(\kappa\). The source path requires \(P_j=P_j(0)+2^{h_j}z\) and the
new one requires \(P_j=2^Kz\). Their common choice
\(z=P_j(0)/(2^K-2^{h_j})\) has \(|z|<\rho\), also for the mirror
at \(-z\). This exhibits a connected positive tight graph with exactly
one cycle and unequal gains \(2^{h_j}\) and \(2^K\) on its two paths.
Unlisted pairs can be assigned an intercept below every used comparison
by the same propagation bounds.

To enforce nonbalance in the joined query, perturb its independent
margins by distinct binary digits \(2^a\gamma\) and put the balancing
change on the last good. Choose \(\gamma\) dyadic after a common integer
scaling of margins and quota. For any proper subset difference use its
complement if necessary to avoid the balancing good. Old nonzero
differences survive, and an old zero difference gets a nonempty signed
sum of distinct digits. These perturbations can preserve strict
positivity of every constructed flow at the same quota and potentials.
Indeed the margin matrix on the unicyclic graph, with one redundant
row deleted and the quota row added, is square and nonsingular. Its
only possible margin-preserving change is a circulation on the cycle,
whose quota coefficient has magnitude \(K-h_j>0\). In dimension \(s\),
after the common scaling the right sides are integral, and both the
determinant and all cofactors have magnitude at most \(s!\). Every
positive coordinate is at least \(1/s!\), and changing the right sides
by total absolute amount less than \(1/(2(s!)^2)\) keeps it positive.
Use the known number of vertices for \(s\). All choices require only
polynomial bit length.

For any other quota clearing, different prices would be strictly ordered.
Equality of quotas in \eqref{eq:quota-order} would force
\(t_i=b_{ik}d_k\) on every constructed positive edge. The two unequal
cycle gains make that impossible. The prices and hence the tight graph
are therefore fixed. Delete the new path from the returned cycle and
propagate the \emph{original} tight equalities from price zero at \(D\)
along the remaining source path. This gives the original canonical price
of \(j\). Use the same source perturbation for every coordinate so all
queries are independent. At those prices recover original trades by
Theorem~\ref{thm:extreme}. For the converse, find a quota clearing by
Lemma~\ref{lem:quota-back}, test offer equalities, and read the cycle
from its tight graph.
\end{proof}

\chapter{Ordinal matching and ordered Boolean search}\label{sec:boolean}

The market also has a discrete readout. In an ordinal matching instance,
a row asks for the identity of one partner and ranks its neighbors
strictly. In an ordered Boolean instance, the answer may be a fixed
point or a pair of points witnessing a crossing. We give direct
reductions for specified families of these problems.

\section{Stable partners from fractional trades}
\reviewmark{Endpoint rankings, support rounding and the converse lift.}

Let a simple bipartite graph list acceptable pairs \(e=ij\), strictly
ordered at each endpoint. Each acceptable partner is preferred to
remaining unmatched. Write \(r_i(e)\) and \(r_j(e)\) for the numbers of
neighbors ranked below \(e\) by its two endpoints. Give every buyer and
good unit margin. With zero reservations set
\[
 a_e=4^{r_i(e)},\qquad T_e=4^{r_j(e)},\qquad b_e=a_e/T_e .
\]
These numbers need only linearly many bits in the largest degree.

\begin{theorem}\label{thm:ordinal}
At every equilibrium, selecting each active buyer's favorite positive
trade gives a stable integral matching. Every stable matching is obtained
in this way at some equilibrium of this same market.
\end{theorem}
\begin{proof}
If the buyer at \(e\) leaves mass unspent or uses a worse partner, its
utility \(u_i\) is at most \(a_e/4\), by the reservation or a used offer.
If the good is slack or receives positive mass from a worse partner,
its price \(P_j\) is at most \(T_e/4\). Both failures together contradict
\(u_i+b_eP_j\ge a_e\). At least one endpoint therefore fills entirely
with partners at least as good as \(e\).

Let \(e\) be a buyer's best positive edge. If \(f_e=1\), its good
has \(e\) as its sole partner. Otherwise the buyer is slack or uses a
worse partner, so the good is full and \(e\) is its worst positive
edge. Thus the choices inject nonempty buyers into full goods. Reversing
the roles injects nonempty goods into full buyers. All full vertices
are nonempty, so these injections give bijections on the active vertices.
The choices form a matching. An omitted edge with positive flow has a
better chosen competitor at its buyer. For an omitted zero-flow edge
one endpoint fills on strictly better edges, so that endpoint has a
better chosen partner. This proves stability. This rounding step is the
one in Roth, Rothblum and Vande Vate~\cite[Lemma~12]{rrvv}.

Given a stable matching, assign flow one to each matched edge, and
set \(u_i=a_e/2,\ P_j=T_e/2\) there. Unmatched vertices have zero
potential. Every selected edge is tight. At each omitted edge a better
partner at an endpoint gives either \(u_i\ge2a_e\) or \(P_j\ge2T_e\).
This verifies all offer and filling conditions. The decoder returns the
original matching.
\end{proof}

The same market realizes Boolean equations on the edges. Let \(x_e\)
be one precisely for selected pairs and let \(F_e(x)\) be the NOR of
the bits for pairs strictly preferred to \(e\) by either endpoint.
The empty NOR is one. Two selected pairs cannot meet, and each
unselected pair must have a better selected competitor. These are
exactly the fixed-point equations \(x=F(x)\). For example, on buyers
and goods numbered modulo \(n\ge3\), use the three offsets \(0,1,2\)
at each buyer, in decreasing preference order, and the reverse order
at each good. Each bit then reads two competitors and is read by two
others. The two preferences at a vertex, the tight inequalities, and
the integral readout all use the same physical trade.

\section{A fixed point or a crossing in the given coordinate order}
\reviewmark{Pivot invariant, slice feasibility and witness recovery.}

Let a Boolean map \(F\) have coordinates \(1,\ldots,n\). An ordered
crossing at \(i\) consists of two points with the same suffix above \(i\)
and opposite \(i\)-bits, each stationary on all rows below \(i\), with
the \(i\)-row pointing up at zero and down at one. Such a pair is a
violation witness in the ordered contraction problem~\cite{fgms}.

\begin{theorem}\label{thm:bool}
For \(F(x)=Cx+b\) over \(\mathbb F_2\), a market with \(O(n^3)\)
goods, buyers and listed offers returns a full fixed point or an ordered
crossing. For explicit Boolean row tables with a supplied width-\(w\)
tree decomposition containing each row and all its arguments in a bag,
the corresponding market has size \(2^{O(w)}\operatorname{poly}(L)\).
In particular the latter is a polynomial reduction at fixed width.
\end{theorem}
\begin{proof}
For affine rows, solve \((I+C)x=b\) by eliminating rows in the supplied
order. Reduce each row by previous pivots in their creation order and
choose its leftmost nonzero coefficient as its pivot. Mark the first
row \(i\) that reduces to \(0=1\) or has pivot to the right of \(i\).
Continue eliminating to the end. If consistent, back substitution
returns a fixed point. If inconsistent, an inconsistent row is eventually encountered and a
first mark therefore exists. Earlier rows with residual \(0=0\) impose
no condition. Every earlier pivot is strictly below \(i\). In reverse
creation order, assign those pivot bits to satisfy their rows, taking
nonpivot lower bits arbitrarily. A later residual has zero coefficients
at all earlier pivots, so any assignment at and above \(i\) extends to a
stationary lower prefix. The marked residual involves only bits strictly
above \(i\). Choose their common suffix to violate it
(if it is \(0=1\), any suffix does). Set \(x_i=0\) and \(x_i=1\) in
turn and back-substitute the earlier rows. The two points both violate
the original \(i\)-row since the previously eliminated equations hold.
This gives exactly the crossing in the stated order.

Store one-hot pivot indices and implement conditional row additions
by XOR and masking. There are \(O(n^2)\) row-pair updates each using
\(O(n)\) Boolean gates, and at most three back substitutions. The
minimum and maximum price gates in Appendix~\ref{sec:compile} implement
each Boolean operation with constantly many goods and buyers using
paired bits and their complements. The reverse sizing of their capacities
gives polynomial encoding length even though the numbers themselves can
be large.

For the tree decomposition, consider a slice with a fixed suffix and
require stationarity only on free rows. Assign each row test to a bag
containing its scope. For each assignment to a bag's separator with
its parent, record whether it extends to a bag assignment meeting its
local tests and all child messages. Computing these messages from the
leaves gives slice feasibility. A downward pass selecting the first
compatible bag assignment yields a witness, because the bags containing
any one variable form a connected subtree.

If the full slice is infeasible, inspect its two children obtained by
fixing the greatest free coordinate \(i\) and dropping its row.
Descend into an infeasible child if one exists. An empty prefix is
always feasible, so eventually both children of an infeasible parent
are feasible. Their two witnesses satisfy the earlier rows and share
their higher suffix. Neither can satisfy row \(i\), or the parent would
be feasible. They give the required crossing. This needs at most
\(2n+1\) slice tests, each of size \(2^{O(w)}\operatorname{poly}(L)\).
Unroll the tests and the choices in the same Boolean price gates.
Every market equilibrium supplies the chosen witnesses. In both
constructions the source description restricts the input map.
\end{proof}

\section{Kernels in directed graphs with restricted feedback}
\reviewmark{Coloring, monotone slices and the two exceptional cases.}

For a directed graph let \(F_i(x)\) be the NOR of the bits at the
outneighbors of \(i\). At a fixed point, the vertices with bit one
form a kernel: no arc joins two selected vertices and each unselected
vertex has an arc to a selected one. Here the direction of an arc is
from a row to a bit it reads. Two further restrictions give direct
reductions to \(\Market\).

\begin{corollary}\label{cor:nor}
An ordinary market returns a kernel or an ordered crossing in each of
the following cases.
\begin{enumerate}
\item A vertex \(h\) is supplied such that deleting the arcs whose head
is \(h\) makes the underlying graph bipartite. The coordinate order
is arbitrary.
\item Deleting the last two vertices in the supplied order leaves a
bipartite underlying graph. The deleted vertices may have arbitrary
dependencies.
\end{enumerate}
\end{corollary}
\begin{proof}
For the first case, two-color the remaining graph by \(c_i\) and use
variables \(y_i=x_i\oplus c_i\). A color-zero row becomes AND and a
color-one row becomes OR on the signed neighbor variables. An
opposite-color read contributes \(y_j\), and a same-color read
contributes \(1-y_j\). Every same-color read must have \(j=h\).
Suffix bits fixed in a slice supply constants.

If \(h\) is free, guess \(t=y_h\) and replace only the same-color reads
by \(1-t\). The resulting Boolean map is monotone, including its row
for \(h\). For \(t=0\) iterate from all zeros to the least fixed point.
For \(t=1\) iterate from all ones to the greatest. Each takes at most
the number of free bits in strict updates. A passing extremum has its
own \(h\)-bit equal to \(t\), so it solves the original slice. Conversely,
any solution with \(h\)-bit zero bounds the least fixed point above,
and any with \(h\)-bit one bounds the greatest fixed point below.
Thus at least one extremum passes precisely when the slice is feasible.
For pinned \(h\) use its supplied value. The descent of
Theorem~\ref{thm:bool} now finds a full root or crossing.

In the second case pin the last two coordinates \(u,v\) to each of
their four possible pairs of bits. The other rows are monotone after
the coloring, so obtain one stationary interior witness for every pair.
For fixed \(v\), evaluate row \(u\) on the two witnesses. If neither is
stationary at \(u\) they give a crossing there. Otherwise retain one
stationary witness for that \(v\). With one retained witness at each
value of \(v\), either one is a full root or the two give a crossing at
\(v\). All these choices and tests have polynomial-size Boolean circuits
and hence single-market implementations as above.
\end{proof}

\section{One additional input per prefix row}
\reviewmark{Binary clauses, suffix constraints and the last-coordinate crossing.}
\begin{corollary}\label{cor:binary-prefix}
Order the bits as \(x_1,\ldots,x_m,b\). Each prefix update may read
\(x_i\), one specified prefix bit \(x_{p(i)}\), and \(b\). The last
update can be an arbitrary Boolean circuit on all the bits. One
polynomial-size market returns a fixed point or an ordered crossing.
This includes every update map on three ordered bits.
\end{corollary}
\begin{proof}
Fix \(b=t\). Evaluate the circuit for each prefix row on its at most four
local input pairs. For each forbidden pair \((r,s)\) for which row \(i\)
is not stationary, impose
\((x_i\ne r)\lor(x_{p(i)}\ne s)\). For a single distinct input use a
unit clause for each forbidden value. In any slice with \(b\) fixed,
keep exactly the stationarity clauses for its free prefix and impose
its suffix by unit clauses. This 2-CNF has a satisfying assignment
precisely when the slice is feasible. The implication-graph algorithm
for 2-SAT finds such an assignment or decides infeasibility in polynomial
time~\cite{apt}.

Test the two slices \(b=0,1\). If either is infeasible, descend inside
it as in Theorem~\ref{thm:bool} until both children of an infeasible
slice supply a crossing. Otherwise take one prefix witness from each
slice and evaluate the last row on both. If either is stationary there,
it is a full root. If both fail, the two witnesses cross at \(b\).
Fixing tie choices gives a polynomial-time Boolean computation, which
the price gates of Appendix~\ref{sec:compile} encode in one market.
\end{proof}

\section{Recurrent signed comparators}
\reviewmark{Damping error, support repair and odd-cycle readout.}

Suppose each Boolean wire is produced once and consumed in at most one
input position. A comparator takes two signed wires and returns their
AND and OR. Outputs may be complemented or discarded, and signed copies
and constants are allowed. Directed cycles are permitted. Unlike the
affine construction, the following market routes trades along the given
cyclic wiring.

\begin{theorem}\label{thm:comparator}
For these signed comparator networks, a polynomial-size ordinary market
returns a Boolean fixed point or an ordered crossing.
\end{theorem}
\begin{proof}
Extend a comparator to real inputs \(a,b\in[0,1]\) using
\(L(s)=\max(0,s-1), H(s)=\min(1,s)\) with \(s=a+b\). Its outputs
are \(L(s),H(s)\). They are monotone in \(s\) and their sum equals
\(s\). Write \(G\) for the resulting real update. It is nonexpansive
in \(\ell_1\) because each input is used at most once. Use two market
deliveries for each wire, intended to carry \(x_i,1-x_i\). Swap them
for negation. Run \(L,H\) on the two positive inputs and \(H,L\) on
the two negative inputs to produce respectively the AND and OR pairs.
The doubled update \(\Phi\) is also nonexpansive and preserves balanced
pairs because \(H(2-s)=1-L(s)\).

The ordered-tranche good of Appendix~\ref{sec:cayley} with two unit
tranches sends its first unit to one successor and the rest to the
other. Its outputs give \(H(s),L(s)\) for every \(s\in[0,2]\).
Signed copies exchange or retain the two outputs, and a dedicated buyer
delivers a constant. Separate one-unit tranche goods cap each wire
delivery at one, including on cycles. The capacity and price guards
in the appendix work against the summed possible incoming buyer masses
without assuming any order of evaluation.

On each output implement \(D(s)=(1-\alpha)s+\alpha/2\) approximately.
Choose a dyadic \(\alpha=2^{-a}\) and put \(a\) binary splitters in
series. At each stage collect one half and forward the other, then
discard the final remainder. Use the scaled alternating-switch market
of Appendix~\ref{sec:cayley} with dyadic unit \(\epsilon\). Its two
outputs sum to its input and differ from exact halves by at most
\(\epsilon/2\). The forwarded mass differs from \(\alpha s\) by
less than \(\epsilon\), by the geometric error sum. Add the constant
\(\alpha/2\) to the collected outputs and cap at one. Capping cannot
increase error against \(D(s)\in[0,1]\).

Thus the paired deliveries \(z\) in any equilibrium satisfy
\(\|z-D(\Phi(z))\|_1\le2n\epsilon\). The exact map \(D\circ\Phi\) has
contraction factor \(1-\alpha\). Its unique fixed point is balanced
by iteration from balanced pairs, and has the form \((x^*,1-x^*)\).
The distance of \(z\) from this point is at most \(2n\epsilon/\alpha\).
With \(\hat x_i=(z_i^++1-z_i^-)/2\), nonexpansiveness gives
\[
 \|\hat x-G(\hat x)\|_\infty\le4n\epsilon/\alpha+\alpha/2 .
\]
For exact repair record the closed branch \(s\le1\) or \(s\ge1\)
of every comparator at \(\hat x\). On that cell the equations
\(x=G(x)\) are affine with integer coefficients of bounded size.
Impose the cell and box inequalities, relax the fixed-point
equalities by a common slack and minimize it. The determinant
argument of Lemma~\ref{lem:repair} gives an input-computable \(Q\)
of polynomial bit length for \emph{all} such LPs. Take
\(\alpha<1/(4Q)\) and \(\epsilon<\alpha/(16nQ)\). The displayed
bound is less than \(1/Q\), so the LP returns an exact fixed point
\(y\) in that cell. Both \(a\) and the switch depth
\(O(\log(1/\epsilon))\) are polynomial in the input length.

Every fractional wire of \(y\) must be produced. A comparator with
one or more fractional inputs produces at most one fractional output,
and an integral-input gate produces none. Since no wire is read twice,
counting forces every fractional wire to be read exactly once and
prevents two of them from meeting at a comparator. Fix all integral
coordinates. Each remaining gate now passes its single fractional
input, possibly complemented, to one retained output. For example, a
second input zero forces AND to zero and passes the first input to OR.
This remains true at the ends of the closed interval. The fractional
wires therefore form signed permutation cycles. With even sign parity
on every cycle, propagate bits to obtain a Boolean fixed point.
Otherwise select the odd cycle with smallest maximum index \(j\). Pin
all higher free coordinates identically in the two witnesses and choose
\(x_j=0,1\), respectively. Every other odd cycle meets a pinned higher
coordinate. Propagate into the unpinned parts of the cycles, choosing
either consistent assignment on an untouched even cycle. All rows
below \(j\) are stationary and the chosen odd cycle forces
\(F_j=1-x_j\) at both witnesses, giving the crossing.
\end{proof}

\chapter{From a single switching walk to equilibrium}\label{sec:arrival}

Does a chip starting at the origin reach the destination if each switch
alternates between its two successors? Make the destination the YES
terminal and redirect vertices with no path to it to NO. An infinite
walk would visit some remaining switch, both its successors, and hence
a terminal infinitely often. The resulting walk therefore stops.

\begin{theorem}[Lin, {\cite[Theorems~4.1 and~4.11]{book}}]\label{thm:arrival}
One ordinary market equilibrium, or one BNE of a one-resource game with
unit defender gains, determines the ARRIVAL answer. For \(n\) switches
the direct game construction has at most \(12n^2+6n+6\) types and
\(12n^2+6n+7\) targets.
\end{theorem}

The alternative construction in \cite[Appendix~C, Theorem~C.1]{book}
uses at most \(12n^2+6n+2\) types and \(12n^2+6n+3\) targets. Both bounds
refer to explicitly listed games; the construction below is the one that
also supplies the ordinary market.

\section{The terminal inflow is the answer}
\reviewmark{Real-valued flow comparison, terminal conservation and the switching-flow search decoder.}

Interpolate the two exit counts. For \(x=2k+s\), \(k\ge0\) integral
and \(0\le s<2\), put
\[
 h_0(x)=k+\min(s,1),\qquad h_1(x)=k+\max(s-1,0).
\]
Put one unit at the start \(o\). If it is already a terminal, return
its answer directly. Otherwise the fixed-point equations on switches are
\[
 x_v=\mathbf{1}_{\{v=o\}}+
      \sum_{(u,i):\,s_i(u)=v}h_i(x_u).
\]
Here \(s_0,s_1\) are the ordered successors. The sum counts both ports
when the successors coincide. Run counts solve these equations. The
continuous update and the fractional-cycle observation are due to
Haslebacher~\cite[Lemma~14]{has}. The switching-flow certificate of
Dohrau et al.~\cite[Lemma~1]{arrival} gives the single-chip decoder in
\cite[Chapter~4]{book}. The following comparison proves the terminal
statement directly for real flows and arbitrary integral supplies.

\begin{lemma}[Terminal counts of every real solution]\label{lem:real-terminal}
Let \(b_v\) be the initial number of chips at vertex \(v\), including
terminals, and assume every switch has a directed path to a terminal.
Define the update \(F\) on nonterminal coordinates by
\[
 F_v(z)=b_v+\sum_{(u,i):\,s_i(u)=v}h_i(z_u).
\]
For every nonnegative real fixed point \(x=F(x)\), the count at each
terminal \(t\) in any completed legal routing equals
\[
                       b_t+\sum_{(u,i):\,s_i(u)=t}h_i(x_u).
\]
\end{lemma}
\begin{proof}
Start with dispatch counts \(z^0=0\). At each parallel round dispatch all
chips present at the switches, in rotor order at each switch, postponing
newly received chips until the next round. The cumulative counts satisfy
\(z^{k+1}=F(z^k)\): the total number available for dispatch through the
next round equals the initial supply plus everything received through the
current round. Both \(h_i\) are nondecreasing, so \(0\le z^k\le x\) by
induction. These integer counts increase and are bounded coordinatewise,
so they stabilize at \(z^*\le x\), with every chip at a terminal. Terminal
receipts computed from \(z^*\) are therefore at most those computed from
\(x\), coordinate by coordinate. For either fixed point their sum is
\(\sum_v b_v\): sum the switch equations, cancel the internal transfers,
and use \(h_0(w)+h_1(w)=w\). Hence equality holds at every terminal.

For an arbitrary legal order, the dispatch counts of any prefix are
bounded by \(z^*\). If the next dispatch at \(v\) would first exceed
\(z^*_v\), then \(v\) has already dispatched \(z^*_v\) and has received,
including its initial supply, at most \(F_v(z^*)=z^*_v\). It holds no chip.
Thus every maximal legal order terminates. Its final dispatch vector is
itself a fixed point and dominates \(z^*\) by the parallel-round induction,
so the two agree. Existence of a completed routing also follows directly
by running one chip at a time: an infinite walk would have a nonempty
set of switches visited infinitely often, closed under both successors
and disjoint from the terminals. The path hypothesis rules this out.
\end{proof}

\paragraph{Recovering a switching-flow certificate.}
For \textsc{S-Arrival} in the sense of G\"artner et
al.~\cite[Definitions~3 and~13]{ghhkms}, build their graph with
absorbing vertices \(d,\bar d\) and both exits of each dead vertex
redirected to \(\bar d\). For \(o=d\) return the zero flow. Otherwise
use \(d,\bar d\) as terminals in the construction below and read its
fixed point \(x\). Give labeled exit \((v,i)\) the integer
\(h_i(\lfloor x_v\rfloor)=\lfloor h_i(x_v)\rfloor\), and give both
loops at each terminal zero. To check conservation, each fractional
switch sends exactly one fractional exit and must receive at least
one fractional entrance because its initial supply is integral.
Counting shows it receives exactly one. No fractional exit reaches
an integral switch or a terminal. Hence flooring commutes with the
sum of incoming flows at every switch and preserves every terminal
inflow. The two exit counts are the ceiling and floor of half the
rounded dispatch count, so they satisfy switching parity even on
parallel exits. For an encoding with one edge per ordered vertex
pair, add the counts of coincident exits; their sum preserves
conservation and the repeated-edge parity inequality is automatic.
This decodes a polynomial-size switching flow to the
reached terminal from one market or one game equilibrium. Composing
with the market-to-P-LCP map gives the same search reduction to P-LCP.
Run-profile recovery in Definition~14 is a different output requirement.

\section{Implementing the switch with prices and purchases}
\reviewmark{Residual assignment LP, strict exit ordering, coincident ports, saturation and bit sizes.}

Here are all the offers needed for the market in Theorem~\ref{thm:arrival}.
We use the two-copy topology of~\cite[Chapter~4]{book} with the
arc fees of Lemma~\ref{lem:slopes}. The original book uses a different
fee sequence with the same minimum gap. The interleaving of two
offset copies follows Disser and Skutella~\cite[\S3.2]{ds}.
We allow a depth \(m\ge1\) and put \(K=2^{m+1}\). Use two copies of
\(G_m\) from Figure~\ref{fig:count} with \(\delta=1\), and multiply
all fees by six. Their \(2^m\) unit marginal costs are now
\(0,6,\ldots,6(2^m-1)\). Identify the sources \(s\); denote the two
sinks by \(t_0,t_1\). An arc \(e\) has fee \(\kappa_e\) and capacity
\(c_e\). Set
\[
 \kappa'_e=\kappa_e+\mathbf1_{\{\operatorname{head}(e)\in\{t_0,t_1\}\}},
 \qquad M=6+\sum_e\kappa'_e,\qquad M_p=M+1 .
\]
Each arc is a private good of capacity \(c_e\) offered to its two endpoints.
The switch has one input good \(p\) of capacity \(c>K\). Its port buyers
purchase on the two successor goods. Write \(P_0,P_1\) for their prices.

\begin{center}
\small
\begin{tabular}{@{}>{\raggedright\arraybackslash}p{.32\linewidth}
  >{\raggedright\arraybackslash}p{.63\linewidth}@{}}
\toprule
Buyer and mass & Available offers\\ \midrule
Source \(s\), mass \(c\) & \(M_p(1-P_p)\) on the input and
 \(M-\kappa_e-P_e\) on each outgoing arc.\\
Internal vertex, total incoming capacity & \(M-P_e\) on each incoming
 arc and \(M-\kappa_e-P_e\) on each outgoing arc.\\
Sink \(t_i\), mass \(K/2\) & \(M+1-P_e\) on each incoming arc and
 \(g_i=M-3i-P_i\) on its successor good.\\ \bottomrule
\end{tabular}
\end{center}

\begin{lemma}[Local exit law]\label{lem:exit-offers}
Fix \(0\le x<c\) units of other purchases on \(p\) and port prices
\(P_0,P_1\in[0,1]\). In every clearing of the private goods and the
residual input, port \(i\) receives \(h_i(\min(x,K))\).
\end{lemma}
\begin{proof}
Freeze the port offers at \(g_i\), giving each an unrestricted outlet.
The source faces residual input capacity \(c-x\). On that good use the
dual price \(M_pP_p\). With this rescaling, every local clearing is an
optimal primal--dual pair of the assignment LP maximizing intercept
times purchase subject to buyer masses and column capacities. The prices,
buyer values and used offers are exactly its complementary pairs.

Every private arc fills. If an arc is slack, it is priced zero and its
head buys less than its capacity. At an internal head this gives the
maximum possible utility \(M>0\). Incoming purchases alone cannot fill
the row, so the head buys a positive amount on a zero-fee, zero-price
outgoing arc. Its next head again buys less than that arc's capacity.
The network is acyclic. At a sink a zero-price incoming arc gives utility
\(M+1>g_i\), so its entire mass must go to incoming arcs. Their total
capacity is exactly its mass, contradicting the strict deficit.

Write \(f_e\) for the tail's purchase; the head then buys \(c_e-f_e\).
Let \(y\) be the source's input purchase and \(z_i\) sink \(i\)'s port
purchase. Up to a constant the objective is
\[
                      M_py-\sum_e\kappa'_ef_e+\sum_{i=0}^1 g_i z_i .
\]
The constraints are \(0\le f_e\le c_e\), \(0\le y\le c-x\),
\(y+\operatorname{out}(s)\le c\), outflow at most inflow internally,
and \(0\le z_i\le\operatorname{in}(t_i)\). Conversely these quantities
give a feasible allocation with full private goods. As
\(g_i\ge M-4>\sum_e\kappa'_e\), each \(z_i=\operatorname{in}(t_i)\)
at an optimum. Also \(y=c-x\). Otherwise increase \(y\) using spare
source mass, or remove a small flow along a positive path from the source
to an internal vertex with excess inflow or to a sink. This frees source
mass for \(y\), at a sacrifice per unit at most \(\max_i g_i\le M<M_p\).
Thus \(\operatorname{out}(s)\le\min(x,K)\).

If an internal vertex has excess inflow, follow positive arcs backward
to reach \(s\) in its residual network. If no sink were reachable, the
residually reachable set would separate \(s\) from both sinks. Every
leaving arc would be full and every entering arc empty. Its cut capacity
would equal the source outflow minus the sum of its internal excesses,
strictly less than \(K\). Yet the two copies together admit a flow of value
\(K\). A sink is therefore reachable. Push a small amount along a simple
residual path from the excess vertex to that sink. The increase in port
reward per unit exceeds the sum of all possible forward arc fees
\(\sum_e\kappa'_e\). This contradicts optimality. With all internal
flows now conserved, the same cut argument from \(s\) yields a profitable
increase whenever \(\operatorname{out}(s)<\min(x,K)\), since the source
then has spare mass.

At fixed value \(F=\min(x,K)\), write \(A+B=F\) for the two sink inflows.
The objective now minimizes the sum of the two flow costs plus
\((3+P_1-P_0)B\); the sink-entry fees contribute the fixed amount \(F\).
The offset lies in \([2,4]\), strictly between zero and the spacing six.
The two lists of unit marginal costs therefore interleave in strict order,
starting with copy zero. Any use of a dearer interval while a cheaper
one remains available admits a cost-reducing transfer. The last occupied
interval may be fractional. This gives \(A=h_0(F)\) and \(B=h_1(F)\),
including integer endpoints and \(F=K\).
\end{proof}

\begin{lemma}[No mass creation at arbitrary port prices]\label{lem:no-creation}
For the two-copy market above, the sum of the two port purchases is at
most \(x\) for all nonnegative port prices and \(0\le x<c\).
\end{lemma}
\begin{proof}
The private-goods fullness argument in Lemma~\ref{lem:exit-offers} used
only \(g_i\le M<M+1\). Thus the same residual LP applies when a port
reward is negative. Its optimum still has \(y=c-x\). To see this, if the
source has no spare mass and \(y<c-x\), remove a small positive flow
along the path used in that proof and increase \(y\) by the same amount.
At its endpoint reduce the port purchase only if needed to keep it below
the new inflow. This loses at most \(\max(0,g_i)\le M\) per unit, while
the entry gains \(M_p>M\). The spare-mass case just increases \(y\).
The source outflow is consequently at most \(x\). Internal outflow is at
most inflow and port purchases are at most sink inflows. Summing cancels
the internal arcs and gives the claim.
\end{proof}

\paragraph{Composition and binary size.}
Supply \(b_v\) chips at each switch or terminal, with total \(S\ge1\).
Use one copy of the switch market at each of the \(n\ge1\) nonterminals.
Port \(i\) is the input good or terminal good at \(s_i(v)\).
A positive supply \(b_v\) gets a dispatch buyer of mass \(b_v\) with sole
offer \(2-P_v\). Count incoming arcs with their labels and take
\[
 c_v=\max(2K,\deg^-(v)K/2+b_v+1),\qquad
 C_t=\deg^-(t)K/2+b_t+1 .
\]
Terminal good \(t\) has capacity \(C_t\) and its own buyer of mass
\(C_t+1\) with sole offer \(1/2-P_t\).
At a switch every buyer except its source can supply in total at most
\(\deg^-(v)K/2+b_v<c_v\) on its input good. At positive price its source
must therefore buy some, giving \(P_v\le1\) by the outside option.
At a terminal a price below \(1/2\) would attract more than capacity
from its own buyer, and a price above \(1/2\) would leave too little
demand to fill it. Hence all dispatch buyers spend their entire mass.
When a successor is the switch itself or both successors coincide,
each port still belongs to a distinct sink buyer. Count its purchase
as incoming to the designated good and freeze its payoff in the
local LP. Only the source's own input purchase is residual. This
justifies applying Lemma~\ref{lem:exit-offers} also in these cases.

One may set \(C_t=S+1\) at every terminal instead. First bound the switch
input prices by the source argument above, which uses only the masses of
the predecessor and dispatch buyers. Even before bounding terminal
prices, Lemma~\ref{lem:no-creation} and cancellation at switches bound
total purchases at terminal goods from ports and dispatch buyers by
\(S\). Here every dispatch contributes at most its mass, including those
at terminals. A terminal of capacity \(S+1\) cannot fill above price
\(1/2\) and its own buyer exceeds capacity below \(1/2\). Thus its price
is \(1/2\) and every dispatch spends in full, as before.

Let \(x_v\) be nonsource purchases at an input and \(Y_t\) purchases
at a terminal excluding its dedicated buyer. The sum of the outputs at
switch \(v\) is \(\min(x_v,K)\). Counting dispatches and canceling all
internal port purchases gives
\[
                    \sum_tY_t=S+\sum_v(\min(x_v,K)-x_v)\le S .
\]
Each port outputs at least \(\phi(x)=\min((x-1)/2,K/2)\).
If \(x_v>K\), follow a shortest path to a terminal. Its \(i\)-th edge
carries at least \((K+1)/2^i-1\): start from \(K\), apply the
nondecreasing \(\phi\) and observe that the next bound is below \(K/2\).
The path length is at most \(n\). Taking \(K\ge 2^n(S+1)\) makes even
its last load strictly greater than \(S\). This contradicts the terminal
sum. All switches therefore implement their unsaturated exit law, so
\(x=F(x)\) and Lemma~\ref{lem:real-terminal} gives the exact vector \(Y\).

For ordinary ARRIVAL take \(S=1\), \(m=n\), and read the YES purchase.
Each switch uses \(4n+3\) buyers, \(8n+2\) arc goods and one input
good. Two terminal goods, two terminal buyers and the single dispatch
give \(4n^2+3n+3\) buyers and \(8n^2+3n+2\) goods. Every number has
\(O(n)\) bits. The embedding in Chapter~\ref{sec:bridge} adds one type
per good and one control type to the source types, and two targets to
the private and real targets. This proves the dimensions in
Theorem~\ref{thm:arrival}; a dense game has \(O(n^4)\) payoff entries.
Decoding sums the indicated attacks in their original mass units.

\chapter{Reading a binary terminal count with one chip}\label{sec:multi}
\reviewmark{Count-bit identity, stopping counters and direct recovery of the full vector with binary supplies.}

A switch has two ordered outgoing arcs and uses them alternately, starting
with the first. In ordinary ARRIVAL a single chip starts at an origin, and
the question is whether it ever reaches a designated target. In
Multi-Arrival, prescribed supplies of chips are routed to absorbing
terminals. Every nonterminal has a path to a terminal, and any legal order
of routing ends with the same terminal counts~\cite[Proposition~4.3]{glm}.
We allow an initial supply at any vertex. In the terminal-dispatch
formulation of~\cite{glm}, a terminal initially dispatching \(q\) chips on
two arcs contributes \(\lceil q/2\rceil,\lfloor q/2\rfloor\) initial chips
at the respective successors. With one out-arc it sends all \(q\) there.
Subsequent receipts at the terminal stay there. This conversion uses binary
arithmetic. A transit vertex with one out-arc can use two identical
successors.

\begin{theorem}[Binary Multi-Arrival]\label{thm:multi}
Given a binary Multi-Arrival instance, a terminal, and a nonnegative bit
index \(p\), one can construct in polynomial time a single ordinary ARRIVAL
instance whose answer is bit \(p\) of the terminal count. All terminal counts
reduce to a polynomial batch of independent ordinary ARRIVAL decisions.
Conversely, one binary Multi-Arrival instance with a single chip suffices
for ordinary ARRIVAL.
\end{theorem}
\begin{proof}
Let \(T\) be the sum of the supplied counts and \(k\) the unknown count at
the designated terminal. Bits are indexed from zero. If \(T=0\) or \(p\)
is at least the binary length of \(T\), output a fixed NO instance. Otherwise
put
\[
 M=2^{p+1},\qquad L=\lfloor T/M\rfloor,\qquad
 c=M-1-(T\bmod M),\qquad \tau=T+M/2+c.
\]
Route the \(k\) receipts at the designated terminal into one filter and
the other \(T-k\) receipts into a second filter. Add initial dispatches
of \(M/2\) to the first and \(c\) to the second. A filter consists of
\(p+1\) switches. At every switch odd receipts go to a discard collector,
and even receipts go to the next switch or, at the last switch, to a
monitored collector. Out of \(u\) receipts precisely \(\lfloor u/M\rfloor\)
reach the monitor. The two filters together would yield
\begin{equation}\label{eq:bit}
 Y=\left\lfloor\frac{k+M/2}{M}\right\rfloor+
   \left\lfloor\frac{T-k+c}{M}\right\rfloor
   = L+\bit_p(k).
\end{equation}
Indeed the second term equals \(L-\lfloor k/M\rfloor\) because
\(T+c=LM+M-1\). The shift by \(M/2\) in the first term gives the required
bit. The discard count would be \(\tau-Y\).

We implement these dispatches using a single chip that returns to a common
source between excursions. Set \(B=\lceil\log_2\tau\rceil\). Label \(2^B\)
leaf positions of a complete binary tree in consecutive blocks, one per
positive dispatch of its prescribed size and one for any padding.
Each internal node alternates between its children. Over \(2^B\) source
visits, each leaf is selected exactly once, by induction on the depth.
Only the labels matter. Replace each dyadic subtree wholly within one
block by its destination. At any depth every remaining internal node
contains a strict block boundary, so \(s\) blocks take \(O(1+sB)\) nodes.
Binary interval arithmetic constructs this tree without constructing its
leaves. Padding returns the chip immediately to the source.

At the monitor install a counter sending the chip to YES on receipt \(L+1\).
At the discard collector send it to NO on receipt \(\tau-L\). Both counters
return earlier receipts to the source. For threshold \(q\ge2\) use
\(d=\lceil\log_2q\rceil\) switches as a binary counter initialized to
\(2^d-q\). A zero bit toggles and returns the chip to the source. A one bit
toggles and carries to the next bit, with final carry going to the outlet.
Swapping a switch's two successors sets its initial bit as needed. The
first overflow is on receipt \(q\). A threshold of one uses the outlet
directly.

To verify the construction, temporarily let both outlets return to the
source. Complete one selector period. The excursions process a legal serial
order of the dispatches, so they all finish and their collector totals are
\(Y,\tau-Y\). When \(Y=L+1\), only the monitor reaches its threshold. When
\(Y=L\), only the discard collector reaches its threshold. The actual
one-chip walk agrees with this auxiliary walk until that first overflow.
Make NO an absorbing loop and designate YES as the ARRIVAL target.

The graph uses \(O(|V|+1+p+s\log(1+\tau))\) switches, with \(s=O(|V|+1)\).
Every counter and dispatch size has polynomial binary length. There are at
most \(|\mathrm{terminals}|\,\lceil\log_2(T+1)\rceil\) relevant bits and
the constructions are independent. In the converse direction, collapse
vertices having no directed path to the ordinary target into a second
terminal. Every remaining vertex has a path to a terminal. Injecting one
chip at the original origin gives target count one exactly in a YES
instance. An origin already at a terminal is handled directly.
\end{proof}

\begin{corollary}[One game returns all counts]\label{cor:multi-one}
The complete binary Multi-Arrival output reduces many-one to one
\(\Market\) instance and hence to one one-resource BNE. It also reduces
many-one to a single parallel-road traffic equilibrium. With \(n\)
switches, \(d\) terminals and total supply \(S\ge1\), the direct market
uses \(O(n(n+\log(S+1))+d)\) buyers and goods in total.
\end{corollary}
\begin{proof}
For \(n\ge1\) use the construction of Chapter~\ref{sec:arrival} with
\[
 m=n+\lceil\log_2(S+1)\rceil-1,\qquad K=2^{m+1}.
\]
Then \(m\ge1\) and \(K\ge2^n(S+1)\), so the path bound in that chapter
applies. Read the incoming port and dispatch purchases at every terminal
from the same equilibrium. Lemma~\ref{lem:real-terminal} shows that all
these numbers are the correct integers, even if individual trades are
fractional. There are \(4m+3\) node buyers and \(8m+3\) goods per
switch, one good and one dedicated buyer per terminal, and at most \(n+d\)
dispatch buyers. All capacities, fees, masses, and offer coefficients
have \(O(m+\log(n+1))\) bits, since indegree is at most \(2n\).
One application of Theorem~\ref{thm:bridge} or Theorem~\ref{thm:traffic}
transfers the result. When \(n=0\) the supplies already give the answer,
and \(S=0\) gives the zero vector; either decoder can ignore a fixed
dummy instance.

There is also a separate proof from the bit reduction. Turn each bit
instance of Theorem~\ref{thm:multi} into its own market and take their
disjoint union. Each equilibrium restricts to one on every component,
so all bits are available in one answer. This proof uses one copy per
bit. The direct construction uses one copy per original switch.
\end{proof}

\chapter{One resource and attack-support forests}\label{sec:forests}

Supplying attack probabilities is one way to specify how the types mix.
In the constructions below, listing the positive edges suffices. The
first reduction isolates an equilibrium and gives its forest a canonical
root in each component. The second works on the original positive-effect
table and recovers coverages by looking at component cutoff values.

\section{Two slope profiles and an isolated forest}
\reviewmark{Rectangle embedding, component roots and support reconstruction.}
\begin{theorem}\label{thm:forms}
Each of the following search tasks is many-one equivalent to \(\BNE\).
\begin{enumerate}
\item Find an equilibrium in a one-resource security game with equal
defender gains and prior ratio at most \(1+\omega\) for any given rational \(0<\omega\le1\). The attacker slopes may
be required to equal one except on one rectangle of types and targets,
where they equal two.
\item Find the positive trade edges in a market produced by the isolation
construction of Lemma~\ref{lem:isolate}, or find the positive attack edges
in its one-resource security-game embedding.
\end{enumerate}
In the second task the instance has a unique full equilibrium and its
positive edges form a forest. Both restrictions can be imposed together.
The perturbed market data (or its game embedding) are the input to
forest search. The constructor retains the source instance for the
subsequent repair to that instance.
\end{theorem}
\begin{proof}
Theorem~\ref{thm:compile} first yields one ordinary market from the input
game. Lemma~\ref{lem:rectangle} puts its slopes in the indicated rectangle.
Lemma~\ref{lem:isolate} preserves slopes, isolates the prices and the
allocation, and provides a repair LP. The embedding keeps the slopes on
listed real pairs and appends entries of slope one. Fill any missing pair
first with a negative intercept and its prescribed slope. The bounded
perturbation in isolation keeps these artificial offers negative.

In the isolated market every tight edge carries positive flow. In its
tight forest, deleting an unused tight edge would leave a tree on the
side away from the unique zero-utility or zero-price vertex. All margins
on that side fill, and all its flow is internal. If it contains both kinds
of vertices this contradicts the unequal subset margins in
Lemma~\ref{lem:isolate}. If it contains only one kind, a positive margin
cannot fill. The same reasoning shows that the one marked vertex in each
tree has slack in its margin. To recover that root from the edges, compare
total buyer mass with total capacity in the tree. The larger side contains
the slack vertex. Propagate the
tight equations from one potential. Increasing a buyer utility increases
every buyer utility in the tree and decreases every good price, because
all slopes are positive. The minimum feasible parameter makes a buyer
zero and the maximum makes a good zero. The side just identified therefore
determines the endpoint imposed by nonnegativity within the tree. This
recovers its potentials. Each leaf away from the root puts its remaining
margin on its sole edge. Successive leaf deletions recover all flows.
Single vertices have potential zero.

In the game embedding each real dedicated type and each private target
adjoins a leaf. An underfilled buyer attacks the common outside target.
At most one such buyer belongs to each market tree, so these edges cannot
create a cycle. The control type and its target form a separate edge.
Deleting the auxiliary vertices recovers the market forest, and all
remaining mixed strategies follow from the embedding. In the reverse direction, query the game embedding of the supplied
isolated market, or the supplied game itself. In the market case retain
only the decoded real trades. In either case their positive edges are
read from one equilibrium without knowing how the instance was isolated.
\end{proof}


\begin{corollary}\label{cor:close-slopes}
Fix any rational \(k>1\). The exceptional rectangle in the one-resource
reduction can instead have slope \(k\), with every other slope still one.
In particular the two slope profiles can be required to take values in
\([1,1+\kappa]\) for any fixed rational \(\kappa>0\).
\end{corollary}
\begin{proof}
In Lemma~\ref{lem:rectangle}, replace each unary doubling gate by a
multiplication-by-\(k\) gate. Its buyer now uses slope \(k\) on the
exceptional input good and one on its ordinary output good. All other
gates, including minimum, maximum, translation and the approximate
mean, use only unit slopes. For a rescaled source coefficient
\(\beta\ge1\) use the consecutive multiples \(k^jP,k^{j+1}P\) with
\(k^j\le\beta\le k^{j+1}\). Their interpolation weight is the rational
\((\beta-k^j)/(k^{j+1}-k^j)\). Truncate its binary expansion so that
the weight error times \(H k^{j+1}\) is below a quarter of the
source tolerance, where \(P\in[0,H]\). The same repeated averages and
support repair as in the lemma then implement this edge to the required
precision. Since \(k>1\) is fixed, \(j=O(1+\log\beta)\), the powers
have polynomial binary length, and all gates have polynomial-bit data.
All listed slopes are at least one, so Chapter~\ref{sec:bridge} embeds
the resulting market with its slopes intact and unit-slope auxiliary
edges. Take \(k=1+\kappa\) for the last assertion.
\end{proof}

\section{A forest on the original payoff table}\label{sec:same-forest}
\reviewmark{Cutoff cases, component endpoints and the positive-support test.}

The isolation construction changes the instance so that its equilibrium
is unique. There is also a support certificate on an unmodified game.
Here \(p_i>0\), each defender gain \(d_j>0\), and each attacker slope
\(b_{ij}>0\). The certificate lists type--target edges alone.

\begin{theorem}\label{thm:givenforest}
On these positive-effect games, finding a forest that is the exact
positive attack support of some equilibrium is many-one equivalent to
finding an equilibrium of the same game. Given any proposed edge set
\(S\), one can test in polynomial time whether it is the exact positive
support and recover a compatible full profile when it is.
\end{theorem}
\begin{proof}
Starting from an equilibrium, cancel cycles of positive joint attack
mass by alternately adding and subtracting the minimum mass on the
decreasing edges. This preserves every row total, every target load and
all best responses. Each cancellation deletes a positive edge, giving a
forest of at most \(m+n-1\) edges.

For the decoder, impose \(v_i\ge a_{ij}-b_{ij}x_j\) on every pair,
equality on \(S\), zero flow off \(S\), and row sums \(p_i\). Reject
if \(S\) omits a row. A nontrivial connected component \(C\) of \(S\)
has a one-dimensional or singleton set of local potentials after imposing
its edge equalities and \(0\le x\le1\), unless these are infeasible.
Indeed a reference coverage determines all potentials by propagation,
and increasing the reference strictly increases every coverage in the
component and decreases every row utility. Minimize and maximize the
reference by LP to obtain the two endpoint coverages.

Set \(w_j=\sum_i f_{ij}\). For a cutoff \(\lambda\), positive coverage
requires \(d_jw_j\ge\lambda\), and coverage below one requires
\(d_jw_j\le\lambda\). In \(C\) the total mass is known. Define
\[
 h_C=\frac{\sum_{i\in C}p_i}{\sum_{j\in C}1/d_j}.
\]
If \(\lambda<h_C\), some \(d_jw_j>\lambda\), forcing \(x_j=1\) and the
upper endpoint of \(C\). If \(\lambda>h_C\), a strict deficit forces
\(x_j=0\) and the lower endpoint. When \(\lambda=h_C\) and the segment
has positive length, a strict surplus would imply a strict deficit
elsewhere by summing the loads. Those two targets would force the
reference to its two distinct endpoints simultaneously. Hence all
\(d_jw_j=\lambda\) in such a component. A singleton has fixed coverages
on which to impose the cutoff inequalities. An isolated target has zero
load, so it must be uncovered if \(\lambda>0\).

Sort \(0,\max_jd_j\) and the component breakpoints in that range. Test
each breakpoint and the intervals between them. On an interval use the
forced endpoints and their cutoff inequalities. At each breakpoint use
the equality rule on flexible tied components. The interval LPs may use
closed intervals because their respective conditions still imply the
cutoff inequalities at the ends. In each case impose the global
comparisons, budget and flow constraints, and maximize \(0\le t\le1\)
subject to \(f_{ij}\ge t\) on \(S\). An optimum \(t>0\) in some case
is necessary and sufficient. The number of rational LPs is linear in the
number of components. Divide the joint masses by the priors and realize
the coverage as a lottery on pure defenses to obtain a full profile.
\end{proof}

\section{Testing a prescribed forest with flat offers}\label{sec:flat}
\reviewmark{Vertex Cover construction and membership in NP.}

The next reduction has a different output. The input already specifies
the forest, and the question is whether \emph{that forest} can be realized.
A single coverage-independent option in a payoff row changes this
recognition problem.

\begin{theorem}\label{thm:flat}
Given a game and \(S\), each of the following decisions is NP-complete:
whether an equilibrium uses only \(S\), whether it does so with every
edge of \(S\) tight, and whether its exact positive support is \(S\).
This holds with unit defender gains, at most one flat offer per type,
unit slope elsewhere, and \(S\) a disjoint union of six-edge paths and
single edges. The reduction also works if a prescribed cutoff is part
of the input.
\end{theorem}
\begin{proof}
Reduce Vertex Cover~\cite{karp} on a graph with \(n\ge2\) vertices and a
nonempty edge set, asking for at most \(k\) vertices, \(0\le k\le n\).
List each edge twice and let \(m\) be the even length of this list.
For each vertex \(j\), make \(m\) targets \(v_{j,1},\ldots,v_{j,m}\).
For listed edge \(i=\{j,\ell\}\) call \(v_{j,i},v_{\ell,i}\) its two
occurrences and add a hub \(h_i\). All other copies are padding. Add
an anchor \(a\) and give the defender \(km+m/2+1\) resources.

Use raw type weights with \(\theta=1/8\) and normalize by their total
\(W\). Every unspecified payoff at \(t\) is \(-1-x_t\).
\begin{itemize}
\item For each occurrence \(v\) of edge \(i\), a router of weight
\(3\theta/2\) gets \(1\) at \(v\), \(3/2-x_{h_i}\) at its hub, and
\(2-x_a\) at the anchor. Its allowed edges in \(S\) go to \(v,h_i\).
\item Each copy \(v\) has a private type with offers \(1-x_v\) and
\(1-x_{v'}\) at itself and the next copy of the same vertex in cyclic
order. Only \(v\) is in \(S\). The weight is \(\theta/4\) at an
occurrence and \(\theta\) at a padding copy.
\item A type of weight \(\theta\) has offer \(1-x_a\) at the anchor
and only that edge in \(S\).
\end{itemize}
The two routers and their occurrence targets connect the two private
types through their shared hub, giving exactly six edges. Different
listed edges have distinct occurrences.

In any equilibrium using only \(S\), private-type comparisons around
each cycle force the copies of \(j\) to share a coverage \(x_j\).
Routers can get at most \(3/2\) on their allowed targets, so
\(2-x_a\le3/2\). The positive anchor coverage and its sole allowed load
\(\theta\) bound the raw cutoff by \(\lambda\le\theta\).
If \(x_{h_i}<1/2\) both routers strictly prefer the hub among their
allowed choices. Its load \(3\theta>\lambda\) would force coverage one.
At \(x_{h_i}=1/2\) the hub receives only \(\lambda\le\theta\). At least
one occurrence then gets at least \(\theta\) from a router plus positive
private mass, forcing full coverage. Above \(1/2\) both routers strictly
prefer their occurrences, likewise forcing full coverage. In each case
some router uses an occurrence of utility one, and its anchor comparison
implies \(x_a=1\). Subtracting the anchor and the \(m\) half-covered or
better hubs from the budget gives \(\sum_jx_j\le k\).
The fully covered vertices meet every input edge.

Conversely enlarge a cover to size \(k\), cover its copies and the
anchor fully, and set each hub coverage to \(1/2\). Each private type
uses its own copy and the anchor type uses the anchor. For a listed edge
with one selected endpoint send \(4\theta/3\) from its router to that
occurrence and \(2\theta/3\) from the other router to the unselected
one. If both endpoints are selected, each occurrence receives \(\theta\)
from its router. Send the rest to the hub. All specified edges are used
at a best response, and each hub has load \(\theta\). An uncovered
occurrence has load \(11\theta/12<\theta\), while each covered occurrence
has load greater than \(\theta\). The remaining targets have load
\(\theta\), so the defender best responds with raw cutoff \(\theta\).
This also proves the claim for the specified normalized cutoff \(\theta/W\).

For membership in NP guess the positive edges and each target's
zero, interior or full coverage status. The resulting best-response,
flow and cutoff conditions are rational linear constraints. Enforce a
required strict inequality with a common positive slack, maximized in
a bounded LP. A feasible strict pattern has a polynomial-bit witness.
The reduction concerns an extra input \(S\) on either answer of Vertex
Cover. A search oracle for an arbitrary equilibrium has no obligation
to realize the requested \(S\).
\end{proof}

\chapter{Games with two attainable targets per type}\label{sec:sparse}
\reviewmark{Control scores, port loads at both price bounds and degree counts.}
A target is \emph{attainable} for a type if it maximizes that type's payoff
at some legal coverage, including a tie.

\begin{theorem}\label{thm:degree}
General BNE search reduces to one game with at most two attainable targets
per type and at most three types for which any target is attainable. The
attacker slopes equal one except on a rectangle of twos. The number of
defender resources in this construction depends on the instance.
\end{theorem}
\begin{proof}
Start with the rectangular-slope market of Lemma~\ref{lem:rectangle}.
Give buyer \(i\) a private outside good of capacity \(2m_i\) and utility
equal to minus its price, and require every buyer to spend in full.
For each of buyer \(i\)'s \(k_i\) options, including that good, make a
separate buyer of mass \(m_i\) choosing between this option and a shared
hub with offer \(C_i-P_{\rm hub}\). Choose
\(C_i>1+\max(0,\max_j a_{ij})\). One more buyer of mass \(m_i\) uses
only the hub, which has capacity \(k_i m_i\).

For each good of capacity \(c\) with \(d\ge1\) incident buyers, make \(d\)
copies of capacity \(c\) on a path, one per incident buyer. A new full
buyer of mass \(c\) between consecutive copies chooses the cheaper one,
with unit slopes. Drop goods with no incident buyer. Put all prices in
a common box \([0,R]\), with \(R\) larger than every \(C_i\), every positive
\(a_{ij}/b_{ij}\), and one. Give each bridge and each single-option buyer
intercept \(R+2\).

The box conditions require load at least capacity above price zero and at
most capacity below price \(R\). To realize them, normalize all masses
and capacities by the total full buyer mass. Use each physical good as
a target with coverage \(P/R\). Divide its offers by \(R\) after this
substitution to preserve their slopes. Give the full buyers half the
total prior in their relative proportions and each such target defender
gain \(2/c\) for its normalized capacity \(c\). For \(N\) physical
targets add \(N+2\) control targets, each with a strictly dedicated type
of prior \(1/(2(N+2))\) and defender gain \(2(N+2)\). The budget is \(N+1\).
Control targets have score one and their total coverage lies between
\(1\) and \(N+1\), forcing the cutoff to be one. This gives the box
conditions exactly. Unintended offers have intercepts below every
intended payoff on the box. Assign each copy of an exceptional good to
the exceptional group and each option's buyer to its source buyer's
group. Hubs and outside goods are ordinary, and bridge and control rows
have unit slopes. Each physical target meets at most one external buyer
and two bridges.

To decode, replace the prices on each path by their maximum \(Q\). A
maximal contiguous run of \(l\) copies with price \(<Q\) is below \(R\)
throughout. Its \(l-1\) internal bridges send all their mass there, as
does at least one boundary bridge. Thus it fills with bridge flow and
has no external purchases. Its total load is \(lc\): all its copies have
price below \(R\), and its bridges already supply at least \(lc\).
Thus every lower-priced copy fills, even at price zero. Every used
external option keeps its price when set to \(Q\), while other options
weakly worsen. If \(Q<R\), all \(d\) copies have load at most \(c\).
If \(Q>0\), the copies at \(Q\) have load at least \(c\), and every
other copy fills. Subtracting the bridge mass \((d-1)c\) gives the
respective upper and lower bounds \(c\) on external load.

A zero hub price would attract all \(k_i+1\) buyers in its group and
exceed capacity. Hence the hub receives at least \(k_i m_i\), leaving
at most \(m_i\) for all external options combined. The private outside
good has price zero since its capacity is \(2m_i\). If the hub price is
below \(R\), it fills exactly. If it equals \(R\), the outside-option
buyer strictly prefers its zero offer to \(C_i-R\) and spends all \(m_i\)
externally. Thus exactly \(m_i\) goes to external options in either case.
If an external offer strictly beats the hub, that buyer spends the whole
\(m_i\) there and no other external offer can strictly beat the hub.
Otherwise every used external option ties the hub. This proves optimality
for the original buyer. At a real price \(R\) all original offers are
strictly negative and the outside is available, so the required positive
capacity could not fill. The real goods therefore clear in the ordinary
sense. Removing the private outside purchases recovers the market
allocation. All tables have polynomial binary length, and the reverse
reduction is inclusion.
\end{proof}

\section{Adding a prescribed graph of two-target types}\label{sec:overlay}
\reviewmark{Attainability, normalized loads and fixed-coverage repair.}

In a game with at most two attainable targets per type, draw an edge
between the two targets attainable by each two-option type. Counting
these edges gives a graph on targets. Its topology can be prescribed in
part while retaining a decoder for the original game.

\begin{lemma}\label{lem:overlay}
Let \(0<r<n\). Given a game and any simple graph \(H\) on its \(n\)
targets, one can add a type for each edge of \(H\), attainable exactly at
its two ends, so that any equilibrium of the augmented game yields a full
old equilibrium at the \emph{same coverage}. Old types retain their
attainable targets.
\end{lemma}
\begin{proof}
Omit zero-prior rows until recovery. If \(s=|E(H)|>0\), give each new
type raw weight \(\eta/s\), with the old raw priors still summing to one.
Normalize all weights together. The type for \(uv\) has offers
\(-x_u,-x_v\) at its ends and \(-2-x_j\) elsewhere. At uniform coverage
both ends tie and a small budget-preserving change favors either one.
Nothing changes an old row's comparisons. If \(s=0\) the original game
already suffices.

Put \(d_*=\max_j d_j\). In the answer rescale joint flows and the cutoff
back to raw units. Write \(w_j\) for the old-row loads. Removing the new
rows changes any score by at most \(d_*\eta\). After clamping the raw
cutoff to \([0,d_*]\), the old rows satisfy their original cutoff
inequalities up to \(d_*\eta\). The upper bound after clamping follows
also from \(d_jw_j\le d_*\).

Fix the supplied coverage and record the support of the old rows.
Minimize a common slack \(0\le t\le1\) in the original score inequalities,
with variables for old flows and a cutoff in \([0,d_*]\). Impose old row
sums and prohibit flows outside the recorded support. Which score rows
are needed depends only on whether each coverage is zero, interior or
full. Clearing these finitely many possible row types gives an
input-computable determinant bound \(Q\) as in Lemma~\ref{lem:repair},
independent of the returned coverage. Use
\(\eta=1/(2Q\max(1,d_*))\). The raw rows certify \(t<1/Q\), hence the
minimum is zero. Every allowed positive pair is still optimal at the
same coverage. The repaired rows and the supplied defense therefore
give an original equilibrium. Restore best replies to zero-prior types.
\end{proof}

For instance, apply the lemma after Theorem~\ref{thm:degree} and use an
explicit expander of fixed even degree on the targets for \(H\)
\cite[Theorem~1.3]{alon}. For example take degree four and an error parameter below
\(4-2\sqrt{3}\) in that theorem. Even degree permits every sufficiently
large target count. Deleting loops preserves cut sizes, and merging
parallel edges loses at most a factor four. Each target gains at most \(\deg_H(j)\) attainable types.
The resulting unweighted graph contains every edge of the expander, so
its edge boundary across any cut is at least the expander's boundary.
The sizes that precede the explicit expander family's threshold can use
complete graphs, a fixed bound on degree there. The original instance
still decodes even though the number of new edges crossing a cut can be
large and their \emph{total prior mass} is small.

\chapter{Uniqueness of a security-game equilibrium}\label{sec:unique}

The input game is as in Chapter~\ref{sec:models}. Two full profiles
can differ either in their marginal coverage and attack strategies or in
the defender lotteries that realize the same coverage. A uniqueness
test must detect both forms of nonuniqueness.

\begin{theorem}[Testing and witnessing uniqueness]\label{thm:unique}
Given the game and one rational equilibrium, a deterministic
polynomial-time algorithm either certifies that it is the only real
equilibrium or returns a different rational equilibrium. This holds for
marginal profiles and for full profiles presented with a support list for
the defender lottery. Consequently the corresponding problem on the game
alone requires at most one call to an exact equilibrium algorithm.
\end{theorem}

\section{Why changing one side suffices}
\reviewmark{Conditional complementarity system and a direct proof of bounded connectedness.}

\begin{lemma}\label{lem:connected}
The set of marginal equilibrium profiles \((x,\sigma)\) is connected,
including games with payoff ties, zero slopes, zero gains, or zero priors.
\end{lemma}
\begin{proof}
For \(r=0,n\) the coverage is fixed and the attack strategies form a product
of simplices. Suppose \(0<r<n\). We use the unperturbed game
complementarity system from~\cite[Lemmas~5.14--5.15]{book}, in conditional
coordinates. Set \(A_{ij}=a_{ij}+1+\max_k(b_{ik}-a_{ik})\), so
\(A_{ij}-b_{ij}x_j\ge1\) on the box. The notation \(v\perp w\) requires
\(v,w\ge0\) and \(vw=0\). Consider
\begin{equation}\label{eq:compact}
\begin{aligned}
 \sigma_{ij}&\perp u_i+b_{ij}x_j-A_{ij},&
 t_j&\perp 1-x_j,\\
 \lambda&\perp r-\textstyle\sum_jx_j,&
 u_i&\perp 1-\textstyle\sum_j\sigma_{ij},\\
 x_j&\perp \lambda-1+t_j-\pi_j(\sigma).&&
\end{aligned}
\end{equation}
Its matrix in the variable order \((\sigma,t,\lambda,u,x)\) is \(P_0\),
meaning that every principal minor is nonnegative. For positive priors
this follows from~\cite[Lemma~5.15]{book} by substituting the joint mass
\(p_i\sigma_{ij}\) in each attack column and dividing each utility-slack
row \(i\) by \(p_i\). The positive row and column scalings preserve the
sign of every principal minor. The entries of the displayed matrix depend
continuously on the prior, which extends the conclusion to priors with
zero entries.

Every solution has \(0\le x\le1\) and \(u_i\ge1\), so each attack row
sums to one and attains its best value. Some \(x_k<1\) since
\(\sum_j x_j\le r<n\). Hence \(t_k=0\) and \(\lambda-1\ge\pi_k\ge0\),
which forces \(\sum_j x_j=r\). Some \(x_j>0\) now gives
\(\lambda-1\le\pi_j\le\max_k d_k\). The score and cap rows give
\(t_j=(\pi_j-(\lambda-1))_+\). Together with
\(u_i=\max_j(A_{ij}-b_{ij}x_j)\), these formulas bound the entire solution
set. Its projection onto \((x,\sigma)\) is precisely the equilibrium set.
Indeed, the last row and the caps say that positive coverages have score at
least the cutoff \(\lambda-1\) and nonfull coverages have score at most that
cutoff. This is sufficient for optimal coverage by an exchange of mass.
Conversely, at an optimal coverage take the \(r\)-th largest score as
cutoff, set \(t_j\) and \(u_i\) as above, and obtain a solution
of~\eqref{eq:compact}.

Finite-game existence supplies at least one equilibrium, filling in
zero-prior best responses if necessary. Jones and Gowda~\cite{jg} prove
that a nonempty bounded solution set of a \(P_0\) linear complementarity
problem is connected. Its continuous projection gives the claim.

For completeness, the bounded case used here also follows directly from
diagonal perturbation. Write \(M,q\) for the LCP matrix and offset, and
introduce the continuous map on all of \(\mathbb R^N\)
\[
  H_\varepsilon(y)=(M+\varepsilon I)[y]_+-[-y]_+ +q .
\]
The maps \(z=[y]_+\), \(w=[-y]_+\) and \(y=z-w\) give a continuous
bijection between the zero set of \(H_0\) and the LCP solutions. The bounds above make that zero set compact. For \(\varepsilon>0\), \(M+\varepsilon I\) is a P-matrix by
principal-minor expansion. Its LCP has a unique solution for every offset
\cite[Proposition~2.6]{book}. Thus \(H_\varepsilon(y)=b\) has a unique
solution for every \(b\). To see continuity of the inverse, observe that on each of the
finitely many closed orthants the map is linear up to translation with
nonsingular linear part, since its determinant is a principal minor of
\(M+\varepsilon I\). The affine inverses agree on common boundaries.

If the compact zero set split into two nonempty disjoint compact pieces,
take disjoint bounded open neighborhoods of them. On the boundary of
their union \(\|H_0\|\) has positive minimum. Choose a zero in each
piece. Uniform convergence on the closures makes the images of these two
points under \(H_\varepsilon\) arbitrarily close to zero, while the
boundary images stay bounded away from zero. The inverse image of the
line segment joining the two images is a continuous path between the two
neighborhoods. It must cross their boundary, which the bounds exclude.
Hence the zero set, and its projection onto profiles, are connected.
\end{proof}

Fix nonempty attack menus \(S_i\) and a nonempty face \(D\) of \(H_r\).
Let \(X\) be the points of \(D\) at which every \(j\in S_i\) is a best
target for type \(i\). Let \(Y\) be the attack distributions supported on
the \(S_i\) for which every point of \(D\) is optimal for the defender.
Both are closed polytopes. For \(Y\) one can compare vertices of \(D\) to
vertices of \(H_r\) to see linearity. Every point of \(X\times Y\) is an
equilibrium. Conversely, take the actual attack supports at an equilibrium
and the smallest face containing its coverage. The defender's maximizing
face contains this smallest face. Thus the equilibrium set is a finite
union of such products.

If a supplied point \(e\) were alone in each product containing it, the
finite union of all remaining closed products would separate it from any
other equilibrium. By Lemma~\ref{lem:connected}, whenever there is a
second equilibrium some product containing \(e=(x^*,\sigma^*)\) contains
another point \((x',\sigma')\). That product also contains
\((x',\sigma^*)\) and \((x^*,\sigma')\). At least one of these differs
from \(e\). This observation lets us avoid enumerating faces or menus.

\section{Two polynomial-size tests}
\reviewmark{Both fixed-strategy LPs and their marginal uniqueness guarantee.}

Hold \(\sigma^*\) fixed. Impose \(x\in H_r\),
\(\pi(\sigma^*)\cdot x=\pi(\sigma^*)\cdot x^*\), and
\[
  U_{ij}(x)\ge U_{ik}(x)\quad
  (i\in[m],\ j\in\supp(\sigma_i^*),\ k\in[n]).
\]
This describes exactly the equilibrium coverages compatible with the
supplied attacks. Next hold \(x^*\) fixed. Require \(\sigma_i\) to be a
probability distribution on the maximizing targets for \(U_i(x^*)\), and
impose
\[
       \pi_j(\sigma)\ge\pi_k(\sigma)
       \quad\text{whenever }x^*_j>0,\ x^*_k<1 .
\]
These conditions express exactly the absence of an improving transfer of
coverage from \(j\) to \(k\). In each polytope minimize and maximize each
coordinate. In at most \(2n(1+m)\) rational linear programs, we either
find a different rational point or establish that both polytopes are
singletons. The preceding argument certifies global marginal uniqueness
in the latter case.

\section{Different lotteries with the same coverage}
\reviewmark{Hypersimplex faces and construction of a second full profile.}

Let \(x\) be a rational coverage and \(\nu\) the supplied lottery, with
distinct positive-weight supports. If a marginal test has already found a
different profile, the new coverage can be realized as a sparse lottery.
For completeness, lay consecutive intervals of lengths \(x_1,\ldots,x_n\)
on \([0,r)\). For \(a\) uniform on \([0,1)\), select the targets whose
half-open intervals contain one of \(a,a+1,\ldots,a+r-1\). Each interval
has length at most one, so exactly \(r\) targets are selected. Target \(j\)
is selected with probability \(x_j\). The fractional parts of the interval
endpoints split \([0,1)\) into at most \(n\) pieces and give a polynomial-time
rational support list~\cite[Lemma~2.2]{book}. Endpoint budgets have just
one pure defense.

It remains to test a fixed \(x\). Coordinates equal to zero or one are fixed
in every pure defense in the support. Write \(k\) for the number of
fractional coordinates and \(s=r-|\{j:x_j=1\}|\). The possible support
vertices are the \(s\)-subsets of those \(k\) coordinates, with the fixed
coordinates reattached. For \(k=0\) the face is a point. For \(s=1\) or
\(s=k-1\) it is a simplex and its barycentric coordinates are unique.

Otherwise \(2\le s\le k-2\) and the face has \(\binom{k}{s}>k\) vertices.
If the supplied lottery uses more than \(k\) vertices, the first \(k+1\)
incidence columns restricted to the fractional coordinates have a nonzero
rational linear dependence. Its coefficients sum to zero since all columns
have sum \(s>0\). Move the probabilities along this dependence until one
first hits zero. This changes the lottery and preserves its marginals.

If the support has size \(h\le k\), enumerate the first \(h+1\) distinct
\(s\)-subsets to find a missing vertex \(v\). On the fractional coordinates
take \(\eta\) to be the minimum of \(x_j\) over \(v_j=1\) and \(1-x_j\)
over \(v_j=0\). Then \(0<\eta<1\) and
\((x-\eta v)/(1-\eta)\in H_r\), including the fixed coordinates.
Realize this point sparsely, give its lottery weight \(1-\eta\), and give
\(v\) weight \(\eta\). This constructs a second lottery with the same \(x\).
All operations have polynomial bit cost in the game and the supplied
profile. This completes the proof of Theorem~\ref{thm:unique}.

\chapter{Decision problems that encode an equilibrium}\label{sec:decisions}
\reviewmark{Decision quantifiers, clone interval, batch recovery and the unambiguous witness.}

Let \(\Unique\) ask whether a specified game has exactly one full real
equilibrium. By Theorem~\ref{thm:unique}, one search call followed by
rational LP decides \(\Unique\) from the game alone. The other direction
can use a particularly restricted yes/no question.

\begin{theorem}\label{thm:shape}
Every language in \(\mathsf P^{\BNE}\) reduces many-one both to
\(\Unique\) and to its complement (on valid games). The constructed game
has one resource and a unique defender strategy. Its full equilibrium set
is either a point or a closed line segment. The four attack probabilities
that vary on the segment are specified by the construction.
\end{theorem}
\begin{proof}
An oracle decision algorithm must give the same bit on every valid
transcript. Replace each game call by its adaptive market algorithm. Apply
Theorem~\ref{thm:compile}, retaining its two output goods for the decision
bit \(b\) and its complement. By a translated identity gate give them
prices \(p_A=1+b,\ p_B=2-b\). Include the following buyers in the reader
budget when sizing the gates. In the last module add a good \(T\) of
capacity \(3h/2\) and two buyers \(A,B\) of the same mass \(h>0\). Each
has an offer \(4-P_T\) and an offer \(4-p_A\) or \(4-p_B\) on its own
comparison good. Both spend in full. Below \(\min(p_A,p_B)\) good \(T\)
would be overdemanded. Above it at most \(h\) could reach \(T\), which
must fill at a positive price. Hence \(P_T=1\). The buyer with the higher
comparison price sends all \(h\) to \(T\). The other puts \(h/2\) on each
of its two options. Buyer \(B\) splits precisely when \(b=1\).

Isolate this entire market by Lemma~\ref{lem:isolate}. Every zero flow in
the isolated solution stays zero on repair, and every positive isolated
edge stays tight on repair. When \(b=1\) neither purchase by \(B\) can
vanish, since all original solutions use both. When \(b=0\) its comparison
offer is strictly suboptimal in every original solution and cannot be used
after isolation. A zero isolated utility would stay zero on repair, so \(B\)
still fills its entire mass \(h'\). Split \(B\) into two labeled buyers
of mass \(h'/2\) with identical offers. Every equilibrium merges to the
unique isolated solution. If the merged purchase of \(T\) is \(z\), the
first clone's purchases there can be any amount in
\[
       [\,\max(0,z-h'/2),\ \min(z,h'/2)\,].
\]
The other three purchases follow from the row and column totals. The
interval has positive length exactly when \(0<z<h'\), that is, when \(b=1\).
Prices and utilities are fixed. The one-resource embedding from
\cite[Theorem~3.16]{book} preserves these profiles bijectively, scaling each
source row by its fixed game weight. It fixes the defense and leaves only
the indicated two-by-two block of attack probabilities free. Complement
the output bit to exchange the two conventions. Theorem~\ref{thm:unique}
gives the converse oracle bound.
\end{proof}

\begin{corollary}\label{cor:batch}
A polynomial batch of independent point-versus-segment queries recovers
an exact full equilibrium of any input game. A single such decision can
also encode an ARRIVAL answer. Both full-profile uniqueness and
nonuniqueness belong to \(\mathsf{UP}\cap\mathsf{coUP}\).
\end{corollary}
\begin{proof}
For search, compile the input to one market and isolate it. For every
listed edge, construct the predicate ``its flow in this isolated market
is positive.'' One BNE suffices to evaluate each predicate under the
market embedding, and uniqueness makes the answer independent of the
oracle. Apply Theorem~\ref{thm:shape} to every predicate before reading
any answer. The resulting forest recovers the full profile by
Theorem~\ref{thm:forms}. For ARRIVAL use the one-query reduction of
\cite[Theorem~4.1]{book} as the predicate instead. For the unambiguous
upper bound, guess the unique solution of the genuine P-LCP produced
from the input game by~\cite[Theorem~5.19]{book}. Each coordinate of that
vector is a polynomial-bit rational by the basis determinant bound.
Require a canonical reduced-fraction encoding, verify complementarity,
decode with a fixed deterministic LP algorithm, and apply
Theorem~\ref{thm:unique}. Exactly one witness certifies the correct side.
\end{proof}

There is also a decision problem on an individual edge. Keep buyer \(B\)
unsplit in the isolated market in the proof of Theorem~\ref{thm:shape},
embed it, and designate its real attack edge to the comparison good.
Every such game has a unique positive attack forest by the forest argument
of Theorem~\ref{thm:forms}. The designated edge occurs exactly when
\(b=1\). One search call reads its membership in the reverse direction.
Thus deciding membership of one named edge on this promise is complete
for \(\mathsf P^{\BNE}\) under polynomial-time many-one decision reductions.
Indeed \(\mathsf P^{\BNE}\subseteq\mathsf{UP}\cap\mathsf{coUP}\), by
Theorem~\ref{thm:shape} and Corollary~\ref{cor:batch}.

\chapter{P-matrices and contraction fixed points}\label{sec:upper}

\(\PLCP\) asks for \(z,w\ge0\) with \(w=Az+q\) and \(z_kw_k=0\) for all
\(k\), where every principal minor of \(A\) is positive. The direct
game construction has dimension \(mn+m+2n+1\) for \(m\) types, \(n\)
targets and an interior budget. Endpoint budgets give fixed defenses.
Fearnley, Gordon, Mehta and Savani place P-LCP in \(\mathsf{UEOPL}\),
Unique End of Potential Line~\cite{fgms}. Affine matching also has a
direct reduction to P-LCP~\cite[Theorem~5.6]{book}.

\section{The complementarity pairs of the game}
\reviewmark{Signs of all principal minors, the uniform coordinate bound and exact support repair.}

For an interior budget \(0<r<n\), shift each type's intercepts by a
row constant to \(A_{ij}\) so that \(A_{ij}-b_{ij}\ge1\). This preserves
best responses. Add one to every defender score, which changes its
objective on the budget hyperplane by the constant \(r\). In the direct
construction of \cite[Theorem~5.19]{book} write \(w_j=\sum_i f_{ij}\).
The nonnegative variables and paired nonnegative slacks are
\[
\begin{array}{c|l}
 \text{variable}&\text{slack}\\ \hline
 f_{ij} & u_i+b_{ij}x_j-A_{ij}\\
 t_j & 1-x_j\\
 \lambda & r-\sum_jx_j\\
 u_i & p_i-\sum_jf_{ij}\\
 x_j & \lambda-1+t_j-d_jw_j .
\end{array}
\]
Each paired product must vanish. Positive \(u_i\) forces the prescribed
row sum, and \(f_{ij}>0\) makes the chosen attack optimal. If
\(\lambda=0\), every last slack would require \(t_j>0\), so every \(x_j\)
would be one, exceeding \(r\). Thus the budget is tight. The cap
multipliers and final slacks are precisely the defender cutoff
conditions. A full profile follows by realizing the coverage and choosing
best responses for any zero-prior rows.

In the variable order \((f,t,\lambda;u,x)\) the coefficient matrix is
\(M=\left(\begin{smallmatrix}0&B^{\mathsf T}\\-D&0\end{smallmatrix}\right)\).
Rows of \(B,D\) correspond to types and targets. In attack column \(ij\),
the two entries in \(B\) are \(1,b_{ij}\), and those in \(D\) are
\(1,d_j\). A cap column is \(-1\) in its target row. The budget column
is \(-1\) in every target row. These last two kinds of columns agree in
both matrices. Unspecified entries vanish.

\begin{lemma}[Principal minors of the game matrix]\label{lem:game-minors}
Every principal minor of \(M\) is nonnegative, including at zero slopes
or defender gains.
\end{lemma}
\begin{proof}
If a principal selection takes different numbers of coordinates from
the two blocks, the off-diagonal matrix has deficient rank. With equal
numbers the determinant is \(\det B'\det D'\), where the primes denote
the same square selection of rows and columns in the two arrays. The
block-column swap and the signs on \(-D'\) cancel.

First assume \(d_j>0\). Prove that product nonnegative by induction
on the selection size, with empty product one. If a column of \(D'\)
is zero, the conclusion follows. If it has exactly one entry, \(B'\)
has nonzeros only in the same position, and the two entries have
nonnegative product. Expand there in both determinants and use induction.
Otherwise every selected attack column retains both its ends, no cap
column occurs, and there may be one budget column. Put \(y_i=1\) on
selected type rows and \(y_j=-1/d_j\) on selected target rows.
Then \(y^{\mathsf T}\) annihilates every attack column of \(D'\).
Without a budget column \(\det D'=0\). With that column in position \(k\),
\[
 y^{\mathsf T}D'=\beta e_k^{\mathsf T},
 \qquad \beta=\sum_{\text{selected targets }j}1/d_j>0 .
\]
Multiplying by \(\operatorname{adj}(D')\) gives, for each row \(j\),
\(\det D'\,y_j=\beta\,\operatorname{cof}_{jk}(D')\).
Expand \(\det B'\) along its budget column. On target row \(j\),
its entry is \(-1\), and the corresponding term of the product is
\[
 \frac{-\beta}{y_j}\operatorname{cof}_{jk}(B')\operatorname{cof}_{jk}(D')
 \ \ge\ 0 .
\]
The coefficient is positive and the cofactor signs cancel, leaving the
smaller determinant product covered by induction. All other terms are
zero. Finally let positive gains decrease to the required nonnegative
values; the determinants are continuous. Zero attacker slopes were
already allowed in the induction.
\end{proof}

For every \(\delta>0\), expansion in the diagonal columns gives
\[
 \det(M_{JJ}+\delta I)=\sum_{H\subseteq J}\delta^{|J|-|H|}\det M_{HH}>0 ,
\]
where the empty minor equals one.
Thus \(M+\delta I\) is a genuine P-matrix and its LCP has a unique
solution~\cite{book}. For \(0<\delta\le1\) put
\[
 U=\max_{ij}A_{ij},\qquad W=1+mU,\qquad V=1+W\max_j d_j,\qquad
 K=\lceil\max(U,W,1+V)\rceil.
\]
The perturbed slack of a variable is its displayed slack plus \(\delta\)
times that variable. If \(u_i>0\), its row spends \(p_i+\delta u_i\),
so a positive \(f_{ij}\) implies \(u_i\le A_{ij}\le U\). Thus total flow
is at most \(W\). If \(t_j>0\) then \(x_j=1+\delta t_j>0\), whose
tight slack gives \(t_j\le1+d_jw_j\le V\). Hence \(x_j\le1+V\).
If \(\lambda>0\), its tight slack forces positive total coverage, so
a positive \(x_j\) similarly implies \(\lambda\le V\). All coordinates
lie in \([0,K]\).

On the original coefficients choose a support \(J\), set \(z_k=0\)
off \(J\), impose \(q+Mz\ge -t\mathbf1\) and
\((q+Mz)_k\le t\) on \(J\), and minimize \(0\le t\le1\) in the box
\([0,K]^N\). If \(H\ge1\) bounds the integer entries after clearing
each possible row, a positive optimum is at least \(1/Q\) with
\(Q=((N+1)H)^{N+1}\). The individual rows, \(K\) and \(Q\) can be
computed in polynomial bit time before the support is known. Query
\(\delta=1/(2KQ)\) and take \(J\) to be the positive coordinates of
the returned solution. It satisfies the original rows at \(t=\delta K\),
so the optimum is zero. Any exact optimizer now has a zero variable or
zero slack in every pair and gives an equilibrium by the displayed
system. This proves the direct upper reduction of \cite[Theorem~5.19]{book}.

\section{The fixed point and the reverse matrix family}
\reviewmark{Earlier contraction theorem and the promise in the reverse direction.}

Theorems~5.27--5.28 of \cite{book} give a piecewise-linear circuit on the
\((n+1)\)-dimensional unit cube and an explicit rational \(\delta>0\) with
\[
    \|F(s)-F(t)\|_1\le(1-\delta)\|s-t\|_1 .
\]
The exact fixed point
determines the attack support and the zero, interior and full coverage
labels for an original-game LP. Exact feasibility on that support
recovers a BNE even with zero gains or slopes. In the total-search
formulation of PL-Contraction, a violation of the promised
contraction can also be an output. Here the constructed circuit
satisfies the inequality everywhere, so every valid answer is its fixed
point. This and the P-LCP construction provide two routes into
\(\mathsf{UEOPL}\)~\cite{book,fgms}.

In the other direction Appendix~\ref{sec:cayley} starts with a rational
P-matrix \(A\) and a supplied rational \(\lambda>0\) for which
\[
 \rho\!\left(\left|(I+\lambda A)^{-1}(I-\lambda A)\right|\right)<1 .
\]
The absolute value is entrywise. In that family the transformed
absolute-value equations can be encoded by market purchases. The
spectral-radius condition supplies the error bound used to recover the
exact complementary signs. Lemma~\ref{lem:cayley-qp} also gives a
direct polynomial-time algorithm for this promised family and exhibits
a P-matrix excluded for every \(\lambda>0\).

\chapter{Reductions to stopping stochastic games}\label{sec:ssg}
A simple stochastic game has Max, Min and fair-coin vertices and two
sinks paying zero and one. The value is the probability of reaching the
one sink under optimal play~\cite{condon}. Explicit SSGs also reduce to
P-LCP~\cite[\S1]{fgms}. Here we construct the stochastic game directly.
We use stopping games, in which
absorption occurs almost surely for every strategy pair. For an
\emph{explicit} game the input lists its vertices and successors. For a
\emph{succinct} game Boolean circuits supply them from a binary state name.

\section{The comparison game for two targets}
\reviewmark{Threshold gaps, ties, endpoint play and fair-coin sampling.}
\begin{theorem}\label{thm:ssg}
With two targets and arbitrarily many attacker types, the exact start
value of one polynomial-size explicit stopping game determines a full BNE.
For an arbitrary market with \(N\) buyer and good vertices, all
singleton values in one succinct stopping game determine an exact market
equilibrium. These designated values obey a denominator bound of polynomial
binary length computable from the market. They can also be recovered using
polynomially many adaptive yes/no comparisons with rational thresholds.
The explicit expansion has \(2^N\operatorname{poly}(L)\) states. Composing
with \(\BNE\dd\Market\) gives both reductions for unrestricted BNE.
\end{theorem}

\begin{proof}[Two targets]
Budgets zero and two fix the defense. With one resource write its coverage
as \((x,1-x)\). If \(d_1+d_2=0\), choose any coverage and best replies.
Otherwise put \(h=d_2/(d_1+d_2)\) and
\[
  \Delta_i(t)=a_{i1}-a_{i2}+b_{i2}-(b_{i1}+b_{i2})t,\quad
  L(t)=\sum_{\Delta_i(t)>0}p_i,\quad U(t)=\sum_{\Delta_i(t)\ge0}p_i .
\]
Attainable attack mass on target one is exactly \([L(t),U(t)]\).
The defender is willing to mix when that mass equals \(h\), to cover
only target two when it is at most \(h\), and to cover only target one
when it is at least \(h\).

Let \(C\) consist of \(0,1\) and all distinct interior roots of the
nonconstant \(\Delta_i\). Let \(Q\) be the product of the denominators of
the priors and \(h\), and put \(c(t)=1/2-t/(4Q)\). Max chooses \(t\in C\).
At \(t=1\) pay a Bernoulli reward of mean \(c(1)\). At \(t<1\) Min chooses
between a reward of mean \(c(t)\) and a test with mean
\((1-L(t)+h)/2\). The test tosses a fair coin between sampling a type and
accepting if \(\Delta_i(t)\le0\), and accepting with probability \(h\).
If \(L(t)\le h\), this node has value \(c(t)\). Otherwise \(L(t)-h\ge1/Q\)
and its value is strictly below \(c(1)\). Hence the start value \(V\) yields
\(t^*=4Q(1/2-V)\), the first passing candidate or the fallback \(1\).
If \(t^*>0\), the previous candidate \(s\) failed. Every strict target-one
type there still weakly prefers it at \(t^*\), so \(U(t^*)\ge L(s)>h\).
At an interior solution split tied types to reach \(h\) exactly. At an
endpoint send all ties to the favored endpoint's target (target two for
\(t^*=0\), target one for \(t^*=1\)). Zero-prior rows take any best reply.

For a rational \(k/d\), sample \(\lceil\log_2d\rceil\) fair bits, accept
below \(k\), reject in \([k,d)\), and repeat otherwise. Tracking two prefix
comparisons uses \(O(1+\log d)\) states and the trial terminates almost
surely. Sequential conditional Bernoulli choices implement any listed
rational distribution. Binary choice trees implement Max's finite menu.
All sizes and rational encodings above are polynomial.
\end{proof}
On two targets the decoder can also compute \(t^*\) directly. Sort
the polynomially many rational candidates, evaluate the passing
inequality at each, and apply the same tie split. The stochastic
game gives an explicit value representation of this threshold.

\section{Further explicit games from payoff data}
\reviewmark{One-type coverage construction and the explicit circuit encoding.}

With one attacker type, the two-target restriction can be removed.
The reduction below works from the payoff table and gives a single
explicit start value. It also explains why a fixed bound on both table
dimensions has an elementary explicit representation.

\begin{theorem}\label{thm:ssgsmall}
For a one-type security game with arbitrarily many targets, or a game
with fixed numbers of types and targets, one explicit stopping SSG of
polynomial size has a start value encoding a full BNE. Endpoint budgets
also admit this construction with arbitrary table dimensions.
\end{theorem}
\begin{proof}
At an endpoint budget the defense is fixed and a maximum-payoff
comparison in each row determines a best reply. For one type and
\(0<r<n\), let \(u_0=\max_j(a_j-b_j)\). To make the largest attack
payoff at most \(u\), it is necessary and sufficient that
\[
 u\ge u_0,\qquad
 \sum_{j:b_j>0}\max(0,(a_j-u)/b_j)\le r .
\]
The summands are the least required coverages and \(u\ge u_0\) keeps
each at most one and handles zero slopes. Remaining budget can always
be placed up to the unit caps. Order the positive-slope targets by
decreasing \(a_j\). The least feasible \(u\) is the maximum of \(u_0\)
and, over all nonempty prefixes \(S\), the numbers
\[
 \frac{\sum_{j\in S} a_j/b_j-r}{\sum_{j\in S}1/b_j}.
\]
Each prefix gives a necessary inequality. At a proposed \(u\) the
indices with positive summands are themselves a prefix, proving
sufficiency. If there are no positive slopes take \(u=u_0\).

Give every positive-slope target its least required coverage. Place
any surplus under the caps, giving priority to flat targets with
\(a_j=u_0\). If a best target has zero defender gain or is fully covered,
attack it purely. In either case the chosen defense is optimal.
Otherwise every positively covered target is best. With no surplus
this follows from the least-coverage formula. With surplus, \(u=u_0\):
if \(u>u_0\), a small decrease would remain feasible. A positive-slope
target attaining \(a_j-b_j=u_0\) would already be full and best.
Thus there is a flat maximizer, and the priority rule either makes one
full or uses all surplus on such maximizers. Attack the positively
covered targets in proportion to \(1/d_j\). Their scores are equal
and all other scores are zero, verifying defender optimality.

For fixed numbers of types and targets, enumerate the finitely many
attack supports and target labels. The row optimality, mass, budget and
cutoff conditions on a fixed pattern are linear. A feasible closed
pattern suffices, with zero mass permitted on an edge required to be
tight. A fixed-dimensional rational LP yields a profile. Use a fixed
rule for breaking ties in each construction. Realize any returned
coverage by the polynomial-size lottery construction used in
Chapter~\ref{sec:unique}.

To turn these particular payoff procedures into explicit games, express
their binary computation as a polynomial-size Boolean circuit. Sorting,
rational arithmetic, and the support and lottery operations above have
polynomial bit cost. Pad the selected numerator, denominator and support
fields to a known polynomial length \(k\). Use a wire and its complement
for each bit, with input literals as absorbing zero and one. An AND is a
Min vertex, an OR is a Max vertex, and negation exchanges the two wires.
All gate values are Boolean by induction. If \(b_1,\ldots,b_k\) are the
output bits, put \(V_{k+1}=0\) and make \(V_j\) a fair-coin vertex with
successors \(b_j,V_{j+1}\). Then \(2^kV_1\) is the desired bit string as
a padded integer. The graph is acyclic before its sinks and hence
stopping under every pair of strategies.
\end{proof}

\section{The representation for arbitrary markets}
\reviewmark{Output length of singleton values, threshold queries and the subset-state bound.}
Appendix~\ref{sec:subset} proves the unrestricted statement and gives
the conversion to two fair-coin sinks. Output size matters in the succinct
model. For \(k\ge1\), index \(2^k\) fair-coin states by a \(k\)-bit
counter. Tails loses immediately. Heads advances the counter or wins if
it is at its maximum. Starting at zero the win probability is
\(2^{-2^k}\), whose reduced denominator takes \(2^k+1\) binary digits.
The successor circuit has size polynomial in \(k\) and the game stops
on every play. For our constructed instances the denominator bound in
Appendix~\ref{sec:subset} supplies the output promise needed for an
exact-value query. The threshold-decision version applies to the ordinary
succinct stopping-game threshold problem. On a star the construction reaches
exponentially many subset states even from the central singleton. A
polynomial-size explicit stochastic-game representation for arbitrarily
many targets remains a separate question for these reductions.

\chapter{Cost of using the reductions}\label{sec:scope}
\reviewmark{Oracle costs and the scope of the reported experiments.}

Write \(T(s)\) for a nondecreasing bound on the cost of returning one
polynomial-size rational equilibrium on a game of binary length \(s\).
For matching, parallel-road traffic and ordinary ARRIVAL, the constructions
give a bound \(T(p(L))+q(L)\), where the polynomials \(p,q\) depend on the
reduction. The same bound applies to the hide-and-seek family of
Theorem~\ref{thm:hide}. With \(d\) terminals and total supply \(S\),
access to ordinary ARRIVAL decisions suffices in at most
\(d\lceil\log_2(S+1)\rceil\) independent queries. With a game oracle,
Corollary~\ref{cor:multi-one} builds the binary-supply market directly and
gives the entire vector in \(T(p'(L))+q'(L)\) for polynomials \(p',q'\).
Full uniqueness, including a second profile when it
exists, costs \(T(L)+\operatorname{poly}(L)\) under the stated output
convention. The reductions alone give these bounds relative to an oracle.


The pinned split, fixed-total rent division, and the one-unit update each
use one equilibrium call and polynomial additional bit work. The rent
decoder includes a bipartite matching. Recovering a BNE from uniqueness
uses independent decision queries. The reverse quota reduction bisects
adaptively. The cycle-only reduction asks independently for each price.

Chapter~\ref{sec:upper} queries a genuine P-matrix. The reverse decoder
requires the supplied Cayley parameter and spectral bound of
Appendix~\ref{sec:cayley}. Chapter~\ref{sec:ssg} gives polynomial-size
explicit stochastic games for its restricted families. General BNE uses
a succinct game and one polynomial-bit singleton value per market
vertex, or polynomially many adaptive threshold queries.

A parallel road specifies a whole route. Multiple junctions would need
flow conservation for each class. Binary dispatches and terminal counts
are charged by bit length, so exponentially many chips are permitted.

\paragraph{Computational checks and provenance.}
Vincent Lin coordinated the project's AI agents. The dated notes record
the drafting AI's self-review. The compiler examples test specific
adaptive queries. A separate exact check uses both control-bit settings
at three module scales and quantifies over all nonnegative purchases at
two ports whose sum is at most a quarter of the scale.
Appendix~\ref{sec:compile} gives the general argument. Exact
rational checks of the price construction considered all real external loads
at depths two and three and checked 28 interval or endpoint cases. A supplied
allocation in a 358-good, 184-buyer market at \(\delta=2^{-43}\) gave the
support of the two-road example above. A separate checker compared 2,528
single-chip bit answers with direct routing on 334 Multi-Arrival instances.
The uniqueness tests were compared with unrestricted equilibrium search at
48 supplied equilibria on 27 game instances, and with explicit pure-defense
enumeration on 15 full profiles. Exact rational checks also recovered a
two-target equilibrium in 3,804 cases, including simultaneous ties and
zero effects. On 96 small markets we compared the subset Bellman update
with the solution of the perturbed market equations on all 3,360 subset
states before expansion into fair-coin samplers.

For the new graph formulation, exact checks included 24 bounded-price
best-response systems and the reverse decoder with signed prices. The
unit-update checks included outside-option ties and long binary margins.
For rent division, 30 four-tenant tests checked the total rent, all
assigned best responses and the support matching. On attack supports,
208 proposals yielded 416 comparisons with an unrestricted feasibility
test, and 33 games supplied 198 checks of the flat-offer decisions with
and without a specified cutoff.

\appendix
\chapter{Encoding adaptive queries in a single market}\label{sec:compile}
\reviewmark{Isolation, margin separation, rational perturbations and support repair. The error bound for reading the two signs of every offer slack was checked again on November 27, 2026 (UTC).}

The proof of Theorem~\ref{thm:compile} uses three steps. First, perturb a
query so that it has a unique answer and its set of tight offers determines
an answer to the original query. Next, place the query inside a larger market
whose earlier goods select its data. Its prices need only be accurate enough
to distinguish a zero slack from a positive one. Finally, encode these
comparisons as exact bits and use them to specify the next query. The
capacity estimates below account for the purchases made by each later
reader of an earlier good.

\begin{lemma}[Isolation of a market]\label{lem:isolate}
An ordinary market reduces in one call to a market on the same graph
with unchanged capacities and slopes and a unique real equilibrium.
Its tight graph is a forest. Exactly one vertex of each tree has zero
utility or price. A rational LP recovers an original equilibrium.
\end{lemma}
\begin{proof}
Choose a positive common denominator \(Q\) for the input data. Number
the edges \(1,\ldots,|E|\). Add \(\delta\) to each buyer mass and
\(\delta^e\) to intercept \(e\), for a small dyadic \(\delta>0\).
For a nonempty buyer set \(D\) and any good set \(S\), the difference
of their old total margins is in \(Q^{-1}\mathbb Z\). The change
\(|D|\delta\) lies in \((0,1/Q)\), so their new totals differ.

Suppose two equilibria have price vectors \(p,q\), and let
\(S=\{j:p_j>q_j\}\). If \(S\) is nonempty, it fills at \(p\). Every
buyer who buys in \(S\) at \(p\) strictly prefers \(S\) to all other
options at \(q\), since its used offer improves and every outside offer
weakly worsens. These buyers must therefore have total mass exactly
the capacity of \(S\), a contradiction. Switching \(p,q\) proves unique
prices. The best offers then fix utilities.

Consider all tight offer equations and all zero-coordinate equations at
these potentials. Their normals are independent. Otherwise a minimal
dependence includes an offer, since the zero-coordinate normals alone are
independent. After clearing rows, its integer coefficients can be taken as
signed minors. Consistency of the right-hand sides gives a polynomial in
\(\delta\) whose distinct offer powers have nonzero coefficients. A
sufficiently small positive dyadic cannot be its root. Every connected
component of tight edges has a zero coordinate, since otherwise all its
positive margins would fill internally. A component on \(v\) vertices
has at least \(v-1\) edges and at most \(v\) independent equations
including the zero coordinates. It is a tree with one zero coordinate.
Every other margin fills, so deleting leaves towards that vertex
determines the allocation uniquely.

For recovery use the old inequalities. For each positive perturbed
coordinate impose equality in its complementary inequality, and set
each zero coordinate to zero. Every feasible solution is an original
equilibrium. To see feasibility, clear all rows by \(Q\), keep
nonnegativity, relax each row by a common slack \(t\ge0\), and minimize
\(t\). A positive optimum is at least \(1/B\) by the vertex determinant
bound, whereas the perturbed solution is feasible with \(t=Q\delta\).
For explicit sizes let \(s=|E|+|I|+|J|+1\), let \(H\ge5\) bound \(Q\)
and all absolute cleared entries, and take \(B=(sH)^{s+2}\). This bounds
the minors above as well. After clearing the tight equations the coefficients
of their consistency polynomial are integers of magnitude at most \(sHB\).
Following its least nonzero power, the tail is smaller in magnitude than
its nonzero integer coefficient if \(sHB\,\delta/(1-\delta)<1\).
Choosing \(\delta<1/\max(8,2(|I|+1)Q,2QB,4sHB)\) meets all three
requirements and gives all powers polynomial binary length.
\end{proof}

\section{Computing with the prices of ordinary goods}
\reviewmark{Generator inequalities, clipping, rectangle slopes and source-flow recovery.}

Represent a signed number by a good's price minus a known offset. A gate
will enforce a minimum or maximum of a finite list of increasing affine
functions of earlier prices. Choose the output offset so each physical
threshold \(T_l(p_l)=A_l+B_l p_l\), \(B_l>0\), is at least one on
nonnegative prices. Choose \(U\) at least one above every threshold on
the input price ranges of interest. Make the output good capacity \(K\)
and give a unit-mass guard the sole offer \(U+1-p\) there. For a minimum,
one buyer of mass \(K\) has output offer \(U-p\) and input offers
\(U-T_l(p_l)\) for all \(l\). If two pins of a minimum refer to the
same good, insert distinct unary identity copies before the gate. Each
listed buyer--good pair then has just one offer. These copies count in
the reverse capacity calculation below. For a maximum use one such buyer per input,
each seeing only that input and the output. All buyers have their usual
outside option. Choose capacities in reverse circuit order so all readers
of this output together have mass \(<K-1\), allowing for external
readers as well. For binary gates, taking \(K\) greater by two than the
total mass already allocated suffices, with at most constant-factor
growth per gate. Reserve extra capacity for stated external reader bounds.

At price zero the guard and generators strictly want the output and exceed
its capacity. At a positive price it fills. Readers and the guard supply
strictly less than \(K\), so a generator buys there. This gives \(p\le U\)
and makes the guard buy its full unit. The aggregate generator purchase
is therefore strictly between zero and \(K\).
For a generator \(g\) with input set \(D_g\), put
\(c_g=\min(U,\min_{l\in D_g}T_l(p_l))\). Its output offer is \(U-p\).
If \(p<c_g\), that offer strictly beats all alternatives, including the
outside option, so \(g\) sends all \(K\) units to the output. If \(p>c_g\), it sends zero. In particular, positive purchases imply \(p\le c_g\).
The aggregate bounds rule out \(p<\max_g c_g\) and \(p>\max_g c_g\),
respectively. Thus \(p=\max_g c_g\). One generator with all inputs gives
the clipped minimum. One generator per input gives the clipped maximum.
On the designated range no clipping occurs. A constant \(c>0\) is fixed
by a good of capacity \(K\) with one buyer of mass \(2K\) offering
\(c-p\), subject to total reader mass \(<K\). Translation handles signed
constants. Exact bits and their complements use separate goods. AND and
OR are minima and maxima, and negation exchanges the two goods.

\begin{lemma}[Two slope profiles]\label{lem:rectangle}
Any market reduces many-one to a market with slopes one except on one
buyer-group by good-group rectangle of twos.
\end{lemma}
\begin{proof}
Scale each source buyer's offers by a positive power of two so all slopes
\(\beta\) are at least one. Choose \(H>2+\max(0,\max_e a_e/\beta_e)\).
For each pair make a private circuit with only one external input, the
terminal price \(P_j\). On \([0,H]\) it computes a function \(s_e(P_j)\)
within \(\xi\) of \(\beta_eP_j\). Its gates have slope one except a unary
doubling gate. Doubling uses a group-two buyer on an exceptional input
good and an ordinary output good. A unit-slope identity can produce
an exceptional input before each doubling. All other generators, guards
and source buyers are in group one. Negative intercepts on unused
pairs realize the required slope pattern across the entire table.

Here is a circuit of polynomial size for prescribed binary precision.
Bracket \(\beta\) by successive powers of two and compute both multiples
of \(P_j\) by doubling. Interpolate between them using the binary digits
of the fractional position of \(\beta\), truncated to precision
\(\xi/(4H\,2^{\lceil\log_2\beta\rceil})\). Read the digits least
significant first, starting at the lower multiple and averaging with the
lower or upper multiple according to the bit. To average \(v,w\) with
\(|v-w|\le D\), replace them by
\(\max(v,w)-D/4,\min(v,w)+D/4\). This preserves their sum and reduces the
gap bound to \(D/2\). Share both outputs at the next fold. A polynomial
number of folds approximates the mean within \(\xi/8\) on the full range.
The error of each incoming approximate mean is halved at the next
average. Taking \(D\) equal to the original range plus \(\xi\) covers
these errors throughout. The result has the required error and is
globally nondecreasing in the terminal price even when physical gates clip.

Source buyer \(i\) buys only at its private final goods, with effective
utility \(a_e-s_e(P_j)\). Include a final identity gate on each edge.
Let \(D_e\) be the sum of private buyer masses in the circuit minus the
sum of its internal capacities. Each gate contributes
\(1+(\#\text{generators}-1)K\), which is positive. Give terminal \(j\)
capacity \(c_j+\sum_{e\to j}D_e\). If \(x_e\) is its source buyer's
purchase at the private root, cancellation across the filled internal
goods bounds the actual load delivered by the circuit to \(j\) by
\(D_e+x_e\). For \(\xi<1/4\), at \(P_j\ge H\) every source offer ending
at \(j\) is negative, by monotonicity and its value at \(H\). At such
a price \(j\) cannot fill. Thus all terminal prices are below \(H\).
All private buyers now spend in full and delivered load equals \(D_e+x_e\).
Consequently the source row and column margins already clear exactly.

Keep those flows, their zero labels, and the observed slack margins.
On the original offers solve for nonnegative prices and utilities,
imposing offer inequalities, tightness on used pairs, and zero potential
on each slack margin. Minimize a common relaxation \(t\) of the payoff
rows, bounding prices by \(H\) and utilities by \(2+\max(0,\max_e a_e)\).
At the observed prices and utilities, \(t=\xi\) is feasible. The only
information from the flows entering this LP is a set of row labels.
The finite polynomial menu of integer-cleared rows has a uniform
determinant bound \(B\) of polynomial bit length, as in
Lemma~\ref{lem:repair}. Choose \(\xi<1/(2B)\) before building the circuits.
The minimum is then zero. This gives an exact input equilibrium with
the same source flows. Powers of two, offsets, circuit capacities and
the inverse accuracy all have polynomial bit length. For offsets, products
along a path involve at most one doubling per gate and all shifts have
common denominator of polynomial bit length.
\end{proof}

\section{A market with coefficients selected by input bits}
\reviewmark{Disabled buyers, baseline loads, precision separation and module scaling.}

\begin{lemma}\label{lem:program}
Given polynomial bounds on dimensions and coefficient bit lengths, a fixed
ordinary market module can represent any isolated query of that size. Its
data are selected by exact bit goods in earlier modules. For a prescribed
dyadic \(\zeta>0\), its output prices approximate the unique query prices
within \(\zeta\) whenever the total mass of later buyers is below a
computable positive tolerance. The module, including the binary length
of this tolerance, has size polynomial in the bounds and \(\log(1/\zeta)\).
\end{lemma}
\begin{proof}
Pad the abstract query with positive-margin isolated vertices and encode
missing edges by negative offers. Use a common positive power of two
to scale offers so all possible slopes are at least one. This leaves
good prices invariant. Choose \(H\) uniformly above all possible positive
ratios \(a_e/\beta_e\) by at least two. Build a private unary chain for
each possible pair. Its logical output approximates
\(\beta_eP_j-a_e\) within \(\xi<1/4\) on \([0,H]\). A binary choice
between two affine laws uses two equal-mass generators in place of one.
Each has an extra offer on its own earlier bit good, strictly better than
all local options in the disabling state and below zero in the enabling
state. The complementary bits enable exactly one generator. Such an
offer is decreasing affine in the physical bit price, using its known
offset and the finite upper bounds on the local intercepts. The disabled
buyer's entire mass goes to the earlier good. The active generator and
guard enforce exactly the clipped unary gate law just proved.

For each source copy, offer on its final edge good its offset minus that
good's physical price. The effective source payoff then approximates
\(a_e-\beta_eP_j\). A guard or enabled generator buys only in its own
chain or on its terminal. Disabled buyers are accounted for solely as
readers of earlier bit goods.

To encode a slope \(\beta\ge1\), first program the exponent of
\(2^{\lfloor\log_2\beta\rfloor}\) by its binary digits, selecting a
multiplier \(1\) or \(2^{2^k}\) at position \(k\). To approximate the
remaining factor in \([1,2)\) from below, at step \(k=1,2,\ldots\)
use the largest of zero, one or two factors \(1+2^{-k}\) that keep
the product below or equal to the requested factor. The remaining ratio
is at most \(1+2^{-k}\), since
\((1+2^{-k})^2\ge1+2^{-(k-1)}\). Stop once
\((\max\beta)H2^{-k}<\xi/2\). Translate by dyadic digits for \(-a_e\)
with error \(<\xi/2\). Use program-independent positive offsets and gate
ceilings computed forward for both possible branches. A fixed, padded
chain thereby covers every possible query word. The bit choices follow
from rational comparisons of polynomial-length numbers.

Represent masses by enabled binary-weight copies of each source buyer,
with the same offers on its roots, and approximate each margin to
accuracy \(\xi\) on a common dyadic grid. Take \(\xi\) below half every
possible positive margin. At terminal \(j\) use binary-weight singleton
buyers offering \(H+1-P_j\) to preload the difference
\(C-\widehat c_j\) between a fixed ceiling \(C\) and the programmed
capacity. Disabled copies and singleton buyers go to the earlier
controls. Give the terminal capacity \(C+D_j\), where \(D_j\) is the
sum of the lengths of all private unary chains attached to \(j\).
At each gate its active generator and guard have total mass \(K+1\).
If \(s\) units from readers purchase its filled good of capacity \(K\),
the two private buyers spend \(K-s\) there. They send at most \(1+s\)
to the preceding good, with equality when both spend in full.
Telescoping this inequality along every chain,
the total load on \(j\) is at most the source purchases into its roots
plus \(D_j+C-\widehat c_j+r\), where \(r\) bounds all later reader
mass, including purchases inside the chains. If \(P_j\ge H\), the source
purchases are zero, by monotonicity for the fixed program and accuracy
at \(H\). Choosing \(r<\min_j\widehat c_j\) contradicts filling a
positive-price terminal. Below \(H\) all active private and preload rows
spend in full and the same accounting is an equality, using the actual
external purchases at \(j\) and on its chains. The effective column
margins are consequently within \(\xi+r\) of the desired values, with
the same bound on the absolute fill error at positive prices.

To choose precision uniformly, denote the intended isolated prices by
\(p^*\). To each closed equilibrium pattern append one inequality
\(p_j\ge p_j^*+\zeta\) or \(p_j\le p_j^*-\zeta\). Every such system
is infeasible by uniqueness. Here is one choice of precision computable
before the query bits are known. Let \(d\) bound the number of variables
in a padded query pattern. Let \(B\ge1\) bound the bit lengths of the
numerators and denominators of each row entry, including its right-hand
side, after the uniform offer scaling. Put \(H_0=2^{(d+2)B}\) and
\(D_0=(dH_0)^d\). Multiplication by the product of a row's denominators
makes its entries integers bounded by \(H_0\). The unique equilibrium
is a vertex of any nonempty exact pattern. Cramer's rule bounds both
the reduced numerator and denominator of every \(p_j^*\) by \(D_0\).
Take \(\zeta=2^{-k}\), rounding a requested accuracy down to a dyadic
power if necessary, with \(k\ge0\). The two bad inequalities now have
rational entries with numerators and denominators bounded by
\(2^{k+1}(D_0+1)\). Choose
\[
 B'=B+k+3+\lceil\log_2(D_0+1)\rceil,\qquad
 H_1=2^{(d+3)B'},\qquad
 \kappa=\bigl(2((d+1)H_1)^{d+1}\bigr)^{-1}.
\]
Relax the pattern and the bad inequality by the same nonnegative slack
\(t\) in the original row units, keeping the variables nonnegative.
The enlarged rows, including the coefficient of \(t\), clear to integers
of magnitude at most \(H_1\). The LP minimizing \(t\) is feasible
at large \(t\), is pointed, and attains its optimum at a vertex.
Infeasibility at zero makes the optimum at least \(2\kappa\).
These calculations use only dimension and entry bounds. They call for
neither the unknown \(p^*\) nor a common denominator over all possible
program words.
The actual module prices, aggregate source flows and true best utilities
satisfy a closed query pattern to error at most \(2\xi+r\). Use the
full-row label when all the active copies spend, and the zero-utility
label otherwise. Used offers miss the true best utility by at most
\(2\xi\). If some terminal price differed from \(p^*\) by at least
\(\zeta\), the associated bad inequality would also hold. Choose
\(2\xi+r<\kappa\) to rule this out. Gate capacities grow by a constant
factor per gate in reverse order. Gate offsets and ceilings use only
sums and products along polynomial-length paths of predetermined
rational factors, with polynomially many primitive constants of
polynomial bit length. This gives the asserted binary size bound.
\end{proof}

\begin{theorem}[Adaptive queries]\label{thm:compile}
For a deterministic polynomial-time algorithm with adaptive \(\Market\)
search queries, one ordinary market encodes a valid transcript. From
every one of its equilibria a polynomial-time decoder obtains the
algorithm's result. In particular \(\BNE\dd\Market\).
\end{theorem}
\begin{proof}
Unroll the computation up to its polynomial time bound, padding unused
rounds with a fixed query. Fix the bound for the original machine on
polynomial-length valid oracle answers. Its time bound bounds both the
number of rounds and every submitted query word. Isolation, transposition,
LP decoding and the two sign readouts each have a polynomial size bound
in that word bound. Compute these uniform bounds before constructing any
physical module. A later round receives answer bits for the original
oracle problem, so its abstract input size is still bounded by the
original machine clock. At each round a Boolean circuit computes and
isolates the actual query using the bits returned in earlier rounds.
Make two programmed modules for that isolated market \(H\) and its
transpose \(H^\top\): exchange buyer and good margins and replace
\((a,b)\) on each reversed edge by \((a/b,1/b)\). Dividing its original
offer inequality by \(b\) shows that prices, utilities and flows
transform by \((P,u,g)\mapsto(u,P,g^{\mathsf T})\). Both markets have
unique solutions and the two sets of good prices approximate \(P\) and \(u\).
The following identities give their negatives.
\[
 -u_i=\min(0,\min_j(b_{ij}P_j-a_{ij})),\qquad
 -P_j=\min(0,\min_i (u_i-a_{ij})/b_{ij}).
\]
Indeed positive potentials force a used tight edge and a zero potential
attains the constant branch. An empty inner menu contributes no term.
Here \(a,b,u,P\) describe the isolated query before internal row scaling.
Such scaling preserves the goods' prices in each of the two modules.

The quantities \(u_i,P_j\) and the offer slacks
\(u_i+b_{ij}P_j-a_{ij}\) are nonnegative. Any positive one exceeds a
uniform dyadic \(\Delta>0\) of polynomial bit length by the same basis
bound as above. Let \(C\ge1\) uniformly bound every \(b_{ij}\) and \(1/b_{ij}\).
Choose a dyadic \(e<\Delta/(16(C+1)^2)\) as error bound for each price
module, each programmed affine map, and each sum. A sum uses twice the
approximate mean from Lemma~\ref{lem:rectangle}, computed to error \(e/2\).
Program each unary map from the query bits as in Lemma~\ref{lem:program}.
Dividing the input by a fixed power of two at least \(C\) allows
programmed slopes below one after multiplying the programmed factor by
that power. For a signed input, a known shift and the corresponding
intercept correction bring the argument into the nonnegative range.
The two minimum identities then give \(-u_i,-P_j\) within \((C+1)e\).
Sum \(u_i\) with \(b_{ij}P_j-a_{ij}\) for the slack and \(-u_i\) with
\(b_{ij}(-P_j)+a_{ij}\) for its negative. The larger error is at most
\(((C+1)^2+2)e<\Delta/8\). All precisions have polynomial binary length.
Given approximations
\(s^+,s^-\) of \(w,-w\), clip \(2s^+/\Delta-1/2\) and
\(3/2+2s^-/\Delta\) to \([0,1]\). The results are the exact complementary
bits for positivity and zero, also at ties.

The price, utility and tight-edge labels determine the unique isolated
profile by LP: set untight trades to zero, keep the labeled equalities,
and fill each positive-potential margin. Every feasible point is an
equilibrium. A second LP from Lemma~\ref{lem:isolate} recovers an
equilibrium of the submitted query. Choose a basic feasible answer
in the source support polytope, using a fixed deterministic rational LP
algorithm (successively minimize its coordinates if necessary). Its
encoding obeys the source oracle's witness bound. These computations can be unrolled
as Boolean circuits that form the next query word. This selects one valid
transcript for a search algorithm accepting arbitrary valid answers.

Finally, compute each module's intrinsic precision before scaling it.
Let \(M\ge1\) bound the normalized total mass of any module and \(s\)
be the number of modules. In each readout gate reserve at least one
unit for external readers, and let \(r>0\) be at most this reserve and
the least tolerance of any program module. Treat the two programmed
markets of a round as separate modules and place the readout gates
strictly after both. Purchases by these readout buyers count toward the reader tolerance of
each program module. Order the modules in this dependency order.
Multiply both masses and capacities of module \(t\) by \(\alpha^t\),
where \(0<\alpha<\min(1/2,r/(2sM))\) is dyadic. In an equilibrium
of the combined market, divide all purchases by buyers of module \(t\)
and its capacities by \(\alpha^t\). Its buyer comparisons still use the
same prices, and its local clearing equations now have their original
margins with incoming reader purchases measured in the same units. All
cross-module purchases are by later buyers on earlier goods. The normalized
incoming reader mass of any module is bounded by
\(\sum_{j>t}M\alpha^{j-t}\le sM\alpha<r\), for every choice made by
those buyers. The tolerance bound in Lemma~\ref{lem:program} therefore
applies to its terminal prices. The gate bounds give the intended exact
bits at every stage by induction. The scaled data retain
polynomial bit length. The one-resource game embedding of the combined ordinary market has a
mixed Nash equilibrium. The forward decoder of
\cite[Theorem~3.16]{book} produces an equilibrium of the combined market. Apply the construction to
the adaptive algorithm of~\cite[Theorem~3.1]{book} for the final claim.
\end{proof}

\chapter{A stopping game on subsets of market vertices}\label{sec:subset}
\reviewmark{Bellman equations, stopping probability, sampler size and exact readout.}

Use the diagonal perturbation of the market LCP from
\cite[Theorem~5.6]{book}. It gives \(0<\varepsilon\le1\) of polynomial
bit length. Put \(z_i=-u_i\) for buyers and \(z_j=P_j\) for goods. The
perturbed equations for an edge \(e=ij\) of slope \(b_e>0\) are
\[
 f_e=\varepsilon^{-1}[a_e+z_i-b_ez_j]_+,\qquad
 \begin{cases}
 0\le m_i-\varepsilon z_i-\sum_j f_{ij}\ \perp\ -z_i\ge0,\\
 0\le c_j+\varepsilon z_j-\sum_i f_{ij}\ \perp\ z_j\ge0 .
 \end{cases}
\]
Its solution is unique, and its positive coordinates specify a feasible
LP on the unperturbed data.

Let \(R_e\) send \([a_e+z_i-b_ez_j]_+/(1+b_e)\) from coordinate \(i\)
to coordinate \(j\), subtracting from the former and adding to the latter.
At a subset \(S\) write \(Z_S=\sum_{v\in S}z_v\). The sum stays fixed if
both endpoints belong to \(S\) or neither does. If only the buyer belongs,
its new sum is the first of the following expressions. If only the good
belongs it is the second:
\begin{align*}
 &\min\!\left(Z_S,\frac{b_e Z_{S\cup\{j\}}+Z_{S\setminus\{i\}}-a_e}{1+b_e}\right),\\
 &\max\!\left(Z_S,\frac{Z_{S\cup\{i\}}+b_e Z_{S\setminus\{j\}}+a_e}{1+b_e}\right).
\end{align*}
Subtracting \(Z_S\) from each second choice gives respectively the negative
and the positive of \((a_e+z_i-b_ez_j)/(1+b_e)\). In each row the two
continuation weights sum to one. These formulas thus define Min and Max
options on arbitrary tables, each 1-Lipschitz in the sup norm.

Take
\[
 \tau=\frac1{2(\varepsilon^2+\sum_e(1+b_e))},\qquad
 \theta=\tau\varepsilon^2,\quad w_e=\tau(1+b_e),\quad w_0=\tfrac12 .
\]
Let \(d_i=m_i/\varepsilon\), \(d_j=-c_j/\varepsilon\), and set
\(Y(z)=\theta d+w_0z+\sum_e w_eR_e(z)\). At a buyer, \(Y_i-z_i\) equals
\(\tau\varepsilon(m_i-\varepsilon z_i-\sum_j f_{ij})\). At a good,
\(Y_j-z_j=\tau\varepsilon(\sum_i f_{ij}-c_j-\varepsilon z_j)\).
Clip buyers to \(\min(0,Y_i)\) and goods to \(\max(0,Y_j)\). The fixed
point equations are exactly the displayed complementarity conditions.
For a subset sum, the affine constant is \(\theta\sum_{v\in S}d_v\).
To clip a coordinate \(v\in S\), a player compares the current entries
at \(S\) and \(S\setminus\{v\}\), taking their minimum at a buyer and
maximum at a good. Apply one coordinate per stage, leaving entries with
\(v\notin S\) intact. On additive tables these stages give exactly the
subset sums of the clipped vector. On arbitrary tables they are
1-Lipschitz. The raw update has Lipschitz constant at most
\(w_0+\sum_e w_e=1-\theta\). The resulting Bellman operator has a unique
fixed table. The subset sums of the perturbed solution supply such a
table, so its singleton entries recover that solution.

To realize the offsets with terminal probabilities put
\(\rho=1-\theta/2\) and \(\gamma=1-(1-\theta)/\rho>0\). Follow the clip
choices in reverse order of their application. For the resulting subset
select edge \(e\) with probability \(w_e/\rho\), the identity with
probability \(w_0/\rho\), and otherwise stop with reward
\(\theta\sum_{v\in S}d_v/\gamma\). At the selected comparison, each
action with continuation distribution \(D\) and offset \(k\) continues
according to \(D\) with probability \(\rho\), and otherwise stops with
reward \(\rho k/(1-\rho)\). The identity has \(k=0\). This gives the
intended value \(\rho(\mathbb E_D Z+k)\) for that action. The probability
of reaching another round is exactly \(1-\theta\) under all strategies.
Choose a polynomial-bit \(M>0\) larger than the absolute terminal rewards,
using \(\sum_v|d_v|\) for all subsets. Replace each reward \(r\) by a
Bernoulli win of probability \((r+M)/(2M)\). Almost-sure termination
makes this the same affine change of all state values. The fair-coin
samplers of Theorem~\ref{thm:ssg} implement the rational lotteries.

A state records the \(N\)-bit membership vector, a stage, and a position
in a sampler. Its two successors can be computed in polynomial time from
these fields, including the subset sum in a terminal reward. Unrolling
this computation gives a polynomial-size circuit for the graph. To keep
the explicit state count \(2^N\operatorname{poly}(L)\), each sampler
tracks its bit index and prefix comparisons, recomputing its thresholds
from the mask. Undo the reward normalization at the singleton states and
compute \(f_e\) from their values. The support LP recovers the market.
For an explicit bound, let \(d\) be the dimension of the perturbed market
LCP. Clear denominators in each of its possible tight rows, including
the right-hand side and the possible coordinate-zero equations. Let \(H\ge1\)
bound the absolute integer entries. On the support of the unique solution,
the corresponding principal submatrix is nonsingular, since the perturbed
matrix is a P-matrix. Cramer's rule bounds each coordinate denominator by
\(D=(dH)^d\). Choose the reward bound \(M\) above to be an integer. The
normalized singleton value is \((z_v+M)/(2M)\), so its denominator is at
most \(D'=2MD\). Both \(M\) and \(D'\) have polynomial binary length.
All the singleton states can thus be designated in advance for exact-value
queries of polynomial output length.

A decision oracle can instead compare the normalized singleton value
\(v\) to any rational \(q\in[0,1]\). Add a fair-coin start state choosing
between that singleton and a win of probability \(1-q\), and ask whether
the start value is at least \(1/2\). Its value is \((v+1-q)/2\), so
the comparison also handles equality. This addition preserves stopping
and has a polynomial-size successor circuit. Starting with \([0,1]\),
\(2\lceil\log_2 D'\rceil+2\) such adaptive comparisons give a closed
interval of width less than \(1/(2D'^2)\). Rational reconstruction returns
its unique rational of denominator at most \(D'\). Repeat for the \(N\)
singleton states and apply the decoder above.

Even on a star with a buyer at the center, the Min option at that
singleton can reach the union with any chosen incident good. From there
the same choice can add a further missing good while retaining the
center. Retaining the mask at each clip is always legal. Thus every one
of the \(2^{N-1}\) masks containing the center is reachable in the
explicit construction. This establishes the size distinction in
Theorem~\ref{thm:ssg}.

\chapter{A reverse reduction for a family of P-matrices}\label{sec:cayley}
\reviewmark{Cayley hypothesis, signed routing errors, rational reconstruction, direct convex program and an excluded P-matrix.}

Let \(A\) be a rational P-matrix and \(q\) a rational vector of length
\(n\). Given \(\lambda>0\), put
\[
 R=(I+\lambda A)^{-1}(I-\lambda A),\qquad
 d=-2\lambda(I+\lambda A)^{-1}q .
\]
The inverse exists by expansion of the determinant in principal minors.
Use the classical absolute-value representation of complementarity
\cite{gb}. For \(w=Az+q,\ z,w\ge0,\ z\perp w\), set \(x=z-\lambda w\).
Then \(|x|=z+\lambda w\), so \(x=R|x|+d\). Conversely
\(z=[x]_+\), \(w=[-x]_+/\lambda\) recover the solution.

Suppose \(\rho(|R|)<1\). Compute
\(v=(I-|R|^{\mathsf T})^{-1}\mathbf1\) and
\(t=(I-|R|)^{-1}\mathbf1\). Their entries are positive by the convergent
matrix series. Put \(V=\operatorname{diag}(v)\), \(H=VRV^{-1}\),
\(G=|H|\), \(d'=Vd\) and \(b_i=v_it_i/\max_j v_jt_j\). Then each
column of \(G\) sums to at most one and \(\gamma=b-Gb>0\). With
\(\sigma=\min(1,\min_i\gamma_i/[4(1+|d_i'|)])\), encode \(y_i=(Vx)_i\)
by two incoming allocations \(b_i/2+\sigma y_i,\ b_i/2-\sigma y_i\).
The unique solution \(y^*=H|y^*|+d'\) has
\(\sigma|y^*|\le(I-G)^{-1}\gamma/4=b/4\).

Two market constructions suffice to compute with these quantities.
First, a good can send the first \(r_1\) units of incoming trades to one
successor, the next \(r_2\) units to another, and so on. For \(h\)
successors give buyer \(k\) mass \(r_k\), a home offer \(1-P\), and an
offer \(1-k/[2(h+1)]-P_k/[4(h+1)]\) on successor \(k\).
Add a private singleton of mass \(A_0\) with offer \(1-P\) and give
the home good capacity \(A_0+\sum_k r_k\), with \(A_0\) exceeding all
possible incoming buyer mass. Above price one it could not fill, so
all such prices are in \([0,1]\). There the exit thresholds are strictly
ordered. At a price strictly between zero and one exactly the incoming
mass must leave in that order. At zero no mass enters, and at one all
tranches leave. Excess can be discarded. A constant amount is delivered
by a singleton offering \(2-P\), and a discard good has capacity above
all potential purchases. Second, at any chosen precision \(\delta=2^{-a}\), \(a\ge1\), we can divide
an input \(u\in[0,2]\) into two amounts within \(\delta/2\) of \(u/2\)
which sum to \(u\). Scale the alternating-switch market
of~\cite[Theorem~4.7]{book} by \(\delta\) in its masses and capacities,
choosing saturation \(2/\delta\) in the unscaled units. Its two exits
alternate every \(\delta\) units. The unscaled depth is \(a\). This uses \(O(1+\log(1/\delta))\)
buyers and goods. Choose its input capacity and local input buyer mass
above the possible incoming buyer mass and above two. At a positive input
price the local buyer must then purchase there, bounding the price by
one as in~\cite[Theorem~4.4]{book}. Thus the required port-price bounds
hold also when these constructions are connected in cycles. All padding
uses incoming buyer masses fixed before the private padding is added.

At coordinate \(j\) clip its actual two inputs to \([0,b_j]\), obtaining
\(p,r\), and write \(b=b_j\). Route their first halves and excesses using
the ordered quotas just described, to form the two outgoing quantities
\[
 A'=b/2+[p-b/2]_++[r-b/2]_+,\qquad
 B'=[\min(p,b/2)+\min(r,b/2)-b/2]_+ .
\]
An ideal pair \(b/2\pm s\) gives \(b/2\pm|s|\). On arbitrary inputs
this map is 1-Lipschitz from \(\ell_1\) to \(\ell_1\): a change to one
input splits monotonically between its two tranches, and clipping or
discarding can only reduce their total change. Both outputs lie in
\([0,3/2]\).

Distribute each of them to rows \(i\) in proportions \(G_{ij}\), swapping
the two destinations if \(H_{ij}<0\). For precision \(k\) use consecutive
blocks of \(\lfloor 2^kG_{ij}\rfloor\) leaves of a binary tree of depth
\(k\), with leftover leaves for discard. Split only nodes containing a
boundary, using the alternating divider above. There are \(O(nk)\) such
nodes, with at most two maximal dyadic subtrees per level for any block.
The amount at depth \(a\) differs from its ideal share by at most
\(\delta(1-2^{-a})\) by induction. The absolute error delivered to a
particular destination from one outgoing quantity is at most
\(\eta=(3/2)2^{-k}+(2k+1)\delta\), including rounding the block length.
Add constant sources \(\gamma_i/2\pm\sigma d'_i\) at the two inputs of
row \(i\) and close the circuit. With ideal proportions the encoded
solution gives exactly its starting quantities at each row.

At any actual equilibrium let \(D_i\) be the sum of the two absolute
input errors at row \(i\) relative to that ideal solution. The
nonexpansive calculation and routing error give
\(D\le GD+2n\eta\mathbf1\). Thus
\(\|D\|_\infty\le2n\eta\|(I-G)^{-1}\|_\infty\), even for actual pairs
whose total differs from \(b_i\). A uniform denominator bound \(U\) for
\(x^*\) follows by Cramer's rule on any complementary basis of \(A\),
followed by the expression \(x^*=z-\lambda(Az+q)\). Choose \(k\) and
\(\delta\le2^{-k}/(2k+1)\) so the displayed error divided by
\(\sigma\min_i v_i\) is less than \(1/(2U^2)\). All matrices, bounds and
inverse scales have polynomial bit length. Hence \(k,\log(1/\delta)\)
and the explicit market size are polynomial. Its observed first input
in each row gives \(x_i^*\) by rational reconstruction, and then \(z,w\)
by their positive parts. The construction uses the supplied step
\(\lambda\) and the spectral hypothesis on its Cayley transform.

\begin{lemma}[Scope of the spectral promise]\label{lem:cayley-qp}
For a supplied rational \(\lambda>0\) with \(\rho(|R|)<1\), the
original rational P-LCP can also be solved in polynomial bit time by
one convex quadratic program. There are P-matrices for which the
spectral promise fails for every positive \(\lambda\).
\end{lemma}
\begin{proof}
Put \(B=|R|\) and take \(t,v\) as above, so \(Bt=t-\mathbf1\) and
\(B^{\mathsf T}v=v-\mathbf1\). Write \(D=\operatorname{diag}(v_i/t_i)\).
For \(x\ne0\), weighted Cauchy--Schwarz in each row gives
\[
\begin{aligned}
 x^{\mathsf T}R^{\mathsf T}DRx
 &\le \sum_i (v_i/t_i)(Bt)_i\sum_j B_{ij}x_j^2/t_j\\
 &\le \sum_j (v_j-1)x_j^2/t_j < x^{\mathsf T}Dx .
\end{aligned}
\]
For \(S=I+\lambda A\), commutation gives \(RS=I-\lambda A\).
Expanding \(S^{\mathsf T}(D-R^{\mathsf T}DR)S\) therefore yields
\(2\lambda(A^{\mathsf T}D+DA)\succ0\). Minimize
\[
 z^{\mathsf T}D(Az+q)\qquad\text{subject to }z\ge0,\quad Az+q\ge0 .
\]
The Hessian is positive definite. Feasible points have nonnegative
objective, and the P-LCP solution attains zero. Conversely, at zero
each term \(D_{ii}z_i(Az+q)_i\) vanishes, giving complementarity.
The two rational linear solves defining \(D\) have polynomial bit
length. Rational convex quadratic programming gives the exact optimum
in polynomial bit time~\cite{ktk}; the unique optimizer here is also
rational by the complementary-basis equations.

For the second assertion let \(C\) cyclically permute three coordinates
and put \(A=I+2C\). Every proper nonempty principal minor is one and
\(\det A=9\). Its eigenvalues are \(3\) and \(\pm i\sqrt3\). Thus \(R\)
has an eigenvalue of modulus one for every \(\lambda>0\), and
\(\rho(|R|)\ge\rho(R)\ge1\).
\end{proof}

\cleardoublepage
\phantomsection
\addcontentsline{toc}{chapter}{Bibliography}
\begin{thebibliography}{99}
\bibitem{alon}
N.~Alon. Explicit expanders of every degree and size.
\emph{Combinatorica} (2021).
\href{https://arxiv.org/abs/2003.11673}{arXiv:2003.11673}.

\bibitem{abr}
E.~R.~Arunachaleswaran, S.~Barman and N.~Rathi. Fully polynomial-time
approximation schemes for fair rent division.
\emph{Mathematics of Operations Research} 47 (2022).
\href{https://arxiv.org/abs/1807.04163}{Author preprint}.

\bibitem{apt}
B.~Aspvall, M.~F.~Plass and R.~E.~Tarjan. A linear-time algorithm for
testing the truth of certain quantified Boolean formulas.
\emph{Information Processing Letters} 8 (1979), 121--123.
\href{https://doi.org/10.1016/0020-0190(79)90002-4}{Publisher's version}.

\bibitem{karp}
R.~M.~Karp. Reducibility among combinatorial problems. In
\emph{Complexity of Computer Computations}, 1972, 85--103.
\href{https://doi.org/10.1007/978-1-4684-2001-2_9}{Original chapter}.

\bibitem{rrvv}
A.~E.~Roth, U.~G.~Rothblum and J.~H.~Vande Vate. Stable matchings,
optimal assignments, and linear programming.
\emph{Mathematics of Operations Research} 18 (1993), 803--828.
\href{https://doi.org/10.1287/moor.18.4.803}{Publisher's version}.
\bibitem{conitzer}
V.~Conitzer. Open problem: Computing Bayes--Nash equilibrium in simple
security games. Prepared for the Fall 2015 Simons Institute program
\emph{Economics and Computation}. Revised arXiv version, 2026.
\href{https://arxiv.org/abs/2610.08796v2}{arXiv:2610.08796v2}.

\bibitem{condon}
A.~Condon. The complexity of stochastic games.
\emph{Information and Computation} 96 (1992), 203--224.
\href{https://doi.org/10.1016/0890-5401(92)90048-K}{doi:10.1016/0890-5401(92)90048-K}.

\bibitem{ds}
Y.~Disser and M.~Skutella. The simplex algorithm is NP-mighty.
\emph{ACM Transactions on Algorithms}, 2018.
\href{https://arxiv.org/abs/1311.5935}{arXiv:1311.5935}.

\bibitem{arrival}
J.~Dohrau, B.~G\"artner, M.~Kohler, J.~Matou\v{s}ek, and E.~Welzl.
ARRIVAL: A zero-player graph game in \(\mathrm{NP}\cap\mathrm{coNP}\).
In \emph{A Journey Through Discrete Mathematics}, 2017.
\href{https://arxiv.org/abs/1605.03546}{arXiv:1605.03546}.

\bibitem{has}
S.~Haslebacher. ARRIVAL: Recursive framework \& \(\ell_1\)-contraction.
In \emph{ICALP 2025}, LIPIcs 334, Article 95.
\href{https://doi.org/10.4230/LIPIcs.ICALP.2025.95}{Publisher's version}.

\bibitem{ghhkms}
B.~G\"artner, T.~D.~Hansen, P.~Hub\'a\v{c}ek, K.~Kr\'al, H.~Mosaad
and V.~Sl\'ivov\'a. ARRIVAL: Next Stop in CLS. In \emph{ICALP 2018}.
Definition numbers refer to \href{https://arxiv.org/abs/1802.07702}{arXiv:1802.07702}.

\bibitem{kim}
J.~Kim. Public reply to V.~Conitzer about ARRIVAL and one-resource
security games. November 22, 2026.
\href{https://x.com/JeonseongKim/status/2124443473428537685}{Original post}.

\bibitem{fgms}
J.~Fearnley, S.~Gordon, R.~Mehta and R.~Savani.
Unique end of potential line. \emph{Journal of Computer and System Sciences}
114 (2020), 1--35.
\href{https://arxiv.org/abs/1811.03841}{arXiv:1811.03841}.

\bibitem{gj}
A.~Galichon and A.~Jacquet. The matching problem with linear transfers is
equivalent to a hide-and-seek game. \emph{Games and Economic Behavior}
152 (2025), 333--344.
\href{https://arxiv.org/abs/2402.12200}{arXiv:2402.12200}.

\bibitem{gb}
A.~Griewank, J.~U.~Bernt, M.~Radons and T.~Streubel.
Solving piecewise linear systems in abs-normal form.
\emph{Linear Algebra and its Applications} 471 (2015), 500--530.
\href{https://doi.org/10.1016/j.laa.2014.12.017}{Publisher's version}.

\bibitem{glm}
E.~Ghorbani, R.~Lopes, and M.~Mnich. ARRIVAL in cube-root subexponential
time. Preprint, 2026.
\href{https://arxiv.org/abs/2610.28241v1}{arXiv:2610.28241v1}.

\bibitem{generic}
ChatGPT-6 Astra (attribution in \cite{conitzer}).
\emph{Interchangeability in Bayesian Security Games: Exact Balance and
Generic Uniqueness}. Research note, September 6, 2026.
\href{https://www.cs.cmu.edu/~conitzer/gpt6interchangeability.pdf}{Manuscript hosted by V.~Conitzer}.

\bibitem{jg}
C.~Jones and M.~S.~Gowda. On the connectedness of solution sets in
linear complementarity problems. \emph{Linear Algebra and its Applications}
272 (1998), 33--44.
\href{https://doi.org/10.1016/S0024-3795(97)00282-6}{doi:10.1016/S0024-3795(97)00282-6}.

\bibitem{ktk}
M.~K.~Kozlov, S.~P.~Tarasov and L.~G.~Khachiyan.
The polynomial solvability of convex quadratic programming.
\emph{USSR Computational Mathematics and Mathematical Physics}
20(5) (1980), 223--228.
\href{https://doi.org/10.1016/0041-5553(80)90098-1}{Publisher's version}.

\bibitem{book}
Research project coordinated by V.~Lin. \emph{Bayes--Nash Equilibrium in
Bayesian Simple Security Games: Markets, ARRIVAL, Contractions and Path
Following}. Public manuscript, November 2026, revision
\href{https://github.com/vvncent/bayes-nash-equilibrium-simple-security-games/blob/095bb77fb6f15406d01f1ed2c8cb8038328f100b/main.pdf}{095bb77}.

\bibitem{mp}
F.~Meunier and T.~Pradeau. Computing solutions of the multiclass network
equilibrium problem with affine cost functions.
\emph{Annals of Operations Research} 274 (2019), 447--469.
\href{https://doi.org/10.1007/s10479-018-2817-z}{doi:10.1007/s10479-018-2817-z}.
\bibitem{milchtaich}
I.~Milchtaich. Topological conditions for uniqueness of equilibrium in
networks. \emph{Mathematics of Operations Research} 30 (2005), 225--244.
\href{https://doi.org/10.1287/moor.1040.0122}{Publisher's version}.
\end{thebibliography}
\end{document}
